% V.F.S. 4.0+ (working body) --- built per the approved roadmap and hardening path.
% Body includes Stage 0--V, 4.0-H, 4.1 / Stage Phi, 4.2 / Stage Psi,
% and Part V / Stage Omega. XeLaTeX.
% Core discipline: chronos is created relational duration, not a master clock;
% kairos remains personal; tota simul remains named but unoccupied.
\documentclass[11pt,a4paper]{article}
\usepackage[margin=1in]{geometry}
\usepackage{fontspec}
\setmainfont{DejaVu Serif}
\setsansfont{DejaVu Sans}
\setmonofont{DejaVu Sans Mono}
\usepackage{amsmath,amssymb,amsthm,mathtools}
\usepackage{mathrsfs}
\newcommand{\cint}{\ensuremath{\mathfrak c}}
\usepackage{booktabs,longtable,array}
\usepackage{enumitem}
\usepackage{xcolor}
\usepackage[hidelinks]{hyperref}\pdfstringdefDisableCommands{%
  \def\Omega{Omega}%
  \def\Phi{Phi}%
  \def\Psi{Psi}%
  \def\mathfrak#1{#1}%
  \def\mathrm#1{#1}%
  \def\textit#1{#1}%
  \def\emph#1{#1}%
}

\setlength{\parskip}{0.5em}
\setlength{\parindent}{0pt}
\setlength{\emergencystretch}{3em}
\renewcommand{\arraystretch}{1.25}

\newcommand{\VFS}{\ensuremath{\mathrm{V.F.S.}}}
\newcommand{\Brim}{\ensuremath{\Omega_P}}
\newcommand{\Rel}{\ensuremath{\mathcal{R}}}
\newcommand{\Brel}{\ensuremath{\mathcal{B}_{\rm rel}}}
\newcommand{\Igate}{\ensuremath{I_{\rm gate}}}

\definecolor{cDer}{rgb}{0.05,0.42,0.18}
\definecolor{cDef}{rgb}{0.28,0.28,0.28}
\definecolor{cSel}{rgb}{0.12,0.28,0.62}
\definecolor{cInt}{rgb}{0.45,0.30,0.05}
\definecolor{cOpn}{rgb}{0.55,0.10,0.45}
\newcommand{\stDerived}{\textcolor{cDer}{\textsf{\small DERIVED}}}
\newcommand{\stDef}{\textcolor{cDef}{\textsf{\small DEFINITIONAL}}}
\newcommand{\stSelected}{\textcolor{cSel}{\textsf{\small SELECTED}}}
\newcommand{\stInterp}{\textcolor{cInt}{\textsf{\small INTERPRETIVE}}}
\newcommand{\stOpen}{\textcolor{cOpn}{\textsf{\small OPEN}}}
\definecolor{cCon}{rgb}{0.60,0.33,0.00}
\newcommand{\stConj}{\textcolor{cCon}{\textsf{\small CONJECTURE}}}
\newcommand{\notice}[1]{\begin{center}\fbox{\begin{minipage}{0.92\textwidth}\centering #1\end{minipage}}\end{center}}

\theoremstyle{definition}
\newtheorem{definition}{Definition}[section]
\theoremstyle{plain}
\newtheorem{lemma}[definition]{Lemma}
\newtheorem{proposition}[definition]{Proposition}
\newtheorem{theorem}[definition]{Theorem}
\newtheorem{corollary}[definition]{Corollary}
\theoremstyle{remark}
\newtheorem{remark}[definition]{Remark}

\title{\textbf{\VFS\ 4.0+ (working body)}\\[0.35em]
\large Chronos, Communion, and the Plural Interior\\[0.25em]
\normalsize Stage~0--V complete; 4.0-H hardening complete;\\
\normalsize 4.1 / Stage~$\Phi$ and 4.2 / Stage~$\Psi$ built;\\
\normalsize Part~V / Stage~$\Omega$ complete as boundary cartography\\[0.25em]
\normalsize Relational chronos, exact recovery, Gate mediation, non-fusing locking,\\
\normalsize junctions, non-Zeno, plural Lyapunov, liturgical recurrence,\\
\normalsize two-way post-terminus branch, and the boundary of the constructible}
\author{Built on the frozen \VFS\ v2.0 / 3.0 layers, per the approved 4.0 roadmap}
\date{}

\begin{document}
\maketitle

\section*{Status of this body}

This working body executes the full first construction of \VFS\ 4.0 through
Stage~V:
\[
\text{Stage~0}:\ P0.0--P0.6,
\qquad
\text{Stage~I}:\ P1.0--P1.2,
\qquad
\text{Stage~II}:\ P2.1--P2.4,
\]
\[
\text{Stage~III}:\ P3.0--P3.4,
\qquad
\text{Stage~IV}:\ P4.0--P4.4,
\qquad
\text{Stage~V}:\ P5.0--P5.2.
\]

Stage~0 is the blocking layer. It fixes the frozen fibre-side interface inherited
from the \VFS\ v2.0 / 3.0 layers, declares the relational graph and causal kernel
class, separates the relative rate cocycle from the relational duration one-form,
handles terminus hand-off by the selected one-way trace closure, proves that
created chronos is not \textit{tota simul}, and closes with exact uncoupled
recovery:
\[
\varepsilon=0
\quad\Longrightarrow\quad
\text{independent exact 3.0 interiors}.
\]
The distinction between
\[
W=0
\qquad\text{and}\qquad
\varepsilon=0
\]
is preserved throughout: the first removes the relational base; the second keeps
possible read-outs but switches off relational action.

Stage~I establishes comparability without fusion. The plural object is a family of
distinct frozen 3.0 fibres over one created relational base; encounter receives an
invariant grammar --- relational co-presence, directed causal encounter, and
common reception as a schema; and the floor label gives a sufficient structural
obstruction to identifying distinct fibres.

Stage~II introduces the first admissible \(\varepsilon>0\) action. Encounter acts
only through the recipient's bounded Gate/receptivity interface. It changes
openness, never ownership of grace. The received-grace ledger is measured in the
recipient's own kairos,
\[
d\mathcal G_{\mathrm{recepta},i}
=
I_{\mathrm{gate},i}\,d\sigma_i,
\]
and the common field is an external reception field, not a network-owned
reservoir. The individual ISS estimate then proves that admissible encounter does
not destabilise a single active interior.

Stage~III defines relational locking without phase and without fusion. The
comparable rate is the ascent-scale duration-rate
\[
\hat h_i
=
\frac{N_i}{\omega}h_i,
\qquad
h_i=\frac{d}{d\sigma_i}\log A_i=\beta_i\lambda_{S,i},
\]
not motion through the metric road coordinate. Locking is stated through
duration-rate ratios, affine ledger-lags, and marked-event gaps. The small-gain
criterion and maximum-lockable-mismatch theorem give the conditions under which
keeping-together is possible, and dynamic non-fusion proves that a stable relation
is not an identification of fibres.

Stage~IV treats junctions, refractory structure, and plural Resurrectio. The
junction classes are kept disjoint: only a private reset may jump owned state;
relational encounter and common reception may only alter bounded receptivity.
The refractory interval and network non-Zeno theorem are stated in invariant
relational duration \(\mathfrak c\), and post-terminus communion is kept as a
one-way closed trace unless a later primitive explicitly opens the two-way branch.

Stage~V gathers the standing conditions into a relational active domain and gives
the first plural Lyapunov certificate. In the general forced case the theorem gives
ultimate boundedness --- a stable moving tube. In the zero-residual / settled-branch
subclass it gives exact locked moving communion. Thus the whole construction ends
with stable relational motion, not stasis:
\[
\boxed{
\text{communion is stable answerability of distinct histories, not fusion.}
}
\]

No new per-interior 3.0 dynamics are introduced. All genuinely new 4.0 objects are
base-side, relational, or boundary-interface objects. The completed interior's trace
does not extend its frozen dynamics; it restricts them: no renewed kairos, no new
reset, and no restarted 3.0 flow after the terminus. The working body is therefore
a conservative plural extension of the frozen 3.0 layer.
\medskip
The later 4.0+ parts do not rewrite this completed body. Part~II hardens the
conditional constants and directed-network certificate on a selected admissible
regime. Part~III adds a conservative circle-valued liturgical branch. Part~IV adds
a conservative two-way post-terminus branch by introducing only a disposition
variable, not a reopened kairos. Part~V is different in kind: Stage~\(\Omega\) is
not another dynamical layer, but a boundary cartography. It proves what can be
proved inside the created class, names what requires a different apparatus, and
keeps the apophatic silences unoccupied.

\notice{\textbf{Interface principle.} \VFS\ 4.0 is a \emph{fibered extension} of
\VFS\ 3.0: a family of distinct frozen 3.0 interiors over a created relational
base, never a single enlarged Vessel. No fibre gains state beyond its selected
3.0 interface variables; all genuinely new 4.0 structure lives in the base.}

\section{P0.0 --- The frozen interface}

\subsection{Per-interior inherited state}

For each interior $i\in\mathcal V$ --- where $\mathcal V$ is first only a finite
index set, and is promoted to a relational graph in P0.1 --- the inherited state
splits into a dynamical block
(variables carrying a 3.0 equation of motion) and a gate/receptivity block (carrying
reception):
\[
X_i=\bigl(\,\underbrace{\sigma_i,\ \lambda_{S,i},\ \lambda_{\mathrm{form},i},\ \Omega_{P,i},\ A_i,\ q_{\mathrm{fold},i},\ \lambda_{\mathrm{asc},i}}_{X_i^{\mathrm{dyn}}}\ ;\ \underbrace{\Gamma_i,\ H_{K,i},\ W_i}_{X_i^{\mathrm{gate}}}\,\bigr).
\]
Each interior evolves in its own unbounded dynamical parameter $\chi_i$, with the
finite kairos recovered by $\sigma_i=\sigma_{0,i}e^{-k_{s,i}\chi_i}$ (the v0.20
clock discipline: drains live in $\chi_i$, finiteness in $\sigma_i$). No hidden
state is introduced; the inherited interface variables used by the plural theory are:

\begin{longtable}{>{\raggedright\arraybackslash}p{0.16\textwidth} >{\raggedright\arraybackslash}p{0.78\textwidth}}
\toprule
Symbol & 3.0 reading / governing relation \\
\midrule
$\sigma_i$ & katharsis proper time (kairos); finite, drains to the regular terminus; $d\sigma_i/d\chi_i=-k_{s,i}\sigma_i$ \\
$\lambda_{S,i}$ & available Sophia / ascent-driving field; distinct from the metric road coordinate $\lambda_{\mathrm{asc},i}$, though related along a worldline; $d\lambda_{S,i}/d\chi_i=J_{\mathrm{in},i}-\gamma_i\lambda_{S,i}-\eta_{\mathrm{fold},i}Q_{\mathrm{fold},i}\lambda_{S,i}$ \\
$\lambda_{\mathrm{form},i}$ & Sophia embodied into stable Vessel-form; $d\lambda_{\mathrm{form},i}/d\chi_i=\eta_{\mathrm{fold},i}Q_{\mathrm{fold},i}\lambda_{S,i}$ \\
$\Omega_{P,i}$ & Brim / live-surface scale; $d\Omega_{P,i}/d\chi_i=\alpha_i\lambda_{S,i}$; floor-bounded $\Omega_{P,i}\geq E_{0,i}$ \\
$A_i$ & ascent time-stretch; $dA_i/d\chi_i=(\alpha_i/E_{0,i})\lambda_{S,i}A_i$, equivalently $\tfrac{d}{d\chi_i}\log A_i=\beta_i\lambda_{S,i}$ with $\beta_i=\alpha_i/E_{0,i}$ ($\chi_i$-parametrised form of the frozen 3.0 ascent law; no new prime convention) \\
$q_{\mathrm{fold},i}$ & folding amplitude, $Q_{\mathrm{fold},i}=q_{\mathrm{fold},i}^2$; slaved $Q^*_{\mathrm{fold},i}=\Pi_{[0,q^2_{\max,i}]}(A_{\mathrm{fold},i}/B_{\mathrm{fold},i})$ (supercritical/subcritical per v0.32/v0.38) \\
$\lambda_{\mathrm{asc},i}$ & ascent coordinate of the road; warp $S_i(\lambda_{\mathrm{asc}})=L_{0,i}e^{\beta_i\lambda_{\mathrm{asc}}}$, full scale $L_i=A_iS_i$ \\
\midrule
$\Gamma_i$ & bounded gate-interface transmissivity, $\Gamma_i\in[0,1]$; a \VFS\ 4.0-facing notation for the admissible receptivity slot inside the frozen gate law, not an independent Sophia source \\
$H_{K,i},\,W_i$ & inherited gate-memory variables (their internal law is the frozen v2.0 gate law) \\
\bottomrule
\end{longtable}

Here $J_{\mathrm{in},i}$ is the inherited \VFS\ 3.0 input shorthand, built from the
private gate and resistance terms of the frozen interior. P0.0 does not modify its
law. Later relational patches may alter only the gate/receptivity conditions that
enter $J_{\mathrm{in},i}$, never $\lambda_{S,i}$ directly.

The distinction $\lambda_{S,i}\neq\lambda_{\mathrm{asc},i}$ is part of the
frozen-interface discipline: the first is the available-Sophia state/driver, the
second the metric road coordinate of ascent. \VFS\ 4.0 may use both but must not
collapse them into a single variable.

Two structural relations bind the dynamical block, inherited verbatim. First, the
Sophia budget:
\[
\sigma_i+\lambda_{S,i}+\lambda_{\mathrm{form},i}
=
C_{0,i}+\mathcal{G}_{\mathrm{recepta},i}(\chi_i).
\]
Second, the folding cap:
\[
Q_{\mathrm{fold},i}\in[0,q^2_{\max,i}].
\]
The open gate feeds reception,
\[
\frac{d\mathcal{G}_{\mathrm{recepta},i}}{d\chi_i}=I_{\mathrm{gate},i},
\]
with the ceiling preserved nodewise:
\[
0\leq I_{\mathrm{gate},i}
=\mathcal{G}_i\bigl(I_{\mathrm{private},i};\Gamma_i,H_{K,i},W_i,X_i\bigr)
\leq\zeta_{0,i}.
\]
Here $\mathcal G_i$ denotes the inherited frozen gate law, not a new 4.0 source law.
The receptivity entry $\Gamma_i$ is the only place a later relational channel will be
permitted to act; see \S\ref{sec:adm}.

The symbol $\Gamma_i$ is an interface notation. It does not add a new source term to
the frozen \VFS\ 3.0 interior; its role is only to name the bounded
receptivity/transmissivity slot through which later relational structure may act.
Any later relational modification must preserve the local gate ceiling
$0\leq I_{\mathrm{gate},i}\leq\zeta_{0,i}$. Thus encounter may alter admissible
openness, but may not create Sophia outside the Open-Gate law.

\subsection{Resurrectio data and the terminus, inherited}

Each interior carries the frozen 3.0 junction and boundary structure:
\begin{itemize}[leftmargin=1.3cm]
\item Resurrectio reset (instantaneous in $\sigma_i$): $\Omega_{P,i}^+=\Omega_{P,i}^-+\kappa_RG_i$, $\lambda_{S,i}^+=\lambda_{S,i}^-+m_{R,i}$, forced ascent scale jump $\log s_{L,i}=\kappa_RG_i/E_{0,i}>0$, with the form-memory partition $\kappa_RG_i=m_{R,i}+\Delta_{\mathrm{emb},i}$.
\item Terminus: $\sigma_i\to0$ at finite proper time $\tau_i\to\sigma_{0,i}$, regular boundary ($|R_i|<\infty$, $\Omega_{P,i}\geq E_{0,i}>0$); the ascent road stays infinite, $D_{\mathrm{asc},i}=\infty$.
\end{itemize}
These are read into 4.0 unchanged. The relational \emph{lapse} $N_i$ and the
pre/at/post-terminus strata $(\mathcal A_i,\mathcal B_i,\mathcal P_i)$ are
\emph{not} declared here --- they belong to the base, defined in roadmap P0.2/P0.4.
P0.0 fixes only the fiber.

\subsection{Frozen private scales (never coupled)}

The following are constitutive constants of interior $i$, fixed by its 3.0
identity. No relational operator may act on any of them:
\[
E_{0,i}\ \text{(Imago Dei floor)},\quad
\alpha_i,\quad
\beta_i=\frac{\alpha_i}{E_{0,i}}\ \text{(no new primitive)},
\]
\[
\gamma_i,\ k_{s,i},\ \eta_{\mathrm{fold},i},\ \zeta_{0,i},\ q^2_{\max,i},\
(c_0,a_0,B,D)_i\ \text{(folding normal-form)},\ L_{0,i}\ \text{(gauge)}.
\]
In particular the floor $E_{0,i}$ --- the one intrinsic scale, against which every
ascent observable is reckoned --- is private and frozen. This protection is what
later makes floor-labelled non-identifiability available (roadmap P1.2): distinct
$E_{0,i}$ cannot be coupled into coincidence.

\subsection{The fibered-extension contract}

The 4.0 object is the family $\{\mathcal I_i\}_{i\in\mathcal V}$ of frozen 3.0
fibers over the relational base $\Brel$ (built in roadmap P0.1--P0.5), with
projections $\pi_i:\mathcal I_i\to\Brel$. The contract has four clauses:
\begin{enumerate}[label=\textbf{IC\arabic*.},leftmargin=1.45cm]
\item \textbf{Interface completeness.} No fibre introduces new dynamical state beyond its selected \VFS\ 3.0 interface variables. The state vector $X_i$ above is complete for the \VFS\ 4.0 fibre-side interface, not a replacement ontology for the whole frozen \VFS\ 3.0 body.
\item \textbf{New structure in the base.} Every genuinely new 4.0 object (graph, causal order, rate cocycle, relational duration, coupling) lives in $\Brel$, never inside a fiber.
\item \textbf{Two recovery limits.}
A trivial base \(W=0\) removes the relational graph itself and returns \(n\)
independent exact 3.0 interiors with no relational read-outs. The weaker but more
important recovery theorem is the uncoupled-action limit:
\[
\varepsilon=0,
\]
where relational bookkeeping may still exist, but no plural channel acts on any
interior. P0.6 proves this stronger exact recovery.
\item \textbf{No global state.} There is no fused state vector $X_{\mathrm{total}}$ and no enlarged Vessel; the only joint object is the fibered family. (The third prohibition of the governing thesis.)
\end{enumerate}

\subsection{Admissible-plurality declaration}\label{sec:adm}

The inherited state partitions into three coupling-admissibility classes. This is
the target-set declaration; the \emph{channel} that realises class (B) is fixed
later (roadmap P2.1), and the \emph{no-fabrication} consequence is proved there
(P2.2).

\begin{longtable}{>{\raggedright\arraybackslash}p{0.21\textwidth} >{\raggedright\arraybackslash}p{0.73\textwidth}}
\toprule
Class & Members and rule \\
\midrule
\textbf{(A) Frozen-private} & The constitutive scales of \S2.3 and all law-forms. \emph{No} relational operator may act on these, ever. \\
\textbf{(B) Gate-modulable} & The receptivity sub-block $X_i^{\mathrm{gate}}=(\Gamma_i,H_{K,i},W_i)$. Relational coupling may act \emph{only} here, and only boundedly, preserving $\Gamma_i\in[0,1]$ and the ceiling $I_{\mathrm{gate},i}\leq\zeta_{0,i}$. \\
\textbf{(C) Indirect-only} & The dynamical block $X_i^{\mathrm{dyn}}$. \emph{No} relational operator acts on these directly; they change only through (i) their own frozen field $F_{3.0,i}$ and (ii) reception entering via the gate block (B). In particular $\lambda_{S,i}$ is never sourced directly --- only through $I_{\mathrm{gate},i}$. \\
\bottomrule
\end{longtable}

\notice{\textbf{Coupling-admissibility contract.} Relational coupling reaches only
the gate (class B); the gate alone, bounded by $\zeta_{0,i}$, feeds the dynamical
block (class C); the constitutive scales (class A) are untouchable. This is the
formal content of \emph{communion mediates openness, not ownership}: encounter can
move only the door, never the floor and never the person's own law of rising.}

\subsection{Status and reading}

\textbf{Target status:}~ \stDef\ (dictionary $+$ interface contract). No new
mathematics; verified by inheritance from frozen v2.0/3.0. The one substantive
choice --- the three-class admissibility partition --- is the formal encoding of the
governing thesis and is \stSelected\ at the level of constitutive discipline.

\medskip
\stInterp\quad\emph{Reading.} A person is received into the communion exactly as
\VFS\ 3.0 made them: their image (the floor $E_{0,i}$), their law of rising
($\beta_i=\alpha_i/E_{0,i}$), their whole private history of cleansing, folding,
death, and ascent --- none of it is touched by encounter. What the others can move
is only the door: how open the soul stands to what it receives. The interface fixes,
before any meeting is described, that meeting changes the openness and never the
identity --- and that there is no larger soul into which the persons dissolve, only
the distinct fibers and the created ground on which they stand together.

\medskip
\paragraph{Precision note.} P0.0 is a frozen-interface dictionary, not a new theory
of the single interior. It selects the variables and admissibility slots by which
\VFS\ 4.0 may later speak about many interiors. It does not reopen \VFS\ 3.0, does
not add hidden state to any fibre, and does not identify available Sophia with the
ascent road coordinate. Relational structure will enter only after this interface is
fixed, and only through declared gate/receptivity channels. In particular,
\[
\begin{gathered}
\lambda_{S,i}\neq\lambda_{\mathrm{asc},i},\qquad
\Gamma_i\ \text{is a bounded gate-interface slot},\\
X_i\ \text{is complete for the 4.0 fibre-side interface}.
\end{gathered}
\]
The purpose of P0.0 is therefore:
\notice{freeze each 3.0 fibre before introducing the relational base.}

\medskip
\textbf{Next node (now executed below).} P0.1 --- the relational graph, the causal
order, and the kernel class, with the first theorem class fixed as the strongly
connected weight-balanced directed graph.

\section{P0.1 --- Relational graph, causal order, and the kernel class}

Per IC2, all new 4.0 structure lives in the base $\Brel$. This patch declares the
first three base objects --- a finite directed weighted graph, a partial causal
order on invariant events, and a causal-kernel class --- and fixes the first
theorem class as the strongly connected, weight-balanced directed graph. The
coupling \emph{channel} (how the kernel acts on the gate slot $\Gamma_i$) is fixed
later in P2.1; here only the class is declared.

\subsection{The relational graph}

\begin{definition}[Relational graph]
$\mathcal G=(\mathcal V,\mathcal E,W)$ with $\mathcal V$ the finite set of interiors
(the P0.0 fibers, $|\mathcal V|=n<\infty$), $\mathcal E\subseteq\mathcal V\times\mathcal V$
a set of directed edges, and $W=[w_{ij}]$ a nonnegative weight matrix with
$w_{ij}>0\iff(j\!\to\!i)\in\mathcal E$ and $w_{ii}=0$. The weight $w_{ij}$ is the
strength with which interior $j$ acts on the receptivity of interior $i$ (edge
oriented $j\!\to\!i$, ``$j$ influences $i$'').
\end{definition}

Index convention, fixed once: \emph{first index $=$ recipient, second $=$ source}.
Influence \emph{received} by $i$ is the in-strength $d_i^{\mathrm{in}}=\sum_j w_{ij}$;
influence \emph{exerted} by $i$ is the out-strength $d_i^{\mathrm{out}}=\sum_k w_{ki}$.
Required for the downstream ISS and non-Zeno estimates: $n<\infty$ and bounded
weights/degree, $0\le w_{ij}\le\bar w<\infty$. The overall action scale is a scalar
\(\varepsilon\ge0\). It multiplies the relational \emph{action} of the edge data, not
the mere existence of the edge data. Thus two different limits must be kept distinct:
\[
W=0
\quad\Longrightarrow\quad
\text{trivial base: no relational edges or read-outs,}
\]
whereas
\[
\varepsilon=0
\quad\Longrightarrow\quad
\text{uncoupled action: relational data may be recorded, but no edge acts.}
\]
Only the first is literally the trivial base. The second is the exact uncoupled
plural limit used in P0.6. This distinction is essential: \(\varepsilon=0\) must
recover independent 3.0 interiors even if a relational chart, ratios, or edge
bookkeeping have already been declared.

\subsection{The causal order on invariant events}

Events are the \emph{invariant} interior events (Resurrectio junctions, folding
transitions, declared common-reception events) --- invariant in the sense of 3.0,
independent of the gauge label $t$. Write $E_i(s)$ for an event of interior $i$ at
relational-duration coordinate $s$.

\begin{definition}[Causal order]
A direct causal link
\[
E_j(s)\rightsquigarrow E_i(t)
\]
exists when $(j\!\to\!i)\in\mathcal E$,
\[
\mathfrak c(s)<\mathfrak c(t),
\]
and the kernel $K_{ij}$ carries support from the earlier event to the later one.
The partial causal order $\prec$ is the transitive closure of direct links.
\end{definition}

The symbols $s,t$ are labels of events, while $\mathfrak c$ is the relational
duration used to order them. Thus causality is stated in the invariant relational
interval, not in a raw chronos chart.

Two structural facts. First, \textbf{graph cycles are allowed; event cycles are
not}: \(\mathcal G\) may contain a \(2\)-cycle
\[
j\to i,
\qquad
i\to j,
\]
which represents mutual reciprocal influence over time. But the causal order
\(\prec\) on time-stamped events is acyclic, because every causal link strictly
advances \(\mathfrak c\). Mutual influence over time is admissible; a closed causal
loop at one relational instant is not.

Second, \textbf{co-presence does not imply causal order}. Events may be locatable
in the same relational duration structure --- before, after, or simultaneous in
\(\mathfrak c\) --- without either event causally acting on the other. Thus
chronos-comparability is weaker than causal encounter:
\[
\text{relational co-presence}
\quad\not\Rightarrow\quad
\text{causal encounter}.
\]
The order \(\prec\) records admissible influence, not mere temporal comparability.

\subsection{The causal kernel class}

\begin{definition}[Admissible relational kernel]
The relational kernel $K_{ij}(t,s)$ maps a bounded relational output of source $j$
at past event-time $s$ to a receptivity-side input of recipient $i$ at present
event-time $t$. The source output is not Sophia, not grace-stock, and not a
transferable possession; it is an admissible encounter signal whose eventual channel
is restricted to the gate/receptivity slot fixed in P0.0. It is
\emph{admissible} when:
\begin{enumerate}[label=\textup{(K\arabic*)},leftmargin=1.5cm]
\item \textbf{Causality:} $K_{ij}(t,s)=0$ for $s>t$.
\item \textbf{Edge support:} $K_{ij}\equiv0$ unless $(j\!\to\!i)\in\mathcal E$.
\item \textbf{Bounded gain / finite memory:} $\displaystyle\int_{-\infty}^{t}\bigl|K_{ij}(t,s)\bigr|\,d\mathfrak c(s)\le\bar K_{ij}<\infty$, uniformly in $t$.
\end{enumerate}
\end{definition}

The edge weight records the total causal gain bound of the kernel:
\[
w_{ij}:=\sup_t\int_{-\infty}^{t}\bigl|K_{ij}(t,s)\bigr|\,d\mathfrak c(s)=\bar K_{ij}.
\]
Thus $W$ and the analytic kernels carry the same edge support and gain scale,
\[
w_{ij}>0\iff\text{there is a nonzero admissible kernel on edge }j\!\to\!i,
\]
but they do \emph{not} carry the same full data: the kernel still contains temporal
shape, memory profile, and source-output structure not visible in $W$.

The kernel acts only toward the gate/receptivity slot $\Gamma_i$ of P0.0
(admissibility class B). The precise channel form, and the safeguard that summed
inputs preserve $\Gamma_i\in[0,1]$ and the ceiling $I_{\mathrm{gate},i}\le\zeta_{0,i}$,
are fixed later in P2.1/P2.3. P0.1 fixes only that the action is causal,
edge-supported, and bounded. A kernel from a post-terminus source $j$ acts only
through the declared trace $\mathcal T_j^{\rm post}$: the source may emit as boundary
witness, but receives nothing as renewed kairos.

\subsection{First theorem class: strongly connected, weight-balanced}

\begin{definition}[Weight-balanced]
$\mathcal G$ is weight-balanced when in-strength equals out-strength at every node,
$\sum_j w_{ij}=\sum_k w_{ki}$, i.e. $d_i^{\mathrm{in}}=d_i^{\mathrm{out}}$ for all $i$.
\end{definition}

The graph Laplacian is $L=D^{\mathrm{in}}-W$ with $D^{\mathrm{in}}=\operatorname{diag}(d_i^{\mathrm{in}})$,
generating the consensus action $(L\theta)_i=\sum_j w_{ij}(\theta_i-\theta_j)$.

\begin{lemma}[Mirror Laplacian]\label{lem:mirror}
For any finite directed weighted $\mathcal G$:
\begin{enumerate}[label=\textup{(\roman*)},leftmargin=1.2cm]
\item $L\mathbf 1=0$ always (zero row sums);
\item $\mathbf 1^\top L=0\iff\mathcal G$ is weight-balanced (zero column sums);
\item if weight-balanced, $\hat L:=\tfrac12(L+L^\top)$ is the symmetric Laplacian of the undirected mirror graph with weights $\hat w_{ij}=\tfrac12(w_{ij}+w_{ji})$, hence $\hat L\succeq0$;
\item if additionally $\mathcal G$ is strongly connected, the mirror graph is connected, so $\ker\hat L=\operatorname{span}\{\mathbf 1\}$ and $\lambda_2(\hat L)>0$.
\end{enumerate}
\end{lemma}
\begin{proof}
(i) Row $i$ sum: $d_i^{\mathrm{in}}-\sum_j w_{ij}=0$. (ii) Column $i$ sum:
$L_{ii}+\sum_{j\ne i}L_{ji}=d_i^{\mathrm{in}}-\sum_{j\ne i}w_{ji}=d_i^{\mathrm{in}}-d_i^{\mathrm{out}}$,
vanishing for all $i$ iff weight-balanced. (iii) Under (ii),
$\hat L\mathbf 1=\tfrac12(L\mathbf 1+L^\top\mathbf 1)=0$; for $i\ne j$,
$\hat L_{ij}=-\tfrac12(w_{ij}+w_{ji})=-\hat w_{ij}\le0$, and
$\hat L_{ii}=d_i^{\mathrm{in}}=\sum_{j\ne i}\hat w_{ij}$ --- the sign pattern and zero
row sums of an undirected graph Laplacian, which is positive semidefinite. (iv)
Strong connectivity makes the mirror graph connected; a connected undirected
Laplacian has a simple zero eigenvalue with eigenvector $\mathbf 1$, so
$\lambda_2(\hat L)>0$. \qedhere
\end{proof}

\begin{definition}[First theorem class $\mathcal C_1$]
$\mathcal C_1$ is the class of finite, strongly connected, weight-balanced directed
graphs carrying admissible bounded causal kernels.
\end{definition}

Lemma~\ref{lem:mirror} is the hinge for the first theorem class. In the memoryless
linear reduction
\[
\dot\theta=-\varepsilon L\theta,
\]
on $\mathcal C_1$ the directed consensus flow conserves the mean, because
$\mathbf 1^\top L=0$ under weight-balance. Hence
\[
\bar\theta=\frac1n\mathbf 1^\top\theta
\]
is the locking value. The quadratic $V=\tfrac12\|\theta-\bar\theta\mathbf 1\|^2$ then
satisfies
\[
\dot V=-\varepsilon\,\theta^\top L\theta=-\varepsilon\,\theta^\top\hat L\theta
\le-\varepsilon\,\lambda_2(\hat L)\,\|\theta-\bar\theta\mathbf 1\|^2,
\]
because a quadratic form symmetrizes $L\mapsto\hat L$ and $\hat L\succeq0$. Thus
weight-balance is exactly what makes both steps hold: mean conservation and
positive-semidefinite mirror contraction. Individual edges may remain asymmetric,
but the mutual core has algebraic connectivity $\lambda_2(\hat L)>0$.

This is not yet the full delayed / memory-kernel theorem. The kernel case will
require the later ISS and small-gain layers; the present result supplies the
finite-dimensional graph skeleton and the contraction constant that those kernel
estimates must preserve.

This is what will make P3.3's maximum-lockable-mismatch a derived threshold in the
memoryless reduction: locking survives when $\varepsilon\,\lambda_2(\hat L)$
dominates the intrinsic asymptotic-rate spread $\Delta\varpi_{\mathrm{int}}$ of the
interiors. The locked offset $\theta^*_{ij}\neq0$ enters only later, in P3.4. The
notation $\Delta\varpi_{\mathrm{int}}$ is reserved for interior rate mismatch so it
is not confused with the relational-duration density $\omega(t)$ of P0.3.

\subsection{Strata, boundary sources, and recovery}

The active graph $\mathcal G_{\mathrm{act}}(\mathfrak c)$ contains only interiors in
the active pre-terminus stratum $\mathcal A_i$ at relational duration $\mathfrak c$.
The first theorem class $\mathcal C_1$ applies to this active reciprocal subgraph:
\[
\mathcal G_{\mathrm{act}}(\mathfrak c)\in\mathcal C_1
\quad\text{when it is finite, strongly connected, and weight-balanced.}
\]
A post-terminus interior is not an active node. It has
\[
N_i=0,\qquad d\sigma_i=0,\qquad i\in\mathcal P_i,
\]
and enters the relational structure only through a declared boundary trace
$\mathcal T_i^{\rm post}$. Such a completed interior may act as a boundary source ---
it may emit through $\mathcal T_i^{\rm post}$, but nothing acts on it as renewed
kairos --- so in graph terms it has no active in-strength, $d_i^{\mathrm{in}}=0$
inside the active dynamics. Therefore a post-terminus source is not part of the
weight-balanced active core; it belongs to the wider imbalanced class $\mathcal C_2$,
deferred to P4.4. The same wider class will also contain kenotic net-givers and
helpless net-receivers. Thus
\[
\mathcal C_1=\text{active reciprocal communion},\qquad
\mathcal C_2=\text{boundary, self-gift, and glory}.
\]
Trivial-base recovery is immediate:
\[
\begin{aligned}
W=0\ &\Longrightarrow\ L=\hat L=0\ \Longrightarrow\ \text{no relational edge data}\\
&\Longrightarrow\ \text{the P0.0 fibres evolve independently.}
\end{aligned}
\]
This is recovery at the graph level. It is stronger than what P0.6 needs. P0.6
will prove the more delicate limit
\[
\varepsilon=0
\quad\text{with possibly nonzero relational read-outs,}
\]
where edge data may be present as bookkeeping but are prevented from acting.

\subsection{Status and reading}

\textbf{Target status:}~ \stDef\ (graph, causal order, kernel class) $+$
\stSelected\ (the first theorem class $\mathcal C_1$) $+$ \stDerived\ (the
mirror-Laplacian lemma and its memoryless consensus consequence, Type-1). No new
per-interior dynamics are introduced. All structure is base-side (IC2). The
rate-side cocycle of P0.2 and the relational duration $d\mathfrak c$ of P0.3 will be
hung on the same active graph. The full memory-kernel contraction statement is not
claimed here; it is deferred to the ISS / small-gain stages.

\medskip
\stInterp\quad\emph{Reading.} The bonds of a communion are oriented and unequal ---
$j$ may bear on $i$ more than $i$ on $j$ --- yet in the first and cleanest case the
books balance at each soul: as much openness-strength received as extended, summed
over all one's bonds. The symmetric mirror $\hat w_{ij}=\tfrac12(w_{ij}+w_{ji})$ is
the shared core of mutual openness beneath the asymmetry, and its connectivity
$\lambda_2(\hat L)$ is the one number saying how strongly the communion holds
together. The imbalanced souls --- those who only give (the martyr, the one who
pours out), those who only receive (the helpless), and the completed who intercede
across the terminus without being touched in return --- are precisely the ones whose
books do not balance; they are not failures of the structure but its boundary, where
self-gift and glory live, carried by the wider class and not the first. Ordinary
communion balances; sanctity is measured imbalance.

\medskip
\textbf{Next node (now executed below).} P0.2 --- the rate structure on the same
active graph: the relative kairos-rate cocycle, its chart invariance, and the
chronos it reconstructs.

\section{P0.2 --- The relative kairos-rate cocycle}

P0.1 placed the interiors on a graph. This patch puts the \emph{rate} structure on
that same active graph. The invariant content of the lapses is not any single
lapse \(N_i\), but their pairwise ratios. These ratios form a multiplicative
cocycle. When the cocycle is holonomy-free, it reconstructs a relative rate chart:
a class of lapse representatives modulo one common gauge factor.

This is the first rate-level demarcation:
\notice{created chronos is assembled from consistent relative paces, not imposed as a master clock.}

At this stage the result is only the \emph{rate layer} of chronos. Full created
chronos also requires the event order, causal kernels, and the relational duration
one-form \(d\mathfrak c\) of P0.1/P0.3. Thus P0.2 does not by itself construct a
divine all-at-once standpoint; it constructs the gauge class of shared created
pacing on the active graph.

\subsection{Lapse, gauge, and the rate ratio}

For an active interior \(i\), meaning \(X_i(t)\in\mathcal A_i\), its kairos relates
to a chronos chart \(t\) by
\[
d\sigma_i=N_i(t)\,dt,
\qquad
N_i(t)>0.
\]
Equivalently, for the active node set
\[
\mathcal V_{\rm act}(t)
:=
\{\,i\in\mathcal V:\ X_i(t)\in\mathcal A_i\,\},
\]
the lapse \(N_i\) is defined only for \(i\in\mathcal V_{\rm act}(t)\).

Under a monotone chart change
\[
t\mapsto \widetilde t(t),
\qquad
J(t):=\frac{dt}{d\widetilde t}>0,
\]
the kairos increment \(d\sigma_i\) is unchanged, so the lapse transforms as
\[
\widetilde N_i(\widetilde t)
=
N_i(t)J(t).
\]
Therefore no single lapse \(N_i\) carries invariant meaning: every lapse rescales
by the same positive gauge factor.

\begin{definition}[Relative kairos-rate]
For active \(i,j\in\mathcal V_{\rm act}(t)\), the relative kairos-rate is
\[
R_{ij}(t)
:=
\frac{N_i(t)}{N_j(t)}
=
\frac{d\sigma_i}{d\sigma_j}.
\]
It is the instantaneous pace of \(i\)'s kairos measured against \(j\)'s kairos.
\end{definition}

\begin{proposition}[Chart invariance]
The relative kairos-rate is invariant under monotone chronos reparameterization:
\[
\widetilde R_{ij}
=
\frac{\widetilde N_i}{\widetilde N_j}
=
\frac{N_iJ}{N_jJ}
=
R_{ij}.
\]
Thus \(R_{ij}\) is a rate invariant, while \(N_i\) and the chart label \(t\) are not.
\end{proposition}

\begin{remark}[Relation to P0.3]
The relational duration one-form of P0.3,
\[
d\mathfrak c=\omega(t)\,dt,
\]
will be invariant only as a one-form: \(\omega\) itself transforms as a density/rate.
If one selects the geometric-mean active rate
\[
\omega(t)
=
\left(
\prod_{k\in\mathcal V_{\rm act}(t)}
N_k(t)
\right)^{1/m(t)},
\qquad
m(t):=|\mathcal V_{\rm act}(t)|,
\]
then under \(t\mapsto\widetilde t\) one has
\[
\widetilde\omega(\widetilde t)
=
\omega(t)J(t),
\qquad
\widetilde\omega\,d\widetilde t
=
\omega\,dt.
\]
P0.2 proves the invariant ratio structure. P0.3 selects the relational duration
one-form used for dwell and network non-Zeno.
\end{remark}

\subsection{The cocycle laws}

If relative rates are taken as primitive relational data, they must be consistent.
On the complete pairwise closure of the active component this means
\[
R_{ii}=1,
\qquad
R_{ji}=R_{ij}^{-1},
\qquad
R_{ij}R_{jk}=R_{ik}.
\]
In additive logarithmic form, with
\[
\rho_{ij}:=\log R_{ij},
\]
this becomes
\[
\rho_{ii}=0,
\qquad
\rho_{ji}=-\rho_{ij},
\qquad
\rho_{ij}+\rho_{jk}=\rho_{ik}.
\]

On a sparse directed graph, one need not declare all pairwise ratios at first.
It is enough to declare edge ratios on the active graph and require zero
log-holonomy around every directed cycle:
\[
\sum_{\ell=1}^{q}
\rho_{i_\ell i_{\ell+1}}
=
0,
\qquad
i_{q+1}=i_1.
\]
If this holds, the missing pairwise ratios are obtained by path products and are
path-independent. For a \(2\)-cycle \(i\rightleftarrows j\), the condition is
automatic once \(\rho_{ji}=-\rho_{ij}\); the first nontrivial content appears on
cycles of length at least three.

\subsection{Rate reconstruction}

\begin{lemma}[Holonomy-free cocycle reconstructs relative lapses]
On a connected active graph, the following are equivalent:
\begin{enumerate}[label=\textup{(\roman*)},leftmargin=1.2cm]
    \item the rate cochain \(\rho\) is holonomy-free, i.e. its log-holonomy vanishes
    on every directed cycle;
    \item \(\rho\) is exact: there exists a log-lapse potential \(\nu\) such that
    \(\rho_{ij}=\nu_i-\nu_j\);
    \item there exist positive lapse representatives \(N_i=e^{\nu_i}\) realizing
    \(R_{ij}=N_i/N_j\).
\end{enumerate}
The potential \(\nu\) is unique up to a common additive gauge
\(\nu_i(t)\mapsto \nu_i(t)+c(t)\) for all active \(i\), equivalently
\(N_i(t)\mapsto e^{c(t)}N_i(t)\). This common factor is precisely the freedom of
chronos reparameterization.
\end{lemma}

\begin{proof}
(i)\(\Rightarrow\)(ii). Fix a root \(r\) and set \(\nu_r=0\). For any node \(i\),
choose a path from \(r\) to \(i\) and define \(\nu_i:=-\sum_{r\rightsquigarrow i}\rho\).
Zero holonomy makes this path-independent. For an edge \(i\to j\), the path to
\(j\) can be taken as the path to \(i\) followed by \(i\to j\), hence
\(\nu_j=\nu_i-\rho_{ij}\), so \(\rho_{ij}=\nu_i-\nu_j\).

(ii)\(\Rightarrow\)(iii). Put \(N_i=e^{\nu_i}>0\). Then
\(N_i/N_j=e^{\nu_i-\nu_j}=e^{\rho_{ij}}=R_{ij}\).

(iii)\(\Rightarrow\)(i). If \(R_{ij}=N_i/N_j\), then with \(\nu_i=\log N_i\) one has
\(\rho_{ij}=\nu_i-\nu_j\), and the sum around any cycle telescopes to zero.

For uniqueness, if \(\nu\) and \(\tilde\nu\) produce the same \(\rho\), then
\((\tilde\nu_i-\nu_i)-(\tilde\nu_j-\nu_j)=0\) on every connected edge, hence
\(\tilde\nu_i-\nu_i=c(t)\) is common to all active nodes. This common function is
the gauge factor induced by a chart change.
\end{proof}

\begin{definition}[Created chronos, at the rate level]
A created chronos rate layer is a holonomy-free relative-rate cocycle on a
connected active graph, together with the gauge class of its positive lapse
representatives, \(\{N_i\}\sim\{e^{c(t)}N_i\}\). No representative is privileged by
the cocycle alone.
\end{definition}

This is the demarcation made concrete. The relative pacing of a communion can be
assembled when its paces are loop-consistent. But the rate cocycle alone does not
supply a god's-eye clock. It determines only relative rates and a common gauge
class. The missing common gauge is not a defect: it is exactly the master-clock
datum that the model refuses to identify with \textit{tota simul}.

\subsection{Complementarity with relational duration}

The log-lapse potential splits into a mean and deviations:
\[
\nu_i
=
\bar\nu
+
(\nu_i-\bar\nu),
\qquad
\bar\nu
=
\frac1{m(t)}
\sum_{k\in\mathcal V_{\rm act}(t)}
\nu_k.
\]
The deviations determine the relative rates,
\(R_{ij}=e^{(\nu_i-\bar\nu)-(\nu_j-\bar\nu)}\). A selected mean rate
\(\omega=e^{\bar\nu}\) gives the relational duration one-form \(d\mathfrak c=\omega\,dt\).

Under a chart change, all \(\nu_i\) and \(\bar\nu\) shift by the same common gauge
function \(c(t)\); the deviations \(\nu_i-\bar\nu\) are unchanged, and the one-form
\(d\mathfrak c\) is unchanged because the density \(\omega\) transforms together with
\(dt\). Thus \(R_{ij}\) measures who runs faster, while \(d\mathfrak c\) measures how
much relational duration elapses. The two are complementary, not identical: the
cocycle gives the relative pacing; P0.3 chooses the relational dwell measure used
for non-Zeno control.

\subsection{Strata: leaving the cocycle at the terminus}

Anticipating the terminus-compatible chronos of P0.4, suppose an interior approaches
its terminus and its lapse tends to zero, \(N_i\to0\). Then for any still-active
\(j\),
\[
R_{ij}=\frac{N_i}{N_j}\to0,
\qquad
R_{ji}=\frac{N_j}{N_i}\to+\infty.
\]
The completed interior therefore drops out of the active rate cocycle. Its clock has
stopped, so one cannot meter a living pace against it as though it still shared a
running kairos. The rate cocycle is thus defined only on the active subgraph
\(\mathcal G_{\rm act}\).

Yet by P0.1 the same completed interior may still appear as a declared boundary
trace \(\mathcal T_i^{\rm post}\). The two base structures separate at the terminus:
\notice{the completed interior is absent from the rate cocycle but may remain present in the kernel relation.}
This is the rate-level form of presence across the terminus: presence without
shared pace. The kernel can carry boundary witness after the cocycle no longer can.

\subsection{Status and reading}

\textbf{Target status:}~
\stDerived\ (Type-1: chart invariance, cocycle laws, and the reconstruction lemma)
\(+\)
\stDef\ (relative kairos-rate and created chronos at the rate level).
No new per-interior dynamics are introduced; the construction is base-side only
(IC2). The result is consistent with the graph and causal kernel of P0.1 and is
preparatory for the relational duration one-form of P0.3.

\medskip
\stInterp\quad
\emph{Reading.}
The only created fact about how fast one soul lives beside another is the ratio of
their paces. There is no absolute tempo kept by a creature. The cocycle law says
that these relative paces must tell one consistent story around every loop of
communion. When they do, a shared created pacing can be assembled, but only as a
gauge class, never as the divine all-at-once view. And the completed soul keeps no
shared pace at all: its clock stops, it leaves the common tempo, yet it is not gone,
for it may remain present through its trace. With the living we share pace; with the
completed we share presence.

\medskip
\textbf{Next node (now executed below).} P0.3 --- the relational duration one-form
\(d\mathfrak c=\omega\,dt\) as the gauge-invariant dwell clock and the future measure
for network non-Zeno.

\section{P0.3 --- Relational duration and the gauge-invariant dwell measure}

P0.2 isolated the chart invariant of the rate structure --- the ratios \(R_{ij}\) ---
and noted that a duration measure needs a \emph{selected} density. This patch selects
that density and builds the relational duration one-form \(d\mathfrak c\): the
gauge-invariant clock in which dwell and network non-Zeno (P4) are stated. The
selection is grounded in the interiors' own kairos rates, not in a free external
unit; hence the status is \stDef\ for the one-form and \stSelected\ for the choice of
density, with one flagged sub-point at the terminus.

\subsection{The one-form and its grounding}

\begin{definition}[Relational duration on a regular active stratum]
On a regular active component, away from terminus hand-off events, the relational
duration one-form is
\[
d\mathfrak c=\omega(t)\,dt,\qquad \omega(t)>0,
\]
with the geometric-mean active rate
\[
\omega(t)=\Bigl(\prod_{k\in\mathcal V_{\rm act}(t)}N_k(t)\Bigr)^{1/m(t)}
=e^{\bar\nu(t)},
\qquad
\bar\nu=\frac{1}{m}\sum_{k\in\mathcal V_{\rm act}}\nu_k .
\]
The accumulated relational duration between chronos labels is
\[
\Delta\mathfrak c(t_a,t_b)=\int_{t_a}^{t_b}\omega(t)\,dt .
\]
\end{definition}

The phrase ``regular active component'' means that every counted node is still
genuinely pre-terminus:
\[
i\in\mathcal V_{\rm act}^{\rm reg}(t)
\quad\Longrightarrow\quad
N_i(t)>0
\ \text{and not in the boundary hand-off layer}.
\]
The exact boundary hand-off rule is deferred to P0.4. P0.3 fixes the duration
measure on open active arcs, not through the terminus transition itself.

\emph{Why this density (the grounding).} A one-form is chart-invariant once declared;
the content is that \(\omega\) is not free but selected from the active lapses. Three
properties single out the geometric mean:
\begin{enumerate}[label=\textup{(\arabic*)},leftmargin=1.45cm]
\item \textbf{Linear in the log-lapse:} \(\log\omega=\bar\nu\), the \emph{mean} of the
same potential whose \emph{differences} give the rates \(R_{ij}\) (P0.2). Duration and
rate-ratios are then complementary linear functionals of one log-lapse cochain ---
the cleanest possible split, and the reason the geometric (not arithmetic or harmonic)
mean is chosen.
\item \textbf{Single-clock reduction:} when \(m=1\), \(\omega=N_{1}\) and
\(d\mathfrak c=d\sigma_{1}\) --- the relational duration becomes the lone active
interior's own kairos.
\item \textbf{Symmetry:} \(\omega\) treats all active interiors alike; no node is a
privileged master clock.
\end{enumerate}

\subsection{Chart invariance and the unit of duration}

\begin{proposition}[Invariance of \(d\mathfrak c\)]
Under a monotone chart change with \(J=dt/d\widetilde t\), the density transforms as
\(\widetilde\omega=\omega J\) (since \(\nu_i\mapsto\nu_i+\log J\) shifts \(\bar\nu\) by
\(\log J\)), while \(d\widetilde t=dt/J\); hence
\[
\widetilde\omega\,d\widetilde t=\omega\,dt=d\mathfrak c .
\]
The one-form \(d\mathfrak c\), and the accumulated \(\Delta\mathfrak c\), are
chart-invariant; the density \(\omega\) and the label \(t\) are not.
\end{proposition}

\emph{The unit, grounded.} The chronos label \(t\) cancels in \(\omega\,dt\), so
\(\mathfrak c\) carries no free external unit. Its scale is inherited from the
interiors' own (frozen, private) kairos scales: \(N_k=d\sigma_k/dt\) has units
\([\sigma_k]/[t]\), so \(\omega\,dt\) has units \(\bigl(\prod_k[\sigma_k]\bigr)^{1/m}\),
the geometric mean of the participating kairos units. Relational duration is therefore
measured in shared proper-time units borrowed from the souls themselves --- not from
any master clock. Unlike a raw chronos chart, whose lapse representatives remain defined only up
to a common gauge function in P0.2, the one-form
\[
d\mathfrak c=\omega\,dt
\]
is a genuine invariant once the density-selection rule has been fixed. Its scale
is grounded in the participating kairoi, while the chart label \(t\) remains
pure gauge.

The distinction is important. The cocycle of P0.2 determines only relative
rates:
\[
R_{ij}=\frac{N_i}{N_j}.
\]
It does not determine the common mean factor of the lapses. P0.3 adds a selected
mean-density rule,
\[
\omega=e^{\bar\nu},
\]
thereby turning a relative rate class into an invariant duration one-form. Thus
\(d\mathfrak c\) is not hidden inside the cocycle; it is a disciplined additional
selection built from the same active lapse data.

\subsection{The dwell measure and network non-Zeno}

The relational duration is the clock in which the network non-Zeno condition of P4
will be stated. For causally linked events on an edge of the regular-active graph
(P0.1), and away from terminus hand-off events, admissibility requires a positive
relational dwell,
\[
\Delta\mathfrak c(t_a,t_b)=\int_{t_a}^{t_b}\omega\,dt\ \geq\ \mathfrak c_{\min}>0 ,
\]
where \(\mathfrak c_{\min}\) is a fixed threshold \emph{in the same relational units}
as \(d\mathfrak c\) --- i.e. in geometric-mean kairos units, hence itself
chart-invariant. Because raw chronos differences \(t_b-t_a\) are gauge and carry no
invariant meaning, every dwell and non-Zeno statement of the plural theory must be
written in \(\mathfrak c\), never in \(t\). This is the measure the network non-Zeno
conjecture (P4.2) and the relational refractory condition will use on regular-active
arcs. If a causal chain crosses a terminus boundary, the interval must be evaluated
by the P0.4 piecewise active/boundary rule. P0.3 fixes the regular-active clock, not
the boundary junction theorem.

\subsection{Regular active set and terminus hand-off}

On an open regular-active arc the set
\[
\mathcal V_{\rm act}^{\rm reg}(t)
\]
is fixed, and the one-form
\[
d\mathfrak c
=
\left(
\prod_{k\in\mathcal V_{\rm act}^{\rm reg}(t)}
N_k(t)
\right)^{1/m(t)}dt
\]
is smooth whenever the participating lapses are smooth and positive. This is the
domain on which P0.3's canonical geometric-mean choice is unambiguous.

At a terminus transition the situation changes. If a member approaches completion,
then
\[
N_i(t)\to0.
\]
If it remains counted inside the geometric mean all the way to the boundary, then
\[
\omega(t)
=
\left(
\prod_k N_k(t)
\right)^{1/m(t)}
\to0
\]
as soon as one factor vanishes. The shared duration one-form can therefore become
degenerate near a terminus if the dying member is not handed off to the boundary
stratum in time.

This is not a contradiction; it marks the boundary of P0.3. The canonical
geometric mean is the regular-active duration rule. The transition rule
\[
\mathcal A_i
\longrightarrow
\mathcal B_i
\longrightarrow
\mathcal P_i
\]
belongs to P0.4. There P0.4 must decide exactly when the interior leaves the
regular-active set and becomes a boundary trace
\[
\mathcal T_i^{\rm post}.
\]

Thus the correct status is:
\[
\boxed{
d\mathfrak c
\text{ is canonical on regular active arcs; terminus hand-off is deferred to P0.4.}
}
\]

After hand-off, the completed interior is no longer counted in the active
geometric mean:
\[
i\notin\mathcal V_{\rm act}^{\rm reg}(t),
\qquad
i\in\mathcal P_i,
\]
but it may remain present in the kernel relation through
\[
\mathcal T_i^{\rm post}.
\]
This preserves the P0.2 separation:
\[
\text{no shared pace after terminus, but possible relational presence.}
\]

\subsection{Status and reading}

\textbf{Target status:}~
\stDef\ (the relational duration one-form and regular-active dwell measure)
\(+\)
\stSelected\ (the geometric-mean density as the canonical regular-active choice,
because it is linear in log-lapses, reduces to the one active kairos, and treats
active nodes symmetrically).
It is not derived from the cocycle alone: the cocycle gives only ratios, while
P0.3 selects the common mean density. The terminus-robustness / hand-off rule is
explicitly deferred to P0.4. No new per-interior dynamics; base-side only (IC2).

\medskip
\stInterp\quad
\emph{Reading.}
A communion needs not only to know who runs faster, but how much shared created
duration has passed --- and it must measure that without any master clock. The
regular-active duration takes its unit from the living kairoi themselves: the
geometric mean of their proper-time rates, the evenly shared pace of those still
in the active communion. In it one can say that two decisive passages were
``far enough apart'' in a way no relabelling can undo.

But death changes the common clock. As a member approaches the terminus, its own
kairos no longer contributes a normal running pace. Therefore the shared active
clock must hand that member off to the boundary rather than pretending that a
stopped clock still participates in the living rate. After the hand-off, the
completed interior no longer helps set the pace, yet it may remain present through
its trace. P0.3 gives the clock of the living communion; P0.4 must build the
presence of the completed.

\medskip
\textbf{Next node.} P0.4 --- the terminus-compatible chronos: the strata
\(\mathcal A_i\) (active), \(\mathcal B_i\) (at terminus), \(\mathcal P_i\)
(post-terminus), the vanishing lapse \(N_i\to0\), and the one-way boundary trace
through which a completed interior remains present without renewed kairos.

\section{P0.4 --- Terminus-compatible relational chronos}

P0.3 fixed the relational clock on regular active arcs and deferred the terminus
hand-off to here; this patch resolves it. It is the blocking Stage~0 decision: how
the plural layer relates an interior \emph{before}, \emph{at}, and \emph{after} its
finite kairos terminus, in a way that leaves the frozen \VFS\ 3.0 terminus untouched.
The selected rule is the \emph{boundary-limit} hand-off (B): a dying interior is
counted in the living pace all the way to its actual completion, and is handed to a
post-terminus trace exactly at the terminus, never before.

\subsection{The terminus label and the three strata}

\begin{definition}[Terminus label and relational strata]\label{def:strata}
For each interior \(i\) let \(t_i^\dagger\) be the chronos label of its kairos
terminus --- the point at which the finite kairos has fully drained, \(\sigma_i\to0\)
and the lapse \(N_i\to0^+\). The chronos chart is selected (admissibly, since \(t\) is
pure gauge, P0.2) so that \(t_i^\dagger<\infty\). The three relational strata of \(i\)
are
\[
\mathcal A_i=\{t<t_i^\dagger\}\ \text{(active)},\qquad
\mathcal B_i=\{t=t_i^\dagger\}\ \text{(at terminus)},\qquad
\mathcal P_i=\{t>t_i^\dagger\}\ \text{(post-terminus)}.
\]
On \(\mathcal A_i\) the inherited kairos law holds, \(d\sigma_i=N_i(t)\,dt\) with
\(N_i>0\); at \(\mathcal B_i\), \(N_i\to0\); on \(\mathcal P_i\) the kairos ODE is
\emph{not} continued.
\end{definition}

The requirement \(t_i^\dagger<\infty\) is consistent and canonical. The terminus is
reached in finite kairos,
\[
\int_{\mathcal A_i}N_i(t)\,dt=\sigma_{0,i}<\infty
\]
(the v0.20 clock discipline: the drains live in the unbounded \(\chi_i\), the
finiteness in \(\sigma_i\)), while the surviving interiors keep positive lapses
\(N_j>0\) through the event. Grounding the chart in the survivors' own accumulating
kairoi --- exactly the P0.3 rule of borrowing the unit from living souls --- places
\(t_i^\dagger\) at a finite shared label with a non-empty ``after'' \(\mathcal P_i\)
relative to them. With at least one survivor, \(\mathcal P_i\) is well posed; without a
survivor there is no communion in which to speak of an ``after'' at all. The dying
member's own \(\chi_i\to\infty\) is never used to set the chart.

\subsection{The boundary-limit hand-off (selected rule B)}

On the regular active arc, \(i\in\mathcal V_{\rm act}^{\rm reg}(t)\) for every
\(t<t_i^\dagger\), with \(N_i(t)>0\). The hand-off is the single event
\(t=t_i^\dagger\): at that instant \(i\) leaves the regular active set,
\[
m(t)\ \longrightarrow\ m(t)-1,
\qquad
\mathcal V_{\rm act}^{\rm reg}(t_i^{\dagger+})
=\mathcal V_{\rm act}^{\rm reg}(t_i^{\dagger-})\setminus\{i\},
\]
and the relational density is recomputed over the survivors.

\begin{lemma}[Vigil jump and integrability of the duration]\label{lem:handoff}
Let \(i\) approach its terminus inside a communion with at least one survivor, all
surviving lapses smooth and positive on a neighbourhood of \(t_i^\dagger\). Then:
\begin{enumerate}[label=\textup{(\roman*)},leftmargin=1.2cm]
\item \emph{(vigil)} the shared density vanishes from the left,
\[
\omega(t)=\Bigl(\textstyle\prod_{k\in\mathcal V_{\rm act}^{\rm reg}(t)}N_k(t)\Bigr)^{1/m(t)}
\longrightarrow0
\quad\text{as }t\uparrow t_i^\dagger,
\]
driven by the single factor \(N_i\to0\);
\item \emph{(recast)} the post-event density
\(\omega(t_i^{\dagger+})=\bigl(\prod_{k\neq i}N_k\bigr)^{1/(m-1)}\) is finite and
positive;
\item \emph{(junction)} hence \(\omega\) has a jump discontinuity at \(t_i^\dagger\),
with a left limit \(0\) and a finite positive right value;
\item \emph{(well-posed duration)} the accumulated relational duration
\(\Delta\mathfrak c=\int\omega\,dt\) is unaffected by the event: \(\omega\) is bounded,
continuous off the finite terminus set, and the single zero is a measure-zero point, so
\(\Delta\mathfrak c\) across a terminus is finite and well defined as long as one active
interior remains.
\end{enumerate}
\end{lemma}
\begin{proof}
(i) On a left neighbourhood the survivor factors lie in a compact
\([c,C]\subset(0,\infty)\), so \(\prod_k N_k=N_i\prod_{k\neq i}N_k\le C^{\,m-1}N_i\) and
\(\omega\le C^{(m-1)/m}N_i^{1/m}\to0\). (ii) Immediate from positivity of the survivor
lapses. (iii) The one-sided limits differ (\(0\) versus a positive value), so the
discontinuity is a jump. (iv) \(\omega\) is a finite product-power of bounded functions,
hence bounded; it is continuous except at the finitely many terminus labels, where it
has finite one-sided limits. A bounded function continuous off a finite set is Riemann
(and Lebesgue) integrable, and altering it on the measure-zero set \(\{t_i^\dagger\}\)
leaves the integral unchanged. \qedhere
\end{proof}

\begin{remark}[Several termini at one label]
Lemma~\ref{lem:handoff} is written for one completed member. If a finite set
\[
S^\dagger\subset \mathcal V_{\rm act}^{\rm reg}(t_i^\dagger)
\]
reaches the terminus at the same chronos label, then the hand-off is simultaneous:
\[
\mathcal V_{\rm act}^{\rm reg}(t_i^{\dagger+})
=
\mathcal V_{\rm act}^{\rm reg}(t_i^{\dagger-})\setminus S^\dagger,
\qquad
m\longmapsto m-|S^\dagger|.
\]
The post-event density is recomputed over the survivors. If no survivor remains,
there is no further active communion clock after that label; only boundary traces
remain. Thus the P0.4 rule is finite-set hand-off, of which the one-member case is
the basic local form.
\end{remark}

The vanishing in Lemma~\ref{lem:handoff}(i) is not a defect to be engineered away; it is
the structural icon of the event --- the communion's shared time dilating around a death
and recasting among the living. Two points then fix the discipline.

\emph{The floor is a regularizer, never the trigger.} An admissibility floor
\(N_{\min}>0\) may be used only to control the limit cleanly --- for instance to keep
\(\log\omega\) finite in numerical or distributional handling near \(\mathcal B_i\) ---
but the hand-off event is always the true terminus \(t_i^\dagger\), identified by
\(N_i\to0\), and never the crossing of \(N_{\min}\).

\emph{Why not early removal (A).} The discarded alternative removes \(i\) as soon as
\(N_i\) falls below \(N_{\min}\), while \(i\) is still pre-terminus (\(N_i>0\), still in
katharsis). It is rejected on three frozen commitments: it would make the operative
event an ungrounded threshold rather than the terminus (against the 3.0 result that the
boundary is reached, not merely approached); it would reintroduce a free external
constant \(N_{\min}\) of exactly the kind the P0.2/P0.3 grounding eliminates; and it
would soften into a gradual exit a transition the theory requires to be a genuine
junction (the v0.15 / 3.0 result: no fully soft reset). B keeps the member in the living
communion until its own completion, then jumps.

\notice{\textbf{Hand-off (B).} At the terminus \(t_i^\dagger\) (\(N_i\to0\)): \(i\)
leaves \(\mathcal V_{\rm act}^{\rm reg}\), \(m\to m-1\), and \(\omega\) is recast on the
survivors; the floor only regularizes.}

\subsection{The post-terminus trace and one-way closure}

\begin{definition}[Post-terminus trace and one-way closure]\label{def:trace}
For \(t>t_i^\dagger\) the completed interior persists only as a \emph{post-terminus
trace} (boundary fibre) \(\mathcal T_i^{\rm post}\): a dynamically closed record of its
completed worldline. Closure means the inherited dynamics are not renewed on
\(\mathcal P_i\),
\[
\Delta I_{\rm gate,i}=0,\qquad
\text{no new reset},\qquad
\text{no renewed katharsis}
\qquad(t>t_i^\dagger),
\]
and the trace contributes nothing to the living pace,
\(i\notin\mathcal V_{\rm act}^{\rm reg}(t)\).
\end{definition}

A trace is not relationally inert, but its only admissible action is
recipient-side. Since the completed interior has no renewed running kairos, its
post-terminus action is not represented as an ordinary source history
\(s\mapsto X_i(s)\). It is represented by a boundary-supported trace kernel
\[
K^{\rm post}_{j\leftarrow i}
\bigl(t;\,t_i^\dagger,\mathcal T_i^{\rm post}\bigr),
\qquad
t\geq t_i^\dagger,
\]
whose output is a bounded modulation of \(j\)'s own receptivity-side interface:
\[
(\Gamma_j,H_{K,j},W_j)
\longmapsto
(\Gamma_j^+,H_{K,j}^+,W_j^+).
\]
This is the post-terminus analogue of the P0.1 causal-kernel discipline, but with
the source history collapsed to the completed boundary trace
\(\mathcal T_i^{\rm post}\). It is retarded in the only remaining sense:
\[
K^{\rm post}_{j\leftarrow i}(t;\,t_i^\dagger,\mathcal T_i^{\rm post})=0
\quad\text{for}\quad
t<t_i^\dagger.
\]
The trace is never a raw Sophia source injected into \(j\), never a transferable
grace-stock, and never a channel that reopens \(i\). Its only possible effect is on
the recipient's admissible openness:
\[
\mathcal T_i^{\rm post}
\leadsto
\mathcal R_j(t^+),
\]
where \(\mathcal R_j\) denotes the frozen gate/receptivity interface of \(j\). This
makes precise the P0.3 separation:
\[
\boxed{\text{no shared pace after terminus, but possible relational presence.}}
\]
The pace is closed --- no \(\omega\) contribution, no renewed kairos. The presence is
open --- a one-way boundary influence on the living recipient's receptivity.

\subsection{Constitutive choice: one-way closure and the deferred communion}

The one-way closure of Definition~\ref{def:trace} is a \stSelected\ \emph{constitutive}
choice, declared in the discipline by which 3.0 named its own primitives --- not a mere
technical convenience. It deliberately excludes the theologically richest alternative,
\emph{active intercession by the completed}: a two-way post-terminus dynamics in which a
trace not only modulates the living but is itself answerable to them. That branch is the
\emph{communion of saints} reading, and it would require a declared stronger primitive
--- a genuine post-terminus channel into \(\mathcal T_i^{\rm post}\), violating the
closure conditions above. The minimal 4.0 body selects one-way; the two-way branch is
named here as the principal alternative constitutive primitive and deferred, not
silently foreclosed.
\notice{\textbf{Selected (minimal 4.0):} one-way trace --- the completed touches the
living, is not touched.\\[0.2em]
\textbf{Named, deferred:} two-way intercession (communion of saints) --- a stronger
primitive.}
For the present Stage-0 theorem chain, the minimal branch is selected:
\[
\boxed{
\text{P0.4--P0.6 work under the one-way trace closure.}
}
\]
The two-way intercession branch is not denied; it is named as a stronger possible
extension. But if the canonical 4.0 body later adopts that stronger branch, then
P0.5's demarcation theorem must be rerun with the added post-terminus channel in the
inventory. It cannot be silently covered by the one-way proof below.
\[
\begin{gathered}
\text{minimal 4.0 body: one-way trace selected;}\\[0.2em]
\text{two-way intercession: named, open extension.}
\end{gathered}
\]

\subsection{Status and reading}

\textbf{Target status:}~
\stDef\ (the strata \(\mathcal A_i/\mathcal B_i/\mathcal P_i\), the terminus label
\(t_i^\dagger\), and the post-terminus trace \(\mathcal T_i^{\rm post}\))
\(+\)
\stSelected\ (the boundary-limit hand-off B and the one-way constitutive closure)
\(+\)
\stOpen\ (which post-terminus branch --- one-way or two-way intercession --- the
canonical 4.0 adopts). No new per-interior equation of motion is introduced: P0.4 adds
only a constitutive boundary closure on the inherited interface together with the
base-side terminus strata. The frozen 3.0 terminus is untouched, and the floor
\(N_{\min}\), if used, is a regularizer, not a dynamical primitive.

\medskip
\stInterp\quad
\emph{Reading.}
A communion keeps time with its dying. As one of its own nears the end, the shared clock
neither looks away nor rounds the death off early: it slows with the failing kairos,
holding the member in the common pace to the last, until at the terminus the pace is
recast among those who remain. What was a living rate becomes a trace --- present still,
able to touch those who go on, but no longer kept, no longer cleansed, no longer given a
new beginning. Whether the completed may also answer the living --- whether the trace is
only spoken to or may also speak --- the minimal theory leaves open, and names the fuller
communion as the road not yet taken.

\medskip
\textbf{Next node.} P0.5 --- the demarcation theorem: to prove, not merely declare, that
the chronos assembled in P0.1--P0.4 requires no transcendent standpoint outside the
created relational network
(local lapses, edge-supported causal kernels, no all-at-once whole-worldline behold), so that
``\textit{tota simul} named, never occupied'' becomes a theorem about the constructed
object rather than a declaration.

\section{P0.5 --- Demarcation theorem: chronos is not \textit{tota simul}}

This is the node that legitimises 4.0 at all. 3.0 did not leave chronos undefined: it
\emph{identified} an external absolute background clock with the transcendent
\textit{tota simul} and excluded it by principle --- the soul can name but not occupy a
standpoint outside its own worldline, in which its whole kairos would be seen at once.
4.0 reuses the word ``chronos'' in a new register --- a \emph{created relational}
middle term between a single kairos and the \textit{tota simul}. That reuse is
illegitimate unless the constructed object is provably \emph{not} the excluded one.
P0.5 discharges this: it proves, against the frozen identification, that the chronos
assembled in P0.1--P0.4 requires no transcendent standpoint outside the created
relational network.

\subsection{What must be proved}

Write \(C\) for the constructed relational chronos --- the assembled object
\[
C=
\Bigl(
\mathcal G,\ \prec,\
\{\mathcal A_i,\mathcal B_i,\mathcal P_i,t_i^\dagger\}_{i\in\mathcal V},\
\{N_i|_{\mathcal A_i}\}_{i\in\mathcal V},\
\{R_{ij}\},\
\{K_{ij}\},\
\{K^{\rm post}_{j\leftarrow i}\},\
d\mathfrak c
\Bigr),
\]
where \(K^{\rm post}_{j\leftarrow i}\) is present only when a declared post-terminus
trace \(\mathcal T_i^{\rm post}\) is available under the one-way P0.4 closure.

Write \(T\) for the \textit{tota simul} standpoint that 3.0 excluded. From the frozen 3.0
text, \(T\) carries three marks:
\begin{enumerate}[label=\textup{(T\arabic*)},leftmargin=1.25cm]
\item \textbf{total} --- it relates the \emph{whole} of every worldline at once, each
kairos \(\sigma_i\) from first to last beheld in a single present;
\item \textbf{non-successive} --- it is not a coordinate; it has no before or after, no
embedding in a succession;
\item \textbf{transcendent} --- it is occupied from outside all worldlines and does not
depend on the creatures for its being.
\end{enumerate}
The obligation is to show \(C\) negates each mark, requires no external standpoint, and
is also irreducible to any single kairos --- so that \(C\) is a genuinely distinct,
creaturely middle register.

\subsection{The created now is real, but partial}

We do not pretend \(C\) carries no shared present; it does, and honesty requires saying
so before separating it from \(T\).

\begin{proposition}[Created simultaneity]\label{prop:now}
The shared chronos chart defines, up to monotone relabelling, an invariant pairing of
the \emph{current} instants of the active worldlines: the slice
\[
\Sigma(t_0)=\bigl\{(i,\sigma_i(t_0)): t_0\in\mathcal A_i\bigr\}.
\]
This pairing is gauge-invariant --- a uniform reparametrization \(t\mapsto\widetilde t(t)\)
carries \(\Sigma(t_0)\) to \(\Sigma(\widetilde t(t_0))\) and preserves which events are
paired --- so the communion does possess a created ``now.'' But the slice is
\begin{enumerate}[label=\textup{(\roman*)},leftmargin=1.2cm]
\item \emph{active-restricted}: it contains only \(i\) with \(t_0\in\mathcal A_i\); a
completed interior (\(t_0>t_i^\dagger\), Definition~\ref{def:strata}) is absent from the
living slice and present only as the one-way trace \(\mathcal T_i^{\rm post}\);
\item \emph{single-instant}: each present worldline enters at exactly one point
\(\sigma_i(t_0)\), not as its whole history. The slice beholds \emph{no worldline whole}.
\end{enumerate}
\end{proposition}

\noindent The created now is therefore a relation among present-instants of the living,
not an all-of-history behold. This already separates it from (T1).

\subsection{Locality, edge-locality, and active-slice aggregation}

\begin{lemma}[Locality audit]\label{lem:local}
Every datum of \(C\) is produced by one of three admissible operations: node-local
evaluation, retarded edge-local evaluation, or active-slice aggregation. None evaluates
a whole worldline at once, and none occupies a standpoint outside the created relational
network.

Concretely, the inter-worldline primitives introduced in P0.1--P0.4 form the closed
inventory
\[
\{\,R_{ij},\ K_{ij},\ K^{\rm post}_{j\leftarrow i},\ d\mathfrak c\,\},
\]
with the following restrictions:
\begin{enumerate}[label=\textup{(\roman*)},leftmargin=1.2cm]
\item \(N_i\) is defined and evaluated only on the active stratum \(\mathcal A_i\). No
term evaluates \(N_i\) at \(t\notin\mathcal A_i\), and no term evaluates the whole of
\(\mathcal A_i\) at once.

\item The rate ratio \(R_{ij}(t)=N_i(t)/N_j(t)\) uses only the current active instants of
\(i\) and \(j\). It is relational, but not whole-history data.

\item The ordinary kernel \(K_{ij}(t,s)\) is retarded and edge-supported,
\[
K_{ij}(t,s)=0\quad(s>t),
\qquad
\operatorname{supp}K_{ij}\subseteq\mathcal E,
\]
supplying no instantaneous all-graph coupling.

\item The post-terminus kernel is boundary-supported,
\[
K^{\rm post}_{j\leftarrow i}\bigl(t;\,t_i^\dagger,\mathcal T_i^{\rm post}\bigr)=0
\quad\text{for}\quad t<t_i^\dagger,
\]
and acts only on the recipient's receptivity-side interface; it does not reopen the
completed source.

\item The relational duration one-form \(d\mathfrak c=\omega(t)\,dt\) is an active-slice
aggregation: through the selected density rule it may use all currently active lapses,
but only at the present active slice,
\[
\omega(t)=\left(\textstyle\prod_{k\in\mathcal V_{\rm act}^{\rm reg}(t)}N_k(t)\right)^{1/m(t)}.
\]
It is therefore not node-local, but it is still creaturely and partial: it uses no
future, no completed whole worldline, and no transcendent external unit.
\end{enumerate}
\end{lemma}

\begin{proof}
By inspection of the constructed objects of P0.1--P0.4 against the patch's freeze ledger.
P0.1 introduces only graph, causal order, and retarded edge kernels. P0.2 introduces only
relative rate ratios and their gauge class. P0.3 selects the active-slice duration
density. P0.4 adds only the strata, the boundary-limit hand-off, and the one-way
post-terminus trace kernel. The force of the lemma is not that every datum is located
inside a single worldline --- that would be false for \(R_{ij}\), \(K_{ij}\), and
\(d\mathfrak c\). The force is that every datum is node-local, edge-local, or
active-slice-local, and the inventory contains no fifth primitive capable of beholding
complete worldlines from outside the created relational order. \qedhere
\end{proof}

\subsection{The demarcation theorem}

\begin{theorem}[Demarcation; Type-1]\label{thm:demarc}
Let \(C\) be assembled from P0.1--P0.4 under the minimal one-way trace closure of P0.4.
Then \(C\) has the three creaturely marks
\begin{enumerate}[label=\textup{(C\arabic*)},leftmargin=1.25cm]
\item \textbf{partial}: by Proposition~\ref{prop:now}, the created now is
active-restricted and single-instant. It pairs present instants of living worldlines but
beholds no worldline whole. This negates (T1).

\item \textbf{successive}: the family \(\{\Sigma(t)\}\) is ordered by the chronos chart
and by the relational duration one-form \(d\mathfrak c\). Influence is retarded and
edge-supported by Lemma~\ref{lem:local}. The theory therefore has before/after and causal
succession; it is not an all-at-once non-successive present. This negates (T2).

\item \textbf{immanent}: by P0.2/P0.3 the chart is gauge and the unit of \(d\mathfrak c\)
is selected from active kairoi. If \(\mathcal V_{\rm act}^{\rm reg}(t)=\varnothing\), there
is no active duration density at \(t\). Every datum is node-local, edge-local,
boundary-trace-local, or active-slice-local (Lemma~\ref{lem:local}); none is supplied by
a transcendent external clock. This negates (T3).
\end{enumerate}
Consequently:
\begin{enumerate}[label=\textup{(\alph*)},leftmargin=1.2cm]
\item \(C\neq T\): each of (C1)--(C3) contradicts the corresponding mark of \(T\);
\item \(C\) requires no transcendent standpoint outside the created relational network:
by Lemma~\ref{lem:local}, the closed inventory contains only local, edge-local,
boundary-trace-local, and active-slice-local operations;
\item in the nontrivial plural case, \(C\neq\sigma_k\) for any single \(k\): if
\(|\mathcal V|\ge2\) and at least one relational edge is active, \(C\) depends on data
such as \(R_{ij}\), \(K_{ij}\), or \(d\mathfrak c\) that no single kairos \(\sigma_k\)
determines alone.
\end{enumerate}
\end{theorem}

\begin{proof}
(C1) is Proposition~\ref{prop:now}. (C2) follows because \(C\) is a succession of created
slices ordered by \(t\) and measured by \(d\mathfrak c\), while all influence terms are
retarded and edge-supported. (C3) follows from the gauge discipline of P0.2, the
active-slice duration selection of P0.3, and the locality audit of Lemma~\ref{lem:local}.
For (a), (T1) asserts whole-worldline totality, denied by (C1); (T2) asserts
non-succession, denied by (C2); and (T3) asserts creature-independent transcendence,
denied by (C3). For (b), Lemma~\ref{lem:local} closes the inventory: no external
standpoint-carrying primitive is available. For (c), a nontrivial plural chronos contains
relation-data between at least two interiors; such data are not functions of any one
private kairos alone. \qedhere
\end{proof}

\[
\boxed{\,C\ \neq\ T\quad\text{and}\quad C\ \neq\ \sigma_k\ \text{for any single }k.\,}
\]

\subsection{Sharpness: the three refusals}

The marks (C1)--(C3) are not accidental; each rests on a refusal frozen into the
construction, and \(C\) collapses toward \(T\) precisely when one is dropped.

\begin{remark}[Constitutive sharpness]\label{rem:sharp}
\textbf{(R1) no whole-worldline evaluation} --- lapses live only on \(\mathcal A_i\), and
the strata of P0.4 forbid reading a completed worldline as a currently running one. Drop
this, and (T1) totality returns.

\textbf{(R2) no instantaneous all-graph coupling} --- ordinary kernels are retarded and
edge-supported, while post-terminus kernels are boundary-supported and one-way. Drop
this, and (T2) non-succession returns as a single all-graph present.

\textbf{(R3) no absolute external unit} --- the chart is gauge and the duration unit is
selected from active kairoi. Drop this by fixing a common absolute lapse, and (T3)
transcendence returns as a god's-eye clock.

Thus the demarcation is non-vacuous: it isolates exactly the primitives 4.0 must keep
absent, mirroring 3.0's discipline that the boundary is constitutive, not a gap to be
filled later.
\end{remark}

\begin{corollary}[Admissibility against the frozen identification]\label{cor:adm}
Using \(C\) does not reintroduce the object 3.0 excluded. Since the excluded clock is
\(T\) and Theorem~\ref{thm:demarc} gives \(C\neq T\) with no external standpoint, the
4.0 register table --- personal kairos, created chronos, transcendent \textit{tota
simul} --- is earned rather than declared, and
\[
\textit{tota simul named, never occupied}
\]
holds as a theorem about the constructed \(C\), not as a posture.
\end{corollary}

\subsection{Status and reading}

\textbf{Target status:}~
\stDerived\ (Type-1, conditional on the minimal one-way P0.4 trace closure). The
demarcation is a theorem about the constructed object \(C\), resting on the locality
audit (Lemma~\ref{lem:local}), the created-now proposition (Proposition~\ref{prop:now}),
the frozen P0.2/P0.3 grounding, and the P0.4 one-way boundary closure. No new primitive
is introduced in P0.5. If the later canonical 4.0 body chooses the stronger two-way
intercession branch, the demarcation proof must be repeated with that additional
primitive in the inventory.

\medskip
\stInterp\quad
\emph{Reading.}
The communion really does share a present --- there is a genuine ``meanwhile'' in which
the living run alongside one another, and it can be measured without anyone stepping
outside time to keep it. But it is a creaturely present: only those still on the road are
in it, each only at the step they have reached, and it passes as they go. It is not the
single standing present in which a whole life is seen at once and a whole communion held
complete; that view belongs to no soul and to no clock the souls could build. 4.0 gives
the living a shared time without stealing the look from outside time --- which is exactly
why the word ``chronos'' may be used here without taking back what 3.0 refused.

\medskip
\textbf{Next node.} P0.6 --- exact uncoupled recovery: \(\varepsilon=0\Rightarrow
dX_i/dt=N_i(t)\,F_{3.0,i}(X_i)\), each trajectory gauge-equivalent to its isolated
\VFS\ 3.0 evolution, closing Stage~0 before any genuine plural coupling is switched on.

\section{P0.6 --- Exact uncoupled recovery (Theorem~I)}

Stage~0 closes with the obligation that legitimises the whole relational apparatus
from below: when relational action is switched off, nothing of 3.0 may have been
lost or quietly altered. The shared chronos, graph, rate cocycle, relational
duration, terminus strata, and traces must reduce to relational bookkeeping over
a family of interiors that each run \emph{exactly} their isolated \VFS\ 3.0
evolution.

This is not the same as requiring the graph to vanish. There are two limits:
\[
W=0
\quad
\text{removes the relational base,}
\]
while
\[
\varepsilon=0
\quad
\text{keeps possible relational read-outs but switches off their action.}
\]
P0.6 proves the second and stronger recovery statement. This is what makes 4.0 a
\emph{conservative extension} rather than a new theory wearing 3.0's variables.

\subsection{Coupling structure and the action parameter
\texorpdfstring{\(\varepsilon\)}{epsilon}}

Let \(F_{3.0,i}\) be the frozen isolated 3.0 vector field of interior \(i\), written
in its own kairos:
\[
\frac{dX_i}{d\sigma_i}
=
F_{3.0,i}(X_i).
\]
In a chronos chart this same isolated flow reads
\[
\frac{dX_i}{dt}
=
N_i(t)F_{3.0,i}(X_i),
\qquad
t\in\mathcal A_i.
\]

The plural layer is allowed to add only a gate-side relational action, and that
action is multiplied by a single Stage-0 action parameter \(\varepsilon\ge0\):
\[
\boxed{\;
\frac{dX_i}{dt}
=
N_i(t)F_{3.0,i}(X_i)
+
\varepsilon\,
\Phi_i
\Bigl(
X_i;\,
\{K_{ij}\},\,
\{K^{\rm post}_{i\leftarrow\ell}\},\,
\mathcal D_{\rm rel}
\Bigr),
\qquad
t\in\mathcal A_i.
\;}
\]
Here \(\mathcal D_{\rm rel}\) denotes any declared relational read-out data from
Stage~0, for example \(R_{ij}\), \(d\mathfrak c\), or graph labels, when such data
are needed to define the channel.

The recovery theorem assumes three admissibility conditions.

\begin{enumerate}[label=\textup{(\roman*)},leftmargin=1.2cm]
\item \emph{Gate-only action.}
The vector field \(\Phi_i\) acts solely on the receptivity-side slot
\[
\mathcal R_i=(\Gamma_i,H_{K,i},W_i)
\]
of \(X_i\). It touches no private scale and no dynamical variable directly.

\item \emph{\(\varepsilon\)-gating of every plural channel.}
Every ordinary edge kernel \(K_{ij}\), every post-terminus trace kernel
\(K^{\rm post}_{i\leftarrow\ell}\), and every future common-reception channel
acts on the interior only through the single prefactor \(\varepsilon\). No plural
channel bypasses this factor.

\item \emph{Read-outs are not actions.}
The objects
\[
R_{ij},
\qquad
d\mathfrak c,
\qquad
\mathcal G,
\qquad
\mathcal A_i/\mathcal B_i/\mathcal P_i
\]
may still be defined at \(\varepsilon=0\), but at \(\varepsilon=0\) they are
bookkeeping only. They do not modify \(F_{3.0,i}\).
\end{enumerate}

A finer per-edge family \(\{\varepsilon_{ij}\}\) is deferred to Stage~I. For the
Stage-0 recovery theorem, one scalar is sufficient because exact recovery means
simultaneous vanishing of all relational actions.

\subsection{The recovery theorem}

\begin{theorem}[Exact uncoupled recovery; Type-1]\label{thm:recovery}
Assume the coupling structure above and the minimal one-way P0.4 trace closure.
Setting \(\varepsilon=0\) gives, for every \(i\in\mathcal V\) and every
\(t\in\mathcal A_i\):
\begin{enumerate}[label=\textup{(\alph*)},leftmargin=1.2cm]
\item \emph{exact decoupling}:
\[
\frac{dX_i}{dt}
=
N_i(t)F_{3.0,i}(X_i),
\]
with no term depending dynamically on any \(X_j\), \(j\neq i\);

\item \emph{gauge-equivalence to isolated 3.0}: under
\[
\sigma_i(t)=\int_{t_0}^{t}N_i(s)\,ds,
\]
which is monotone on \(\mathcal A_i\) because \(N_i>0\),
\[
\frac{dX_i}{d\sigma_i}
=
F_{3.0,i}(X_i),
\]
so the \(\varepsilon=0\) trajectory is gauge-equivalent to \(i\)'s isolated
\VFS\ 3.0 evolution;

\item \emph{no single-interior content in chronos}: the Stage-0 chronos objects
\[
R_{ij},\quad
K_{ij},\quad
K^{\rm post},\quad
d\mathfrak c,\quad
\mathcal A_i/\mathcal B_i/\mathcal P_i
\]
either remain relational read-outs or are multiplied by \(\varepsilon\). Hence
none alters the \(\varepsilon=0\) single-interior flow.
\end{enumerate}
\end{theorem}

\begin{proof}
(a) Setting \(\varepsilon=0\) deletes the relational action term \(\Phi_i\)
identically. What remains,
\[
N_i(t)F_{3.0,i}(X_i),
\]
depends only on \(X_i\) and on the chart lapse \(N_i\). Any dependence on another
interior \(X_j\) entered only through the relational action term and is therefore
absent.

(b) On the active stratum \(\mathcal A_i\), \(d\sigma_i/dt=N_i>0\). Thus
\(t\mapsto\sigma_i(t)\) is a strictly increasing reparametrization. By the chain
rule,
\[
\frac{dX_i}{d\sigma_i}
=
\frac{1}{N_i(t)}\frac{dX_i}{dt}
=
F_{3.0,i}(X_i),
\]
which is exactly the frozen isolated 3.0 equation.

(c) By the locality audit, the Stage-0 inter-worldline inventory is made from
rate ratios, causal kernels, post-terminus trace kernels, relational duration,
graph data, and strata. The ratios and duration are read-outs unless explicitly
used inside an action channel. The ordinary kernels and post-terminus kernels act
only through \(\Phi_i\), and hence vanish from the equation at \(\varepsilon=0\).
The strata determine where the inherited flow is defined; they do not change the
field on the active arc. Therefore no Stage-0 chronos object survives as a force
or source in the single-interior dynamics. \qedhere
\end{proof}

\begin{corollary}[The frozen limits]\label{cor:limits}
\textup{(1)} \emph{Single-interior limit.}
If \(|\mathcal V|=1\), there are no inter-interior edges, no nontrivial rate ratios,
and no relational recipient for a post-terminus trace. On the active arc,
\[
\omega=N_1,
\qquad
d\mathfrak c=N_1dt=d\sigma_1,
\]
so 4.0 reduces to 3.0 identically, for every \(\varepsilon\).

\textup{(2)} \emph{Trivial-base limit.}
If \(W=0\), then there are no relational edges and no ordinary edge kernels:
\[
W=0
\quad\Longrightarrow\quad
L=\hat L=0.
\]
This removes the graph-level relational base.

\textup{(3)} \emph{Uncoupled plural limit.}
If \(\varepsilon=0\) and \(|\mathcal V|>1\), then the system is exactly
\(|\mathcal V|\) independent \VFS\ 3.0 interiors sharing at most a gauge chart and
relational read-outs:
\[
\varepsilon=0
\quad\Longrightarrow\quad
n\ \text{independent 3.0 interiors}.
\]
The implication holds even when relational bookkeeping is present. That is the
strong recovery theorem.
\end{corollary}

\subsection{Stage~0 closes}

\begin{remark}[Conservative extension]
Theorem~\ref{thm:recovery} and Corollary~\ref{cor:limits} discharge the recovery
requirement of the roadmap. The Stage-0 apparatus is a relational overlay that
has two clean reductions:
\[
W=0
\quad
\text{removes the relational base,}
\]
and
\[
\varepsilon=0
\quad
\text{switches off all relational action while allowing read-outs to remain.}
\]
The second reduction is the decisive one. It says that even after the base has
been declared, the plural theory can be made dynamically invisible to every
interior. Thus the frozen interface (P0.0), relational base (P0.1), rate cocycle
(P0.2), relational duration (P0.3), terminus closure (P0.4), demarcation theorem
(P0.5), and recovery theorem (P0.6) form a conservative Stage~0. Genuine plural
coupling \((\varepsilon>0)\) may now be introduced in Stage~I.
\end{remark}

\subsection{Status and reading}

\textbf{Target status:}~
\stDerived\ (Theorem~I). It depends on P0.0--P0.5: P0.0 fixes the frozen
interface, P0.1--P0.4 define the relational inventory, and P0.5 supplies the
locality audit and demarcation discipline. No new primitive is introduced in
P0.6. The single substantive modelling commitment is the
\(\varepsilon\)-gating of every plural action channel, which is exactly the
condition that makes recovery \emph{exact} rather than approximate.

\medskip
\stInterp\quad
\emph{Reading.}
Switch off the action of the bonds and each soul is left precisely its own 3.0
self, running its own kairos to its own terminus as if no other interior acted on
it. The shared chronos was never a leash and never a second law inside the soul.
It is a way of laying lives side by side so they can be ordered, compared, and
possibly brought into encounter. With \(\varepsilon=0\), the communion is not yet
an acting communion; it is many lives, each entirely itself, under a silent
relational chart. That is the sense in which 4.0 extends 3.0 without disturbing it.

\medskip
\textbf{Next node.} Stage~I --- \emph{comparability without fusion}: the first genuine
coupling (\(\varepsilon>0\)), where paces may be compared and influence may pass along
edges, while the interiors remain distinct --- no identification, no merged state, no
shared interior. Stage~0 has built the standing ground; Stage~I begins the encounter.

\section*{Stage~I --- Comparability without fusion}

\noindent Stage~0 built the relational base, proved its time creaturely (P0.5) and its
whole apparatus a conservative overlay (P0.6). Stage~I begins the encounter --- but
before any interaction is switched on, it must say \emph{what the plural object is}.
The answer is the thesis of the stage: a network of distinct \VFS\ 3.0 fibres above the
relational base, never one enlarged common interior. This stage establishes the arena
and the grammar of encounter and proves that the interiors stay \emph{two}; the living
mutual adjustment (\(\varepsilon>0\)) is the work of Stages~II--III.

\section{P1.0 --- Fibered encounter geometry}

\subsection{The fibered object}

\begin{definition}[Plural fibered interior]\label{def:fibered}
Let the Stage--0 relational base be
\[
\mathcal B_0
=
\Bigl(
\mathcal G,\ \prec,\
\{\mathcal A_i,\mathcal B_i,\mathcal P_i,t_i^\dagger\}_{i\in\mathcal V},\
\{R_{ij}\},\
\{K_{ij}\},\
\{K^{\rm post}_{j\leftarrow i}\},\
d\mathfrak c
\Bigr).
\]
For each vertex \(i\in\mathcal V\), the \emph{fibre} \(\mathcal M_i\) is interior
\(i\)'s frozen \VFS\ 3.0 geometry: its own interior \(3{+}1\) manifold, its own
kairos \(\sigma_i\), ascent road \(\lambda_{\mathrm{asc},i}\), available-Sophia
field \(\lambda_{S,i}\), state \(X_i\), terminus \(t_i^\dagger\), and frozen
gate/receptivity interface
\[
\mathcal R_i=(\Gamma_i,H_{K,i},W_i).
\]

The \emph{plural fibered interior} is the assembly
\[
\boxed{\;
\mathcal P
=
\bigl(\{\mathcal M_i\}_{i\in\mathcal V},\ \mathcal B_0,\ \pi\bigr),
\;}
\]
where \(\pi\) attaches each fibre \(\mathcal M_i\) to its base vertex \(i\).

No interior coordinate of \(\mathcal M_i\) is identified with any interior
coordinate of \(\mathcal M_j\) for \(i\neq j\). The only permitted relations between
distinct fibres are base-side relations from \(\mathcal B_0\), and their only
admissible future dynamical channel is through the recipient's gate/receptivity
interface \(\mathcal R_i\).
\end{definition}

\begin{remark}[Metric convention]
When the fibre metric is displayed schematically, it should not identify the
Stage--0 lapse \(N_i=d\sigma_i/dt\) with a new lapse inside the private proper-time
chart. A safe notation is
\[
ds_i^2
=
-d\sigma_i^2
+
\Omega_{P,i}^2 h^{(2)}_i
+
L_i^2\,d\lambda_{\mathrm{asc},i}^2,
\]
when \(\sigma_i\) is used as the private proper-time parameter. If another private
chart \(\chi_i\) is used, then a private lapse may appear there; the Stage--0 lapse
\(N_i\) remains the relational conversion \(d\sigma_i=N_i(t)\,dt\), not a new
interior dynamical variable.
\end{remark}

\(\mathcal P\) is thus neither a product geometry (the fibres share no interior
coordinates) nor a fusion (there is no common interior): it is a base-tied assembly of
otherwise self-standing 3.0 worlds. The base supplies \emph{comparability} --- the
invariant ratio \(R_{ij}\) records who runs faster --- while the fibres supply
\emph{distinctness}: each carries its own private kairos and its own terminus.

\subsection{Geometric non-fusion}

The structural content of ``without fusion'' is that \(\mathcal P\) is not a disguised
single interior. The invariant that proves this is the count of independent kairoi.

\begin{proposition}[Geometric non-fusion]\label{prop:nonfusion}
For \(|\mathcal V|=n\ge2\), the plural object \(\mathcal P\) is not isomorphic to
any single \VFS\ 3.0 interior. The invariant obstruction is the ownership count of
kairos histories,
\[
\{\sigma_i\}_{i\in\mathcal V}.
\]
A single 3.0 interior has exactly one owned kairos history and one terminus
boundary. The plural object has \(n\) owned kairos histories and \(n\) terminus
boundaries, even when several termini share the same chronos label.
\end{proposition}

\begin{proof}
The count is not the number of distinct coordinate values \(t_i^\dagger\): several
interiors may complete at the same chronos label. The invariant count is the number
of owned kairos histories and their terminal boundary events. Each fibre
\(\mathcal M_i\) has its own \(\sigma_i\) and its own boundary event
\(\mathcal B_i\), closed or active according to P0.4. A single 3.0 interior has only
one such history. Since \(n\ge2\), no relabelling of one interior can produce \(n\)
owned kairos histories. Hence no isomorphism with a single interior exists.
\qedhere
\end{proof}

This is non-fusion at the level of the plural geometry. The later dynamic claim is
conditional: it survives the switching-on of interaction exactly so long as the
interaction respects the Stage--0 admissibility contract. A coupling that literally
identifies fibres would not be a valid \VFS\ 4.0 coupling.

\begin{theorem}[Conditional robustness of non-fusion; Type-1]\label{thm:robust}
Assume that later \(\varepsilon>0\) coupling satisfies the Stage--0 admissibility
contract:
\begin{enumerate}[label=\textup{(\roman*)},leftmargin=1.2cm]
\item every plural action is \(\varepsilon\)-gated;
\item every plural action acts only through the recipient's gate/receptivity
interface \(\mathcal R_i=(\Gamma_i,H_{K,i},W_i)\);
\item no plural action identifies private coordinates, merges kairos histories, or
creates a common interior state vector.
\end{enumerate}
Then the number of owned kairos histories of \(\mathcal P\) remains
\(|\mathcal V|\) for every admissible \(\varepsilon\ge0\). Consequently
\(\mathcal P\) remains non-isomorphic to any single \VFS\ 3.0 interior under all
admissible later couplings.
\end{theorem}

\begin{proof}
Under the admissibility contract, coupling may modulate the recipient-side
receptivity slot but may not rewrite the fibre atlas, identify two private kairoi,
or replace the family \(\{\mathcal M_i\}\) by one common state space. Therefore the
owned kairos histories remain labelled by \(i\in\mathcal V\). Their count is a
gauge-invariant integer and cannot be changed by a gate-side modulation. The
obstruction in Proposition~\ref{prop:nonfusion} therefore survives admissible
coupling. \qedhere
\end{proof}

\noindent The geometric core of the stage is thus secured: interiors may be compared
(through \(R_{ij}\)) and, from Stage~II onward, coupled (through \(\mathcal R_i\)), while
remaining \(n\) distinct kairoi. Comparability acts on the base; distinctness lives in the
fibres; the two never collapse.

\subsection{Status and reading}

\textbf{Target status:}~
\stDef\ (the plural fibered interior \(\mathcal P\), Definition~\ref{def:fibered})
\(+\)
\stDerived\ (geometric non-fusion, Proposition~\ref{prop:nonfusion}, and its robustness
under coupling, Theorem~\ref{thm:robust}). It depends on Stage~0 only; no new primitive is
introduced, and the result is kinematic --- it concerns the structure of the plural
object, not yet any dynamical effect of interaction.

\medskip
\stInterp\quad
\emph{Reading.}
Two players, one stage, one groove that may form between them --- but two instruments that
never become one. The shared base is what lets them be set side by side, compared, and
later brought into encounter; the separate fibres are what keep them two. The point is not
that fusion is forbidden by fiat but that it would be self-defeating: one instrument
playing one line has no groove, because a groove needs the difference between two. So the
refusal to fuse is not a limit on communion --- it is the condition under which there is
anything to commune. Comparability without fusion is the geometry of a duet.

\medskip
\textbf{Next node.} P1.1 --- the invariant encounter hierarchy: replacing the gauge phrase
``the same \(t\)'' with invariant relational predicates,
\[
\text{co-presence}\ \subset\ \text{causal encounter}\ \subset\ \text{common reception},
\]
so that ``these two events met'' becomes a statement no relabelling can undo, and the
arena built here acquires its grammar of encounter.

\section{P1.1 --- The invariant encounter hierarchy}

P0.5 showed that the shared chart label \(t\) is pure gauge: a created now exists, but the
slice \(\{t=t_0\}\) is relabel-movable, so the sentence ``these two events happened at the
same \(t\)'' is not yet an invariant claim about the world. P1.1 supplies the invariant
replacements --- predicates of encounter that no relabelling can undo. They form a ladder
of strengthening requirements,
\[
\begin{aligned}
\text{co-presence}
&\ \xrightarrow{+\,\text{edge}+\text{support}}\
\text{causal encounter}\\
&\ \xrightarrow{+\,\text{common external field}}\
\text{common reception},
\end{aligned}
\]
where each rung adds a requirement to the one before. (The roadmap chain
\(\text{co-presence}\subset\text{causal encounter}\subset\text{common reception}\) is this
strength ladder; as \emph{sets of qualifying pairs} the inclusions run the other way, since
each added requirement admits fewer pairs.)

Throughout, \(E_i(s)\) is a marked invariant interior event (P0.1) at relational-duration
coordinate \(\mathfrak c(s)\), and \(\rightsquigarrow\), \(\prec\) are the direct causal
link and its transitive closure (P0.1).

\subsection{The three rungs}

\begin{definition}[Relational co-presence]\label{def:copres}
Two marked invariant events
\[
e=E_i(s),
\qquad
e'=E_j(s')
\]
are \emph{relationally co-present} when both are locatable in one common created
relational order, so that exactly one of
\[
\mathfrak c(s)<\mathfrak c(s'),
\qquad
\mathfrak c(s)=\mathfrak c(s'),
\qquad
\mathfrak c(s)>\mathfrak c(s')
\]
holds.

Thus co-presence at this rung means shared relational locatability and
comparability, not necessarily equality of duration-values. The special case
\[
\mathfrak c(s)=\mathfrak c(s')
\]
is \emph{simultaneous co-presence}. The general case includes ordered
co-presence: one event is before or after the other inside the same created
relational order.
\end{definition}

\noindent Two refinements will matter later. Co-presence is \emph{live} when both
interiors are active across an overlapping relational-duration window. It is \emph{trace}
co-presence when one interior is no longer active and is locatable only through its
terminus event or post-terminus trace. The trace case is co-present by relational
locatability, not by shared active meantime.

\begin{definition}[Causal encounter]\label{def:causenc}
Let
\[
e=E_a(s),
\qquad
e'=E_b(t)
\]
be relationally co-present marked events. They are in \emph{causal encounter} if,
in addition, at least one directed admissible channel carries influence from one
event to the other.

Concretely, \(e\) encounters \(e'\) in the direction \(a\to b\) when
\[
(a\to b)\in\mathcal E,
\qquad
\mathfrak c(s)<\mathfrak c(t),
\]
and the recipient-side kernel
\[
K_{ba}(t,s)
\]
has support from the source event \(E_a(s)\) to the recipient event \(E_b(t)\).
Equivalently, using the P0.1 convention, the first index of \(K_{ba}\) is the
recipient and the second index is the source.

For a completed source \(a\), the ordinary source history is replaced by the
boundary trace kernel
\[
K^{\rm post}_{b\leftarrow a}
\bigl(t;\,t_a^\dagger,\mathcal T_a^{\rm post}\bigr).
\]

The reverse direction \(b\to a\) is a different possible encounter and must be
checked separately.
\end{definition}

\begin{remark}[Direction matters]
Causal encounter is not symmetric by definition. Mutual encounter requires two
admissible directed channels, or a symmetric theorem class in which both
directions are guaranteed. P1.1 defines the predicate; it does not yet prove
reciprocity.
\end{remark}

\begin{definition}[Common reception schema]\label{def:comrec}
Events \(e=E_i(s)\) and \(e'=E_j(t)\) are in \emph{common reception} when there
exists a declared common external event
\[
\mathcal X\in\mathcal J_{\rm comm}
\]
which is not an owned interior event of \(i\) or \(j\), and whose reception is
registered by both interiors:
\[
\mathcal X\rightsquigarrow e,
\qquad
\mathcal X\rightsquigarrow e'.
\]

If the stricter roadmap ladder is used, common reception is counted as the top
rung only when \(e,e'\) are already in causal encounter. If the weaker sibling is
later admitted, common reception may instead be treated as a second floor-level
relation: two interiors jointly receive the same \(\mathcal X\) without a direct
edge between them.
\end{definition}

Under the stricter Stage--I ladder, the rungs nest as predicates:
\[
\boxed{\;
\text{common reception}\ \Longrightarrow\ \text{causal encounter}\ \Longrightarrow\ \text{co-presence}.
\;}
\]

The weaker sibling, if adopted in Stage~II, will be recorded separately and will not
be silently folded into this theorem chain.

\subsection{Invariance, and one honest exception}

\begin{remark}[Gauge-invariance of the hierarchy]
All three predicates are invariant under the chronos-chart gauge \(t\mapsto\widetilde t(t)\).
Each is built only from objects P0.5 already certified invariant: the relational duration
\(\mathfrak c\), the causal order \(\prec\) and its links \(\rightsquigarrow\), the edge set
\(\mathcal E\) and kernel supports, the active strata, and the external event \(\mathcal X\).
None refers to the chart label \(t\). Hence ``these two events met'' --- at any rung --- is
a statement no relabelling can undo, which is exactly what ``the same \(t\)'' failed to be.
\end{remark}

\begin{remark}[The trace exception]\label{rem:trace}
Causal encounter does not require \emph{live} co-presence. By the one-way closure of P0.4, a
completed interior \(j\) can act on a living \(i\) through the post-terminus trace kernel
\(K^{\rm post}_{i\leftarrow j}\): there is a genuine causal encounter, and co-presence holds
in the locatable sense, yet there is no shared active meanwhile --- \(j\) is no longer on
the stage. The dead reach the living across a gap that is not a common present. This is why
co-presence is defined at the level of locatability, not active overlap: it must remain the
true floor of the ladder even for trace-mediated encounters.
\end{remark}

\begin{remark}[A named sibling]
Common reception is defined here, per the roadmap, as a strengthening of causal encounter:
it presupposes an edge between \(i\) and \(j\). The weaker notion --- two interiors jointly
receiving the same external \(\mathcal X\) with \emph{no} direct edge between them (strangers
under one Word) --- is a natural sibling. It is named here and deferred; adopting it would
make common reception a second floor-level relation rather than a top rung, and is left to
the Stage~II receptivity treatment.
\end{remark}

\subsection{Status and reading}

\textbf{Target status:}~
\stDef\ (relational co-presence, causal encounter, and the common-reception
predicate schema) \(+\) \stDerived\ (gauge-invariance of the hierarchy, immediate
from P0.5). P1.1 introduces no new dynamical primitive. The external common event
field \(\mathcal J_{\rm comm}\) is only a named placeholder here; its admissible
Gate-side implementation belongs to Stage~II. Therefore common reception is a
schema at Stage~I and becomes an acting channel only after the bounded-input and
Gate-mediation rules are fixed.

\medskip
\stInterp\quad
\emph{Reading.}
Three ways for two to be ``together,'' each stronger than the last. To be \emph{co-present}
is to be on the one stage, locatable in the same unfolding --- even if one has already left
it and is heard only as an echo. To be in \emph{causal encounter} is for one to actually
reach the other --- the hearing, the touch, the word that lands. To be in \emph{common
reception} is for both to turn toward the same thing not their own --- the shared downbeat,
the cue neither gave, the gift both receive. The gauge slice could say none of this; it
could only say ``at the same mark,'' and the mark was movable. The invariant ladder says
what the slice could not: that an encounter, once it happens, has happened, and no
relabelling of time can take it back.

\medskip
\textbf{Next node.} P1.2 --- floor-labelled non-identifiability (the parallel Type-1
result): if two interiors carry different floors \(E_{0,i}\neq E_{0,j}\), their
floor-normalised ledgers differ (\(\beta_i=\alpha_i/E_{0,i}\),
\(\log s_{L,i}=\kappa_R G_i/E_{0,i}\)), so they cannot be identified as the same
floor-normalised interior --- structural non-identifiability, short of the full dynamic
non-fusion theorem.

\section{P1.2 --- Floor-labelled non-identifiability}

P1.0 secured non-fusion geometrically (the kairos count) and P1.1 the invariant grammar
of encounter. P1.2 adds one further, deliberately modest, distinctness result, drawn from
the single intrinsic scale each fibre inherits from 3.0 --- the floor. The discipline of
the node is to take exactly what the floor proves, and no more.

\subsection{The floor as a unit-independent scale}

The comparison \(E_{0,i}\neq E_{0,j}\) is made inside the inherited \VFS\ parameter
dictionary fixed by P0.0. It is not an arbitrary comparison of two unrelated unit
systems. Floor-normalisation may set the floor coordinate value to \(1\) inside a
single fibre, but it does not erase the invariant scale label \(E_{0,i}\) carried
by that fibre's frozen 3.0 data.

Each fibre carries, from its frozen 3.0 geometry, a finite floor \(E_{0,i}\) --- the
\emph{Imago Dei}, the one intrinsic scale of the interior core. Through it the ascent
coupling and the reset ledger are tied to the floor with no new primitive (3.0, v0.4):
\[
\beta_i=\frac{\alpha_i}{E_{0,i}},
\qquad
\log s_{L,i}=\frac{\kappa_R\,G_i}{E_{0,i}},
\]
where \(\beta_i\) is the de Sitter ascent coupling (\(A_i'/A_i=\beta_i\lambda_i\)),
\(s_{L,i}=A_+/A_-\) the reset scale-jump, \(G_i\) the reset source of fibre \(i\), and
\(\kappa_R\) the reset coupling. (Here \(G_i\) is the 3.0 reset source, not the relational
graph \(\mathcal G\).)

The floor is not a unit. Rescaling the ascent coordinate \(\lambda_i\to c\,\lambda_i\)
absorbs the floor \emph{value} \(L_{\rm floor}=1\), yet every invariant ---
\(L_i'/L_i\), \(L_i''/L_i\), the reset ratio \(s_{L,i}\), the exotic depth
\(-(\alpha_i\lambda_{\infty,i})^2/E_{0,i}^2\) --- is independent of that choice (3.0,
v0.9). What remains is the \emph{scale} \(E_{0,i}\) itself: the physics is the scale, the
``\(1\)'' is the unit. The floor is therefore a genuine unit-independent invariant of the
fibre.

\subsection{The non-identifiability theorem}

\begin{theorem}[Floor-labelled non-identifiability; Type-1]\label{thm:floor}
Inside the common inherited \VFS\ parameter dictionary, if
\[
E_{0,i}\neq E_{0,j},
\]
then fibres \(i\) and \(j\) cannot be identified as the same floor-normalised
interior.
\end{theorem}

\begin{proof}
Floor-normalisation fixes a coordinate convention; it does not delete the
unit-independent floor scale. Two interiors identified as the same
floor-normalised interior would share the same inherited scale label and all
unit-independent invariants. If \(E_{0,i}\neq E_{0,j}\), this already fails.
The obstruction is exhibited in the ledger: the ascent coupling
\(\beta_i=\alpha_i/E_{0,i}\) and the reset jump \(\log s_{L,i}=\kappa_R G_i/E_{0,i}\) carry
the floor scale through the \(1/E_{0,i}\) factor, and for matched ascent and reset data
(\(\alpha_i=\alpha_j\), \(\kappa_R G_i=\kappa_R G_j\)) they take distinct values when the
floors differ. Even in the non-generic case where \(\alpha\) and \(\kappa_R G\) conspire
so that \(\beta_i=\beta_j\) and \(\log s_{L,i}=\log s_{L,j}\) coincide, the floor scale
itself still differs, and identification fails on that invariant alone. Thus the floor
difference is a sufficient obstruction to identification. It is not a necessary
obstruction to distinction, because equal floors still leave all history, memory,
folding, reset, and relational-neighbourhood data available as individuating data.
\qedhere
\end{proof}

\subsection{What this does not say}

\begin{remark}[Non-identifiability, not non-fusion]\label{rem:notfusion}
Theorem~\ref{thm:floor} is a one-way, structural result and must not be over-read.
\begin{enumerate}[label=\textup{(\roman*)},leftmargin=1.2cm]
\item \emph{Sufficient, not necessary.} A floor difference is sufficient for
non-identification, not necessary. Equal floors \(E_{0,i}=E_{0,j}\) do \emph{not} imply
that \(i\) and \(j\) are the same interior, still less that they fuse.

\item \emph{Distinction survives equal floors.} Two interiors with the same floor may
still differ in initial conditions, in memory and folding history, in folding branch, in
reset history \(G_i\), and in relational neighbourhood (their position in
\(\mathcal G,\prec\)). Any of these keeps them distinct.

\item \emph{Structural, not dynamic.} This is non-identifiability of inherited
floor-labelled data, not a full dynamic non-fusion theorem. P1.0 gives the
kinematic obstruction by kairos count; P1.2 gives an additional sufficient
floor-label obstruction. The full claim that admissible coupled dynamics never
collapses two distinct histories into one belongs to the later locking and
Lyapunov stages.
\end{enumerate}
\end{remark}

This is the careful contribution of the floor: it labels distinctness when present, and
refuses to legislate identity when absent.

\subsection{Status and reading}

\textbf{Target status:}~
\stDerived\ (Type-1): the non-identifiability theorem, drawn from the unit-independence of
the floor scale (3.0~v0.9) with no new primitive. It depends on the inherited 3.0 floor
data and on P1.0; the accompanying \stInterp\ correction (Remark~\ref{rem:notfusion})
bounds the claim explicitly.

\medskip
\stInterp\quad
\emph{Reading.}
The floor \(E_{0,i}\) is the \emph{Imago Dei}, the intrinsic measure of a soul. Two souls
scaled differently in that measure cannot be normalised into one --- the image is a true
individuating scale, not a removable convention. But sameness of the image is not sameness
of the person: souls bearing the image at the very same scale are still told apart by all
they have lived --- their memories, their turnings, their resets, their place among
others. The image individuates; it does not exhaust. The person is always more than its
floor.

\medskip
\textbf{Stage~I closes.}
With P1.0, the plural object is fixed as a family of distinct frozen 3.0 fibres over
one created relational base. With P1.1, encounter is given an invariant grammar:
relational co-presence, directed causal encounter, and the common-reception schema.
With P1.2, the floor supplies a one-way structural distinctness label.

Thus Stage~I establishes \emph{comparability without fusion} at the structural
level:
\[
\boxed{
\text{interiors meet through the base, not by sharing a state layer.}
}
\]
The players are many, encounter is sayable invariantly, and the image-scale may
label them apart without exhausting their personhood. No acting communion has yet
been introduced except conditionally through the admissibility contract. The first
true \(\varepsilon>0\) Gate-side action begins in Stage~II.

\medskip
\textbf{Next node.} Stage~II --- Receptivity, Gate mediation, and bounded inputs, opening
at P2.1. Its governing principle: encounter changes the \emph{conditions of reception}; it
does not produce Sophia directly, transfer owned grace, or bypass the frozen Gate law.
This is where \(\varepsilon>0\) first acts --- where the musicians begin, at last, to
listen.

\section*{Stage~II --- Receptivity, Gate mediation, and bounded inputs}

\notice{\textbf{Governing principle.} Encounter changes the \emph{conditions of
reception}; it does not produce Sophia directly, transfer owned grace, or bypass the
frozen Gate law.}

\begin{remark}[Gate-rate convention for Stage~II]\label{rem:gate-rate-convention}
From Stage~II onward,
\[
I_{\mathrm{gate},i}
\]
denotes the gate-input rate measured in the recipient's own active kairos
\(\sigma_i\), not in an arbitrary chronos chart label \(t\). Thus
\[
\frac{d\mathcal G_{\mathrm{recepta},i}}{d\sigma_i}
=
I_{\mathrm{gate},i},
\]
and in a chronos chart,
\[
\frac{d\mathcal G_{\mathrm{recepta},i}}{dt}
=
N_i(t)\,I_{\mathrm{gate},i}.
\]
Equivalently, relative to the relational duration one-form
\(d\mathfrak c=\omega dt\),
\[
\frac{d\mathcal G_{\mathrm{recepta},i}}{d\mathfrak c}
=
\frac{N_i}{\omega}\,I_{\mathrm{gate},i}.
\]
This keeps the ledger invariant: \(I_{\mathrm{gate},i}d\sigma_i\) is the received
gate one-form of the recipient.
\end{remark}

\noindent Stage~I proved the players many and their meeting sayable. Stage~II switches on
the first genuine action: \(\varepsilon>0\). The discipline of the whole stage is the
principle above --- no influence may add grace, move grace, or open the gate by force; it
may only change how disposed a recipient is to receive through its own gate.

\section{P2.1 --- The receptivity-mediated Gate channel}

\subsection{The deleted channel}

A naive plural coupling would pour grace from neighbour to neighbour by an additive source
term directly on the available-Sophia field,
\[
\dot\lambda_{S,i}\ \supset\ C_{\lambda_S,i}\bigl(\{X_j\}_{j\to i}\bigr),
\]
the recipient's Sophia increased in direct proportion to the neighbours' states. This
channel is \emph{deleted}. It violates each clause of the governing principle at once: it
\emph{produces} Sophia directly (a source not passing the Gate), it \emph{transfers owned
grace} (treating grace as a stock moved between persons), and it \emph{bypasses the frozen
Gate law} (the recipient's own $I_{\mathrm{gate},i}$ no longer mediates what enters). No
admissible \VFS\ 4.0 coupling contains such a term.

\subsection{The receptivity-mediated channel}

What encounter may change is the recipient's \emph{receptivity}: the bounded transmissivity
\(\Gamma_i\in[0,1]\) of its own gate, and, where admissible, the Gate-memory variables
\(H_{K,i},W_i\) --- exactly the frozen receptivity interface \(\mathcal R_i\) of P0.0. The
selected schematic form of the inherited Open-Gate law, now carrying the modulation, is
\[
\boxed{\;
I_{\mathrm{gate},i}
=
\zeta_{0,i}\,e^{-\varphi_i\sigma_i}\,
\frac{
\Gamma_i\bigl(X_i,\ \mathcal N_i,\ J_{\mathrm{comm}}(\mathfrak c)\bigr)
}{
1+\kappa_{\lambda,i}\,\lambda_{S,i,+}
},
\qquad
0\le\Gamma_i\le1.
\;}
\]
Here \(\zeta_{0,i}>0\) is the node's gate amplitude or ceiling,
\(\varphi_i\ge0\) the kairos decay, \(\kappa_{\lambda,i}\ge0\) the
Sophia-saturation coefficient, and \(\lambda_{S,i,+}:=\max(\lambda_{S,i},0)\).
The symbol \(\kappa_{\lambda,i}\) is used deliberately: \(\chi_i\) is avoided here
because \(\chi\) is already tied to the inherited interior dynamical parameter in
the 3.0 layer. The read-out \(\mathcal N_i\) is the bounded relational
neighbourhood reaching \(i\) along declared in-edges, and
\(J_{\mathrm{comm}}(\mathfrak c)\ge0\) is the common external reception field of
P2.3, labelled by relational duration rather than by a raw chart time.

A relational event may change \(\Gamma_i\), or the admissible dynamics of
\(H_{K,i},W_i\). It never adds an unconstrained source to \(\lambda_{S,i}\), and
it never replaces the recipient's own gate law.

The relational dependence is \(\varepsilon\)-gated. At the generic P2.1 level one
may write
\[
\Gamma_i
=
\Pi_{[0,1]}
\Bigl(
\Gamma_i^{(0)}(X_i)
+
\varepsilon\,\delta\Gamma_i
\bigl(X_i,\mathcal N_i,J_{\mathrm{comm}}\bigr)
\Bigr),
\]
where \(\Pi_{[0,1]}\) is projection to \([0,1]\). This is a general bounded
interface statement: at \(\varepsilon=0\) it recovers the isolated 3.0 receptivity
\[
\Gamma_i=\Gamma_i^{(0)}(X_i),
\]
with no neighbour dependence. P2.3 will replace this generic projection by a
smoother headroom-opening formula for the positive common-reception field. The
point of P2.1 is only the admissibility contract:
\[
\text{relation may modulate bounded receptivity, not Sophia directly.}
\]

\subsection{The ceiling survives automatically}

\begin{proposition}[Automatic receptivity ceiling]\label{prop:ceiling}
For every admissible encounter --- any \(\Gamma_i\in[0,1]\), any \(\mathcal N_i\), any
\(J_{\mathrm{comm}}\ge0\) ---
\[
0\ \le\ I_{\mathrm{gate},i}\ \le\ \zeta_{0,i}.
\]
No configuration of encounters can drive the gate input beyond the node's own ceiling.
\end{proposition}

\begin{proof}
Each factor is nonnegative: \(\zeta_{0,i}>0\), \(e^{-\varphi_i\sigma_i}\in(0,1]\),
\(\Gamma_i\in[0,1]\), and \((1+\kappa_{\lambda,i}\lambda_{S,i,+})\ge1\) since
\(\kappa_{\lambda,i},\lambda_{S,i,+}\ge0\). Hence \(I_{\mathrm{gate},i}\ge0\). For the upper bound, each
of \(e^{-\varphi_i\sigma_i}\), \(\Gamma_i\), and \((1+\kappa_{\lambda,i}\lambda_{S,i,+})^{-1}\) lies in
\([0,1]\), so their product does, and \(I_{\mathrm{gate},i}\le\zeta_{0,i}\). The bound holds
pointwise in the encounter data, hence for every configuration. \qedhere
\end{proof}

\begin{remark}[The Stage--I contract is now instantiated]
The channel acts only on the recipient's receptivity slot
\(\mathcal R_i=(\Gamma_i,H_{K,i},W_i)\), is \(\varepsilon\)-gated, and identifies no private
coordinates and merges no kairos histories. It therefore satisfies the admissibility
contract assumed in Theorem~\ref{thm:robust}. Consequently P2.1 coupling preserves
non-fusion: the plural object remains \(|\mathcal V|\) distinct fibres while interacting.
The conditional theorem of Stage~I acquires its first concrete witness.
\end{remark}

\subsection{Status and reading}

\textbf{Target status:}~
\stDef\ (the receptivity-mediated channel, with the deleted direct-source term named and
excluded) \(+\) \stSelected\ (the schematic Open-Gate modulation form). It depends on the
frozen Open-Gate law (v2.0), P0.0, P0.6, and P1.0. The yield --- the ceiling
\(0\le I_{\mathrm{gate},i}\le\zeta_{0,i}\), Proposition~\ref{prop:ceiling} --- is
\stDerived\ and survives automatically from the bounded transmissivity. The common field
\(J_{\mathrm{comm}}\) enters here as a bounded argument only; its definition is P2.3, and
the non-fabrication theorem is P2.2.

\medskip
\stInterp\quad
\emph{Reading.}
Encounter changes how open you are, not what you are given. Another person can dispose you
--- soften you toward grace or harden you against it, encourage or scandalise, draw you out
or close you in --- but cannot be the source of your grace, cannot hand you grace as if it
were a stock to pass, and cannot force your gate. What enters still enters through your own
gate, at your own measure, and never above your own ceiling \(\zeta_{0,i}\): the community
can open you to the full of your measure, but not past it. This is how persons truly affect
one another's inner life without grace becoming a commodity and without anyone becoming the
source who is not the Source. The players may draw each other out; none can play another's
part, or lend another their tone.

\medskip
\textbf{Next node.} P2.2 --- no artificial grace creation: to derive non-fabrication from
the coupling architecture itself rather than impose a closed-system conservation law ---
peer encounter may alter the local Gate state but cannot create an unexplained Sophia term,
every received-grace contribution remaining explicitly accounted through the admissible
Gate input \(\dfrac{d}{dt}\mathcal G_{\mathrm{recepta},i}=I_{\mathrm{gate},i}\ge0\).

\section{P2.2 --- No artificial grace creation}

P2.1 deleted the direct grace-source and left receptivity modulation as the only plural
channel. P2.2 turns that architecture into a theorem: peer encounter cannot fabricate
grace. The discipline of the node is to obtain non-fabrication \emph{from the coupling
structure}, not by imposing a closed-system conservation law --- for the Open-Gate sector
is precisely not closed.

\subsection{The received-grace ledger}

\begin{definition}[Received-grace ledger]\label{def:ledger}
For each active interior \(i\), the received-grace ledger
\(\mathcal G_{\mathrm{recepta},i}\) is the accumulated admissible reception through
its own gate:
\[
\boxed{\;
d\mathcal G_{\mathrm{recepta},i}
=
I_{\mathrm{gate},i}\,d\sigma_i.
\;}
\]
Equivalently,
\[
\frac{d\mathcal G_{\mathrm{recepta},i}}{d\sigma_i}
=
I_{\mathrm{gate},i},
\qquad
\frac{d\mathcal G_{\mathrm{recepta},i}}{dt}
=
N_i(t)I_{\mathrm{gate},i},
\qquad
\frac{d\mathcal G_{\mathrm{recepta},i}}{d\mathfrak c}
=
\frac{N_i}{\omega}I_{\mathrm{gate},i}.
\]
The subscripted \(\mathcal G_{\mathrm{recepta},i}\) is the grace ledger, distinct
from the relational graph \(\mathcal G\).
\end{definition}

\subsection{The non-fabrication theorem}

\begin{theorem}[No artificial grace creation; Type-1]\label{thm:nofab}
Under the P2.1 architecture --- receptivity-mediated channel, direct source
\(C_{\lambda_S,i}\) deleted --- the following hold for every active interior \(i\):
\begin{enumerate}[label=\textup{(\alph*)},leftmargin=1.2cm]
\item \emph{Single accounted source.} The ledger changes only through the gate,
\[
d\mathcal G_{\mathrm{recepta},i}
=
I_{\mathrm{gate},i}\,d\sigma_i,
\qquad
I_{\mathrm{gate},i}\ge0.
\]
Thus the ledger is monotone non-decreasing along the recipient's own kairos.

\item \emph{No peer transfer.} No term moves grace from one interior to another.
Peer states enter only inside \(I_{\mathrm{gate},i}\), through the bounded
transmissivity \(\Gamma_i\); they scale \(i\)'s own admissible reception and never
withdraw from \(\mathcal G_{\mathrm{recepta},j}\).

\item \emph{Accounted and bounded.} Every positive contribution is an admissible
gate input, bounded nodewise in recipient-kairos rate:
\[
0
\le
\frac{d\mathcal G_{\mathrm{recepta},i}}{d\sigma_i}
=
I_{\mathrm{gate},i}
\le
\zeta_{0,i}.
\]
In a chronos chart the same statement reads
\[
0
\le
\frac{d\mathcal G_{\mathrm{recepta},i}}{dt}
=
N_iI_{\mathrm{gate},i}
\le
N_i\zeta_{0,i},
\]
so the chart expression carries the lapse factor and is not mistaken for an
absolute external rate.
\end{enumerate}
Hence grace is neither fabricated (no source outside the gate) nor transferred
(no peer stock movement).
\end{theorem}

\begin{proof}
By P2.1 the only plural channel is the \(\varepsilon\)-gated modulation of
\(\Gamma_i\in[0,1]\) inside \(I_{\mathrm{gate},i}\), and the direct additive source
\(C_{\lambda_S,i}\) is deleted. Therefore the only one-form feeding
\(\mathcal G_{\mathrm{recepta},i}\) is \(I_{\mathrm{gate},i}d\sigma_i\), which gives
(a); its sign is fixed by Proposition~\ref{prop:ceiling}. For (b), peer data appear
in the architecture only as arguments of the recipient's bounded receptivity
interface; no admissible term has the form ``subtract from
\(\mathcal G_{\mathrm{recepta},j}\), add to
\(\mathcal G_{\mathrm{recepta},i}\)''. For (c), the bound is exactly the automatic
ceiling of P2.1, expressed in the recipient's own kairos. \qedhere
\end{proof}

\subsection{Why not conservation}

The temptation is to secure non-fabrication by imposing a closed-transfer rule
\[
\sum_i C_{\lambda_S,i}=0,
\]
making peer grace zero-sum. This is rejected, and its rejection is the point of the node.

\begin{remark}[Non-fabrication is not conservation]\label{rem:notconservation}
The total received grace is \emph{not} conserved:
\[
\frac{d}{d\mathfrak c}
\sum_{i\in\mathcal V_{\rm act}^{\rm reg}}
\mathcal G_{\mathrm{recepta},i}
=
\sum_{i\in\mathcal V_{\rm act}^{\rm reg}}
\frac{N_i}{\omega}\,I_{\mathrm{gate},i}
\ge0.
\]
The total received grace, compared on the common relational-duration measure
\(d\mathfrak c\), is therefore not conserved and is generically strictly increasing
on active arcs. The factors \(N_i/\omega\) are the invariant conversion from each
recipient's kairos-rate ledger to the shared duration measure. The Open-Gate sector is
open: each gate receives from the
Source beyond the network, so the communion as a whole may grow ever richer at once. A
conservation law \(\sum_i C_{\lambda_S,i}=0\) would falsely recast this open communion as a
closed economy in which one interior's gain is another's loss. Non-fabrication is the
weaker and correct property: every reception is accounted through an admissible bounded
gate, and none is conjured or moved as stock --- but the sum is free to grow.
\end{remark}

\subsection{Status and reading}

\textbf{Target status:}~
\stDerived\ (Type-1): non-fabrication follows from the P2.1 architecture and the automatic
ceiling, with no conservation law imposed. It depends on P2.1. The accompanying correction
(Remark~\ref{rem:notconservation}) is the \stInterp\ guard against reading non-fabrication
as closure.

\medskip
\stInterp\quad
\emph{Reading.}
Grace is not a pie. The communion is not a closed economy where one soul's gain is
another's loss, so there is nothing here to envy and nothing to hoard: all may grow richer
together, because each is fed not from the others but from a Source beyond them all. Nor is
grace conjured within the circle --- no encounter makes grace from nothing, none pours its
own into another. What a neighbour can do is open you to receive more through your own gate;
what arrives is then truly received, accounted, and yours, taken from no one. The refusal to
impose a conservation law is the refusal of scarcity: a communion supplied from beyond has
no zero-sum to fight over.

\medskip
\textbf{Next node.} P2.3 --- common reception as a shared external reception field:
to define the one genuinely new positive external field of reception,
\[
J_{\mathrm{comm}}(\mathfrak c)\ge0,
\]
shared without being owned by the network, entering each interior only through a
bounded susceptibility
\[
\upsilon_i(\mathfrak c)\in[0,1]
\]
inside \(\Gamma_i\), with no requirement that
\[
\sum_i\upsilon_i=1
\]
and with the nodewise Gate ceiling preserved.

\section{P2.3 --- Common reception as a shared external field}

P2.1 modulated receptivity; P2.2 forbade fabrication and transfer. P2.3 introduces
the one genuinely new positive \emph{external reception field} of the plural
theory --- and keeps it disciplined by making it external and shared rather than
owned and divided. It is not a source term manufactured by the network and not a
direct source on \(\lambda_{S,i}\). It becomes active only through each recipient's
bounded Gate-side susceptibility. This is the acting form of the common-reception
rung of P1.1, including its edge-free sibling.

\subsection{The shared external field}

\begin{definition}[Common external reception field]\label{def:jcomm}
The \emph{common external reception field} is a nonnegative shared field
\[
J_{\mathrm{comm}}(\mathfrak c)\ge0,
\]
whose marked events are the declared common events
\(\mathcal X\in\mathcal J_{\mathrm{comm}}\) of P1.1. It is \emph{external} ---
supplied from beyond the network, like the Source of the Open Gate --- and
\emph{shared} --- available to be received by any interior. It is \emph{not} a
network-owned reservoir: the network neither holds it as stock nor divides it
among members. Only bounded images of this field may enter \(\Gamma_i\).
\end{definition}

Two interiors that both register \(\mathcal X\in\mathcal J_{\mathrm{comm}}\) are in common
reception (P1.1) \emph{whether or not} a direct edge joins them. Thus \(J_{\mathrm{comm}}\)
realises both the strict top rung (when the pair is already in causal encounter) and the
edge-free sibling (strangers under one field), which P1.1 named and deferred to here.

\subsection{Individual susceptibility and bounded composition}

\begin{definition}[Common-field susceptibility]\label{def:suscept}
Each interior carries an individual common-field susceptibility
\[
\upsilon_i(\mathfrak c)\in[0,1],
\]
its own openness to the common external field. The field enters the interior only
through the recipient's transmissivity \(\Gamma_i\), never as a direct source on
\(\lambda_{S,i}\). The notation \(\upsilon_i\) is used to avoid confusion with the
inherited wound/memory variable \(W_i\).
\end{definition}

Receptivity is assembled by a bounded composition that opens only the recipient's remaining
headroom:
\[
\boxed{\;
\Gamma_i
=
\Gamma_i^{(0)}(X_i)
+
\bigl(1-\Gamma_i^{(0)}(X_i)\bigr)\,\varepsilon\,\Theta_i,
\qquad
\Theta_i
=
\sum_{j\to i} a_{ij}\,\gamma_{ij}(\mathcal N_i)
+
\upsilon_i\,\gamma_i^{\mathrm{comm}}\bigl(J_{\mathrm{comm}}\bigr).
\;}
\]
with the nodewise normalisation (the safeguard)
\[
\sum_{j\to i} a_{ij}+\upsilon_i\le1,
\qquad
\varepsilon\in[0,1],
\]
with edge gains \(\gamma_{ij}\in[0,1]\) and common-field gain
\(\gamma_i^{\mathrm{comm}}\in[0,1]\). The field \(J_{\mathrm{comm}}\) itself may be
large or impulsive; only the bounded gain
\(\gamma_i^{\mathrm{comm}}(J_{\mathrm{comm}})\) enters the receptivity slot.

\subsection{The safeguard and non-rivalry}

\begin{proposition}[Bounded composition; nodewise ceiling preserved]\label{prop:safeguard}
Under the normalisation
\[
\sum_{j\to i}a_{ij}+\upsilon_i\le1
\]
with all gains in \([0,1]\) and \(\varepsilon\in[0,1]\), one has
\(\Theta_i\in[0,1]\) and hence
\[
\Gamma_i\in[\Gamma_i^{(0)},1]\subseteq[0,1]
\]
for arbitrarily many simultaneous encounters and any
\(J_{\mathrm{comm}}\ge0\). Consequently the automatic gate ceiling of P2.1
(Proposition~\ref{prop:ceiling}) survives unchanged:
\[
0\le I_{\mathrm{gate},i}\le\zeta_{0,i}.
\]
\end{proposition}

\begin{proof}
Each summand of \(\Theta_i\) is a product of a nonnegative weight and a gain in
\([0,1]\), so
\[
0\le\Theta_i\le\sum_{j\to i}a_{ij}+\upsilon_i\le1.
\]
Then
\[
\Gamma_i
=
\Gamma_i^{(0)}+(1-\Gamma_i^{(0)})\varepsilon\Theta_i
\]
is a headroom-opening interpolation. Since \(\varepsilon\Theta_i\in[0,1]\), it lies
between \(\Gamma_i^{(0)}\) and \(1\). Thus \(\Gamma_i\in[0,1]\) by construction,
with no clipping and no dependence on the number of contributing encounters.
Proposition~\ref{prop:ceiling} then applies verbatim. \qedhere
\end{proof}

\begin{remark}[The common field is non-rivalrous]\label{rem:nonrival}
The susceptibilities are \emph{not} a partition: there is no requirement that
\[
\sum_{i\in\mathcal V}\upsilon_i=1.
\]
The normalisation
\[
\sum_{j\to i}a_{ij}+\upsilon_i\le1
\]
is internal to each recipient's own composition, not a division of a fixed common
total across recipients. The same \(J_{\mathrm{comm}}\) may be received in full by
every interior at once, each according to its own \(\upsilon_i\), with no member's
reception diminishing another's. Common reception is therefore non-rivalrous, in
keeping with the open, non-conserved accounting of P2.2.
\end{remark}

\subsection{Status and reading}

\textbf{Target status:}~
\stDef\ (the common external field \(J_{\mathrm{comm}}\), the susceptibility \(\upsilon_i\), and
the bounded composition for \(\Gamma_i\)), with a \stDerived\ safeguard
(Proposition~\ref{prop:safeguard}: the nodewise ceiling survives arbitrarily many
encounters). It depends on P2.1 and P2.2 and on the P1.1 common-reception schema. The field
is the one new positive external reception field of the plural layer; it
remains Gate-mediated, ceiling-bounded, and non-rivalrous. It is not a
network-owned reservoir and not a direct source on \(\lambda_{S,i}\).

\medskip
\stInterp\quad
\emph{Reading.}
There is one thing a communion may receive together: a call from beyond it, the same for
all --- a word, a summons, a light not made within the circle. But it is not a stock the
community owns and rations. It is external and shared, and each receives it through their own
gate, at their own openness \(\upsilon_i\), to the full of their own measure. Many may hear the
same word, each wholly according to their attentiveness, and no one's hearing takes from
another's --- which is exactly why \(\sum_i \upsilon_i\) need not be one: the common gift is not
divided, it is shared. Even the common call does not force any gate; it opens only the room
each has left to open. Strangers with no tie between them may yet be joined by turning to the
same thing not their own.

\medskip
\textbf{Next node.} P2.4 --- individual input-to-state stability under relational forcing:
to prove, on an active domain, the bridge estimate
\[
\dot{\mathcal L}_{3.0,i}\le C_i-C_{1,i}\,\mathcal L_{3.0,i}+b_i\,\lVert u^{\mathrm{rel}}_i\rVert^2
\]
for the bounded relational input induced by the Gate modulation --- the layer without which
a plural Lyapunov sum would carry uncontrolled cross terms. Status \stConj\(\to\)\stDerived.

\section{P2.4 --- Individual input-to-state stability under relational forcing}

This is the analytic bridge of Stage~II. The earlier nodes made encounter bounded
(P2.1), non-fabricating (P2.2), and non-rivalrous (P2.3). P2.4 asks the stability
question: does a single interior stay well-behaved when relationally forced? Without an
affirmative input-to-state estimate, a future plural Lyapunov sum
\(\sum_i\mathcal L_{3.0,i}\) would be only a formal sum with uncontrolled cross terms.

\subsection{The inherited certificate and the relational input}

Each fibre inherits, frozen from \VFS\ v2.0 / 3.0, its dissipative certificate on
the active domain. We state it in the recipient's active kairos derivative
\[
D_i:=\frac{d}{d\sigma_i}.
\]
The isolated certificate is
\[
D_i\mathcal L_{3.0,i}
\le
C_i-C_{1,i}\mathcal L_{3.0,i},
\qquad
C_{1,i}>0.
\tag{H1}
\]
The bounded Open-Gate contribution enters the budget constant \(C_i\), not the
contraction coefficient \(C_{1,i}\).

The plural coupling reaches \(\mathcal L_{3.0,i}\) only through the recipient's
gate-side input direction. Splitting the gate input into its isolated value and its
relational deviation,
\[
I_{\mathrm{gate},i}
=
I_{\mathrm{gate},i}^{(0)}
+
u_i^{\mathrm{rel}},
\qquad
u_i^{\mathrm{rel}}:=\varepsilon\,\delta I_i,
\]
the isolated part is already absorbed in \(C_i\) by (H1). By the ceilings of
P2.1--P2.3 (Propositions~\ref{prop:ceiling},~\ref{prop:safeguard}), the relational
input is bounded:
\[
\lVert u_i^{\mathrm{rel}}\rVert
\le
\varepsilon\,\zeta_{0,i}.
\tag{H3}
\]

We add two structural bounds on the active domain:
\[
\lVert\nabla\mathcal L_{3.0,i}\rVert^2
\le
\mu_i\,\mathcal L_{3.0,i},
\qquad
\mu_i>0,
\tag{H2}
\]
and
\[
\lVert g_i(X_i)\rVert
\le
M_i,
\qquad
M_i<\infty,
\tag{H4}
\]
where \(g_i\) is the frozen gate-side input direction by which
\(u_i^{\mathrm{rel}}\) enters the \(i\)-th dynamics.

In a chronos chart the whole estimate is multiplied by the positive lapse \(N_i\)
on \(\mathcal A_i\). Thus the kairos statement is the invariant one; the chart
statement is only its representation.

\subsection{The input-to-state estimate}

\begin{theorem}[Individual ISS under relational forcing]\label{thm:iss}
Assume \textup{(H1)}, \textup{(H2)}, \textup{(H3)}, and \textup{(H4)}. Then on the
active domain there exist \(C_{1,i}'>0\) and \(b_i>0\) such that
\[
\boxed{\;
D_i\mathcal L_{3.0,i}
\le
C_i-C_{1,i}'\,\mathcal L_{3.0,i}
+
b_i\,\lVert u_i^{\mathrm{rel}}\rVert^2.
\;}
\]
The contraction \(C_{1,i}\) is reduced only by a controllable amount; it is not
destroyed.
\end{theorem}

\begin{proof}
Along the forced flow in recipient kairos,
\[
D_i\mathcal L_{3.0,i}
=
D_i\mathcal L_{3.0,i}\big|_{\mathrm{isolated}}
+
\bigl\langle
\nabla\mathcal L_{3.0,i},\,
g_i(X_i)u_i^{\mathrm{rel}}
\bigr\rangle.
\]
The isolated term obeys (H1). For the relational term, Young's inequality with a
free \(\delta>0\) gives
\[
\bigl\langle
\nabla\mathcal L_{3.0,i},\,
g_i u_i^{\mathrm{rel}}
\bigr\rangle
\le
\frac{\delta}{2}\,\lVert g_i\rVert^2
\lVert\nabla\mathcal L_{3.0,i}\rVert^2
+
\frac{1}{2\delta}\,\lVert u_i^{\mathrm{rel}}\rVert^2.
\]
Using (H2) and (H4),
\[
\bigl\langle
\nabla\mathcal L_{3.0,i},\,
g_i u_i^{\mathrm{rel}}
\bigr\rangle
\le
\frac{\delta\,\mu_iM_i^2}{2}\,\mathcal L_{3.0,i}
+
\frac{1}{2\delta}\,\lVert u_i^{\mathrm{rel}}\rVert^2.
\]
Choose \(\delta\) small enough that
\[
C_{1,i}'
:=
C_{1,i}-\frac{\delta\,\mu_iM_i^2}{2}
>0,
\]
and set
\[
b_i:=\frac{1}{2\delta}.
\]
This gives the stated estimate. \qedhere
\end{proof}

\subsection{Consequences}

\begin{corollary}[Bounded state under bounded encounter]\label{cor:bddstate}
Under the hypotheses of Theorem~\ref{thm:iss}, and as long as the trajectory stays
inside the active domain on which the constants are valid,
\[
\limsup_{\sigma_i}
\mathcal L_{3.0,i}(\sigma_i)
\le
\frac{
C_i+b_i(\varepsilon\zeta_{0,i})^2
}{
C_{1,i}'
}.
\]
Hence each interior is input-to-state stable under admissible relational forcing:
bounded encounter produces a bounded interior state on the active domain. No
configuration of admissible encounters drives a single interior to blow up inside
that domain.
\end{corollary}

Two structural points close the node. First, relational forcing enters only as a
bounded input term in the individual estimate; it does not remove the
sign-definite cleansing/alignment contraction inherited from 3.0. Second, with
each individual certificate surviving forcing up to a bounded input, the cross
terms of a future plural sum
\[
\sum_i\mathcal L_{3.0,i}
\]
are controlled rather than free. This is the prerequisite, deferred to Stage~V, for
a legitimate plural Lyapunov theory.

\subsection{Status and reading}

\textbf{Target status:}~
\stConj\(\to\)\stDerived. The estimate is \stDerived\ on the active domain under
the inherited certificate (H1), the quadratic-type gradient bound (H2), the bounded
relational input (H3), and the bounded gate-side input direction (H4). The
derivative is the recipient-kairos derivative \(D_i=d/d\sigma_i\); chronos-chart
versions must carry the lapse factor \(N_i\). The residual conjectural element is
uniformity: that (H2)--(H4) and the constants \(C_{1,i}',b_i\) hold uniformly over
the whole active domain, including corridor walls and the approach to terminus.
It depends on P2.1--P2.3 and the frozen 3.0 certificate.

\medskip
\stInterp\quad
\emph{Reading.}
A soul can be pushed and pulled by those it meets, and still not be thrown. Relational
forcing enters only as a bounded disturbance to the constant terms of the soul's own
balance; it never reaches the rate at which the soul cleanses and aligns itself. So
encounter may unsettle, may add weather, may raise the steady level a little --- but within
bounds the soul remains itself and does not fly apart. This is the first safety result of
the plural theory, and it is the individual one: before a communion can be shown stable as a
whole, each person must be shown unbreakable under being met. The player can be drawn off the
beat and yet keep their feet; only because each can, might the whole band later be shown to
hold together.

\medskip
\textbf{Stage~II closes.}
Encounter is bounded (P2.1), creates no artificial grace and transfers no owned
grace-stock (P2.2), receives a shared external field without rivalry (P2.3), and
does not destabilise the individual on the active domain (P2.4). Receptivity, Gate
mediation, and bounded inputs are in place. The first \(\varepsilon>0\) action is
safe at the level of the single interior:
\[
\boxed{
\text{encounter may alter openness, but never becomes ownership of grace.}
}
\]

\medskip
\textbf{Next node.} Stage~III --- Relational locking: rates, lags, and events, opening at
P3.0 (the asymptotic rate atlas, \(h_i=\beta_i\lambda_{S,i}\)). A standing caution carries
in: the 3.0 ascent is monotone and gapless, so nothing here may be called ``phase
synchronization'' until a genuine circle-valued phase observable is constructed.

\section*{Stage~III --- Relational locking: rates, lags, and events}

\notice{\textbf{Terminology correction.} Nothing in this stage may be called ``phase
synchronization'' until a genuine circle-valued phase observable exists. The 3.0 ascent
is monotone and gapless --- a climb, not an oscillator. What can be compared here are
\emph{rates}, not phases.}

\noindent Stage~II made encounter safe at the level of one interior. Stage~III asks the
comparative question: when can two interiors' ascents be set side by side? P3.0 builds the
atlas of rate observables that admit invariant cross-interior comparison, and marks plainly
what such comparison does and does not license.

\section{P3.0 --- The asymptotic rate atlas}

\subsection{The intrinsic ascent rate}

\begin{definition}[Intrinsic ascent-scale rate]\label{def:rate}
For each active interior \(i\), the \emph{intrinsic ascent-scale rate} is the
logarithmic rate of the ascent scale \(A_i\) in the interior's own kairos:
\[
\boxed{\;
h_i
:=
\frac{d}{d\sigma_i}\log A_i.
\;}
\]
The frozen 3.0 ascent law gives, on each admissible branch,
\[
\boxed{\;
h_i=\beta_i\,\lambda_{S,i},
\qquad
h_{*,i}=\beta_i\,\lambda_{S,\infty,i},
\;}
\]
with
\[
\beta_i=\frac{\alpha_i}{E_{0,i}}.
\]
Here \(\lambda_{S,i}\) is the available-Sophia field and
\(\lambda_{S,\infty,i}\) its asymptotic active-arc value. The coordinate
\(\lambda_{\mathrm{asc},i}\) is the metric road coordinate of epektasis; it is not
being differentiated here. The rate \(h_i\) measures how the ascent scale opens
along the road, not motion through the road coordinate itself.

On a folded interior the rate is read branch-wise: folding may alter the
admissible branch and the effective asymptotic value, but no cross-branch
identification is made.
\end{definition}

Because \(\beta_i=\alpha_i/E_{0,i}\), the ascent-scale rate is floor-labelled:
\(h_i\) is scaled by the inverse image-floor \(1/E_{0,i}\). Each interior's ascent
rate is proper to its own image-scale, not an entry on an absolute universal tempo.
This preserves the P1.2 discipline:
\[
\lambda_{S,i}\neq\lambda_{\mathrm{asc},i}.
\]

\subsection{The comparable duration-rate}

The intrinsic rate \(h_i\) lives in interior \(i\)'s own kairos and is not yet directly
comparable across interiors. The comparison is carried out on the shared relational
duration \(d\mathfrak c=\omega\,dt\) (P0.3):
\[
\boxed{\;
\hat h_i
:=
\frac{d}{d\mathfrak c}\log A_i
=
\frac{N_i}{\omega}\,h_i.
\;}
\]

This is defined on regular active arcs, where \(N_i>0\) and
\(d\mathfrak c=\omega dt\) is the P0.3 active duration one-form. No terminal
quantity \(N_{*,i}\) is introduced: the terminal rate \(h_{*,i}\) is an
intrinsic active-arc asymptotic of the 3.0 fibre, while the comparable rate
\(\hat h_i\) remains a duration-rate on whatever regular active arc exists before
the P0.4 hand-off.

\begin{proposition}[Invariant rate comparison]\label{prop:ratecompare}
The duration-rate \(\hat h_i\) is invariant under chronos reparametrisation, and
the cross-interior duration-rate ratio is
\[
\frac{\hat h_i}{\hat h_j}
=
\frac{N_i}{N_j}\,\frac{h_i}{h_j}
=
R_{ij}\,\frac{h_i}{h_j},
\]
with \(R_{ij}=N_i/N_j\) the P0.2 kairos-rate cocycle. Hence ``which interior's
ascent-scale opens faster per shared duration'' is a well-posed gauge-invariant
question.
\end{proposition}

\begin{proof}
Under a chronos chart change \(t\mapsto\widetilde t(t)\), \(N_i\) and \(\omega\)
transform by the same common density factor, so \(N_i/\omega\) is unchanged.
The intrinsic rate \(h_i=d(\log A_i)/d\sigma_i\) is defined in the private kairos
of the fibre and is chart-independent. Therefore
\[
\hat h_i=(N_i/\omega)h_i
\]
is invariant. The ratio formula follows by cancelling \(\omega\) and using the
cocycle identity of P0.2. \qedhere
\end{proof}

\subsection{What the atlas does not provide}

\begin{remark}[Rate, not phase]\label{rem:noPhase}
On the active ascending domain the ascent-scale rate is positive in the
admissible regime and the accumulated ascent-scale ledger
\[
\Lambda_i:=\log(A_i/A_{i,0})
\]
is monotone. There is no zero-crossing, return map, or period supplied by the
ascent sector itself. Consequently the
atlas supports \emph{magnitude} comparison of rates only. It defines no angle, no period,
and no circle-valued observable; there is no ``same point in a cycle'' because there is no
cycle. Any locking discussed in Stage~III is therefore rate-locking or lag-locking of
accumulated ledgers --- equality, ratio, or bounded offset of \(\hat h_i\) and
\(\Lambda_i\) --- and not phase-locking. A genuine phase observable, were one
ever introduced, would require additional structure absent from the monotone gapless
ascent, and would have to be constructed explicitly before any synchronization language
could be used.
\end{remark}

\subsection{Status and reading}

\textbf{Target status:}~
\stDef\ (the intrinsic ascent-scale rate \(h_i\), terminal active-arc rate
\(h_{*,i}\), and comparable duration-rate \(\hat h_i\)) \(+\) \stDerived\
(Proposition~\ref{prop:ratecompare}: invariance of \(\hat h_i\) and the cocycle
form of the duration-rate ratio, immediate from P0.2--P0.3). The intrinsic rate
form
\[
h_i=\beta_i\lambda_{S,i}
\]
is inherited frozen from 3.0 as the logarithmic ascent-scale law. It is not a law
for \(d\lambda_{\mathrm{asc},i}/d\sigma_i\). The floor label \(1/E_{0,i}\) is
inherited from P1.2. It depends on P0.0, P0.2--P0.3, and the frozen 3.0
ascent law.

\medskip
\stInterp\quad
\emph{Reading.}
One may truly ask whether one soul ascends faster than another, and the question has an
invariant answer on the communion's shared duration. But it is a question of rate, of how
fast each climbs --- never of phase, of where each stands in a cycle, because the ascent has
no cycle. The climb is monotone and gapless: it never returns, never closes, never arrives.
Souls may come to climb in step, at one rate, but they cannot be ``in phase,'' for there is
no beat to be on and no orbit to share. This is the formal shape of an ascent that is
endless --- always further up, never back around --- and the framework will not quietly turn
it into a wheel.

\medskip
\textbf{Next node.} P3.1 --- rate, lag, and event observables: to define the
synchronization \emph{targets} without importing oscillator language --- the comparable
rate already in hand, a duration-valued lag measuring how far apart two ascents sit on the
shared measure, and marked relational events --- so that Stage~III can speak of interiors
drawing level without ever positing a phase the ascent does not have.

\section{P3.1 --- Rate, lag, and event observables}

P3.0 gave one comparable scalar: the rate. P3.1 declares the full set of
\emph{synchronization targets} --- the precise senses in which two interiors might be said
to keep together --- and does so without positing any phase the monotone ascent does not
have. These are definitions of targets, not claims that locking occurs; the conditions
under which they are reached belong to the later small-gain and locking nodes.

\subsection{The accumulated observable}

\begin{definition}[Declared accumulated observable]\label{def:theta}
For each interior \(i\), an \emph{accumulated observable} \(\Theta_i\) is an
explicitly declared, monotone non-decreasing ledger carried along the ascent and
registered on the shared relational duration:
\[
\frac{d\Theta_i}{d\mathfrak c}\ge0.
\]
The standard ascent declaration is the accumulated ascent-scale ledger
\[
\Lambda_i
:=
\log(A_i/A_{i,0}),
\qquad
\frac{d\Lambda_i}{d\mathfrak c}
=
\hat h_i.
\]
A second standard declaration is accumulated reception,
\[
\Theta_i=\mathcal G_{\mathrm{recepta},i},
\qquad
\frac{d\Theta_i}{d\mathfrak c}
=
\frac{N_i}{\omega}I_{\mathrm{gate},i}.
\]
Which observable is meant must be declared. Lag statements are read relative to
that declaration. The road coordinate \(\lambda_{\mathrm{asc},i}\) may be used as a
spatial label of the metric ascent road, but it is not the default accumulated
dynamical ascent observable of Stage~III.
\end{definition}

\subsection{The three locking targets}

\begin{definition}[Relational locking targets]\label{def:locking}
Let \(i,j\) be co-present interiors on a common regular active arc, with marked
relational events \(\{E_{i,n}\}_n\) and a declared pairing \(n\mapsto m(n)\). The
Stage~III targets are:
\begin{align*}
&\text{\emph{duration-rate locking:}}
&&\frac{\hat h_i}{\hat h_j}\ \longrightarrow\ r^*_{ij}>0,\\[2pt]
&\text{\emph{affine ledger-lag locking:}}
&&\Theta_i-r^*_{ij}\,\Theta_j\ \longrightarrow\ \delta^*_{ij},\\[2pt]
&\text{\emph{event locking:}}
&&\Delta\mathfrak c\bigl(E_{i,n},E_{j,m(n)}\bigr)
\ \text{bounded or convergent,}
\end{align*}
all limits taken along the shared duration \(\mathfrak c\to\infty\) on the common
regular active arc, or up to the relevant P0.4 hand-off if the arc terminates,
where \(\Delta\mathfrak c(E_{i,n},E_{j,m(n)}):=\bigl|\mathfrak c(E_{i,n})-\mathfrak c(E_{j,m(n)})\bigr|\).
\end{definition}

Three points fix the meaning. The ratio target uses the comparable duration-rates
of P3.0:
\[
\hat h_i/\hat h_j=R_{ij}\,h_i/h_j,
\]
so rate locking is a gauge-invariant statement on the shared duration. The
constant \(\delta^*_{ij}\) is an affine lag in the declared ledger \(\Theta\) ---
ascent-scale ledger \(\Lambda\), reception ledger, or another explicitly declared
monotone observable --- and never in radians. Event locking is synchrony of marked
events on the shared duration measure, not coincidence of a phase.

\begin{remark}[Lag is offset, not rank]\label{rem:lagnotrank}
A nonzero \(\delta^*_{ij}\) says one interior runs a constant amount ahead of the other in
the declared accumulated observable, neither closing nor widening. It is an offset in
progress, not a ranking of worth; the image-floor labels of P1.2 are a separate datum and
are not what \(\delta^*_{ij}\) measures.
\end{remark}

\subsection{The optional circle-valued reduction}

\begin{remark}[When, and only when, a phase may be introduced]\label{rem:s1}
A circle-valued variable \(\theta_i\in S^1\) may be introduced \emph{only} if the interior
supplies a recurrent event sequence with a regular duration-spacing, or a genuinely
periodic reduced observable; \(\theta_i\) then advances by \(2\pi\) across one recurrence
cycle. In that conditional case a Kuramoto-type reduced equation
\[
\dot\theta_i=\Omega_i+\sum_{j}K_{ij}\sin(\theta_j-\theta_i)
\]
may be used as a \emph{local reduction tool} for the recurrent layer. It is never the
foundation: the foundational observables remain the monotone rate, the affine lag, and the
event gaps. The ascent itself supplies no such recurrence, so no phase exists by default;
\(\theta_i\) is admitted only where genuine periodicity is exhibited.
\end{remark}

\subsection{Status and reading}

\textbf{Target status:}~
\stDef\ (the accumulated observable \(\Theta_i\) and the three relational locking
targets --- duration-rate, affine ledger-lag, and event). The targets are
definitions of what keeping-together would mean, gauge-invariant through P3.0;
whether any is reached is deferred to P3.2--P3.3. The optional circle-valued
reduction is \stOpen: admitted only under an exhibited recurrence, never as a
primitive. It depends on P3.0 and the Stage~0--I event structure.

\medskip
\stInterp\quad
\emph{Reading.}
There are three honest ways for souls to keep together, and none of them makes the ascent a
wheel. They may climb at one rate, growing in step. They may hold a constant distance in
what they have accumulated --- one steadily ahead of the other, the gap neither closing nor
opening, like two pilgrims at one pace with one always a fixed way along the road. And they
may mark the same events together, their milestones falling at a bounded remove on the
shared measure. Only where life genuinely repeats --- where something truly comes round
again --- may one finally speak of a phase, and even then it is a local way of describing the
recurring layer, never the truth of the climb. The framework lets the liturgical wheel be
spoken of without ever claiming the soul's ascent is a wheel.

\medskip
\textbf{Next node.} P3.2 --- incremental receptivity and a small-gain criterion: to give the
first sufficient condition under which these targets are actually reached, bounding the
relational receptivity gain so that rate and lag differences contract rather than grow ---
the analytic engine, built on the bounded composition of P2.3 and the individual ISS of
P2.4, by which keeping-together becomes provable rather than merely definable.

\section{P3.2 --- Incremental receptivity and a small-gain criterion}

P3.1 defined what locking would mean. It does not follow from bounded coupling alone:
boundedness keeps interiors from blowing apart, but does not pull a mismatch down. P3.2
supplies the missing ingredient --- an actual contraction mechanism for mismatch, driven by
the way receptivity responds, incrementally, to relational difference --- and states the
small-gain condition under which that contraction wins.

\subsection{The mismatch observable}

\begin{definition}[Declared mismatch]\label{def:mismatch}
For co-present interiors \(i,j\) with a declared accumulated observable
\(\Theta\) and target ratio \(r^*_{ij}\), the \emph{affine ledger-lag mismatch} is
\[
\Delta_{ij}
:=
\Theta_i-r^*_{ij}\Theta_j-\delta^*_{ij}.
\]
Affine ledger-lag locking is the statement \(\Delta_{ij}\to0\). The corresponding
duration-rate mismatch is
\[
\hat h_i-r^*_{ij}\hat h_j.
\]
Statements below are read relative to the declared \(\Delta_{ij}\) and the chosen
ledger \(\Theta\).
\end{definition}

\subsection{The incremental-receptivity contraction}

Differentiating along the shared duration and collecting terms by origin gives the P3.2
estimate
\[
\boxed{\;
\frac{d}{d\mathfrak c}\Delta_{ij}
=
-\,\kappa^{\mathrm{eff}}_{ij}\Delta_{ij}
+\Delta^{\mathrm{free}}_{ij}(\mathfrak c)
+\mathcal R^{\mathrm{res}}_{ij}(\mathfrak c,\Delta_{ij}).
\;}
\]
where \(\Delta^{\mathrm{free}}_{ij}\) is the uncoupled detuning (the mismatch drift that would
persist at \(\varepsilon=0\)) and \(\mathcal R^{\mathrm{res}}_{ij}\) is the bounded relational
residual.

The contraction coefficient is \emph{not} assumed; it is read off the Gate-mediated
response under one structural hypothesis.

\begin{definition}[Incremental-receptivity hypothesis (IR)]\label{def:IR}
The relational part of receptivity responds to relative lag in a mismatch-reducing
direction: near the target, the gate-mediated contribution to \(d\Delta_{ij}/d\mathfrak c\)
has negative slope in \(\Delta_{ij}\),
\[
\kappa^{\mathrm{eff}}_{ij}
:=
-\,\frac{\partial}{\partial\Delta_{ij}}
\Bigl(\tfrac{d}{d\mathfrak c}\Delta_{ij}\Bigr)_{\!\mathrm{gate}}
>0,
\]
built from the bounded-composition slope of P2.3 (the \(\varepsilon\)-gated gain
\(\partial\gamma_{ij}/\partial\mathcal N_i\), the susceptibility \(\upsilon_i\), and the
headroom \(1-\Gamma_i^{(0)}\)).
\end{definition}

\begin{remark}[The residual is bounded by four sources]\label{rem:residual}
Writing the mismatch-dependent part of the residual as
\[
\mathcal R^{\mathrm{res}}_{ij}
=
\ell_{ij}(\mathfrak c)\Delta_{ij}
+
\mathcal R^{0}_{ij}(\mathfrak c),
\]
the loop gain and offset are assumed bounded on the regular active arc:
\[
0\le \ell_{ij}(\mathfrak c)\le \bar\ell_{ij},
\qquad
|\mathcal R^0_{ij}(\mathfrak c)|\le \bar R^0_{ij},
\qquad
|\Delta^{\mathrm{free}}_{ij}(\mathfrak c)|\le \bar\Delta^{\mathrm{free}}_{ij}.
\]
The bounds come from the same four sources: Sophia saturation, Gate memory,
graph weights, and frozen interior heterogeneity. Each is finite by P2.1--P2.4:
saturation and memory are bounded interface data, weights obey the P2.3
normalisation, and heterogeneity is fixed frozen 3.0 data.
\end{remark}

\subsection{The small-gain criterion}

\begin{theorem}[Small-gain contraction of mismatch; Type-1]\label{thm:smallgain}
Assume \textup{(IR)} with \(\kappa^{\mathrm{eff}}_{ij}\ge\underline\kappa_{ij}>0\),
the bounded residual decomposition above, and the bounded individual states of
P2.4. If the small-gain condition
\[
\boxed{\;
\bar\ell_{ij}<\underline\kappa_{ij}
\;}
\]
holds, then with
\[
\kappa_{ij}:=\underline\kappa_{ij}-\bar\ell_{ij}>0,
\]
one has
\[
\limsup_{\mathfrak c}|\Delta_{ij}|
\le
\frac{\bar\Delta^{\mathrm{free}}_{ij}+\bar R^0_{ij}}{\kappa_{ij}}.
\]
In particular, if the free detuning and offset residual vanish on the arc, then
\(\Delta_{ij}\to0\): the P3.1 affine ledger-lag target is reached.
\end{theorem}

\begin{proof}
Substitution gives
\[
\frac{d}{d\mathfrak c}\Delta_{ij}
=
-\bigl(\kappa^{\mathrm{eff}}_{ij}-\ell_{ij}\bigr)\Delta_{ij}
+\Delta^{\mathrm{free}}_{ij}
+\mathcal R^0_{ij}.
\]
By the small-gain condition,
\[
\kappa^{\mathrm{eff}}_{ij}-\ell_{ij}\ge\kappa_{ij}>0.
\]
Using the upper Dini derivative of \(|\Delta_{ij}|\),
\[
D^+|\Delta_{ij}|
\le
-\kappa_{ij}|\Delta_{ij}|
+\bar\Delta^{\mathrm{free}}_{ij}
+\bar R^0_{ij}.
\]
The scalar comparison lemma yields the stated limsup bound. If the forcing bounds
are zero, the comparison reduces to
\[
D^+|\Delta_{ij}|\le-\kappa_{ij}|\Delta_{ij}|,
\]
so \(\Delta_{ij}\to0\). The bounded states of P2.4 keep the estimate inside the
active domain where the constants are valid. \qedhere
\end{proof}

\subsection{Status and reading}

\textbf{Target status:}~
\stConj\(\to\)\stDerived. The contraction estimate and the small-gain bound are
\stDerived\ under the incremental-receptivity hypothesis (IR), the four-source residual
bound, and the bounded states of P2.4. The residual conjectural element is the sign and
uniformity of \(\kappa^{\mathrm{eff}}_{ij}\) over the whole admissible regime: (IR) is a
structural monotonicity of the Gate response, natural but not unconditional, and carrying it
to fully unconditional \stDerived\ is the work of the explicit response model. It depends on
P2.1--P2.4 and P3.0--P3.1.

\medskip
\stInterp\quad
\emph{Reading.}
This is where keeping-together stops being merely nameable and becomes provable --- and the
mechanism is receptivity itself. When two souls drift apart in their ascent, a responsive
openness tends, of itself, to pull them back: the one that lags grows more receptive, and
the gap closes. But the pull is not a clamp. It prevails only when the responsiveness
outweighs how differently the two are made and how much memory and saturation resist ---
the small-gain condition, which is, plainly, \emph{listen more than you differ}. Where it
holds, drift contracts to a bounded remove or to nothing; where it fails --- too great a
difference, too hardened a heart, too thin a bond --- the souls keep a bounded distance or
drift free. Communion has a built-in restoring tendency, gentle and conditional, never
forced, and it does its work through openness, not through any transfer of grace.

\medskip
\textbf{Next node.} P3.3 --- the maximum lockable mismatch: to turn the small-gain bound into
a sharp threshold, the largest detuning \(\Delta^{\mathrm{free}}_{ij}\) and residual a given
bond can still draw into locking --- the explicit boundary, in heterogeneity and bond
strength, between a communion that holds together and one that comes apart.

\section{P3.3 --- The maximum lockable mismatch}

P3.2 gave a contraction once the effective rate \(\kappa^{\mathrm{eff}}_{ij}\) was positive,
but treated the restoring term as if it acted without limit. It does not: the Gate ceiling
caps how hard receptivity can pull. P3.3 carries that finiteness through and turns it into a
sharp conditional threshold --- the largest detuning a given bond can still draw into
locking.

\subsection{The detuning and the bounded restoring term}

Write the constant or bounded rate detuning as
\[
\Delta\varpi_{ij}:=\Delta^{\mathrm{free}}_{ij},
\]
a duration-rate mismatch on the shared measure, not an angular frequency and not
the density \(\omega\) of P0.3. For an oriented recipient-side channel
\(j\to i\), the mismatch equation is written
\[
\frac{d}{d\mathfrak c}\Delta_{ij}
=
\Delta\varpi_{ij}-F_{j\to i}(\Delta_{ij}),
\]
where \(F_{j\to i}\) is the Gate-mediated restoring term acting through recipient
\(i\)'s receptivity interface: continuous, with \(F_{j\to i}(0)=0\), restoring slope
\(F_{j\to i}'(0)=\kappa^{\mathrm{eff}}_{ij}>0\), and bounded because receptivity is
bounded.

The symmetric capacity bound is the explicit product of the four declared data,
\[
\boxed{\;
\Delta\varpi_{\max,j\to i}
:=
F_{\max,j\to i}
=
\frac{N_i}{\omega}\,
\beta_i\,\zeta_{0,i}\,
\underbrace{\varepsilon\,a_{ij}\bigl(1-\Gamma_i^{(0)}\bigr)}_{\text{relational gain}}\,
\underbrace{\frac{1}{1+\kappa_{\lambda,i}\lambda_{S,i,+}}}_{\text{saturation}}\,
\underbrace{b^{\mathrm{fold}}_{j\to i}}_{\text{branch}},
\;}
\]
on the regular active arc. The factor \(N_i/\omega\) appears because the Gate
ceiling is a recipient-kairos rate, while the mismatch is measured per relational
duration \(d\mathfrak c\). For a reciprocal pair, the usable pair threshold is the
oriented minimum
\[
\Delta\varpi_{\max,ij}^{\rm pair}
:=
\min\{\Delta\varpi_{\max,j\to i},\Delta\varpi_{\max,i\to j}\},
\]
unless the model explicitly chooses a one-sided leader/follower locking
convention.

\subsection{The threshold}

\begin{theorem}[Maximum lockable mismatch; Type-1 / Type-2]\label{thm:maxlock}
Assume a continuous oriented restoring law \(F_{j\to i}\) with
\[
F_{j\to i}(0)=0,
\qquad
F'_{j\to i}(0)>0,
\]
and assume that on the relevant branch it is monotone increasing until saturation,
with one-sided capacities
\[
F^+_{\max,j\to i}:=\sup_{\Delta}F_{j\to i}(\Delta),
\qquad
F^-_{\max,j\to i}:=\sup_{\Delta}\bigl(-F_{j\to i}(\Delta)\bigr).
\]
Then a stable locked state exists whenever
\[
-F^-_{\max,j\to i}
<
\Delta\varpi_{ij}
<
F^+_{\max,j\to i}.
\]
If
\[
\Delta\varpi_{ij}>F^+_{\max,j\to i}
\quad\text{or}\quad
\Delta\varpi_{ij}< -F^-_{\max,j\to i},
\]
no locked state exists on that oriented branch. In the symmetric-capacity case
\[
F^+_{\max,j\to i}=F^-_{\max,j\to i}
=\Delta\varpi_{\max,j\to i},
\]
this reduces to
\[
|\Delta\varpi_{ij}|<\Delta\varpi_{\max,j\to i}.
\]
The explicit product above is the Type-2 estimate for the symmetric capacity.
\end{theorem}

\begin{proof}
A locked state is a zero of
\[
\Delta\varpi_{ij}-F_{j\to i}(\Delta),
\]
i.e. a solution of
\[
F_{j\to i}(\Delta^*)=\Delta\varpi_{ij}.
\]
On a monotone branch, the range of \(F_{j\to i}\) is the interval
\[
(-F^-_{\max,j\to i},F^+_{\max,j\to i})
\]
up to endpoint conventions. Hence the intermediate value theorem gives a solution
precisely inside that range. On the increasing branch,
\[
\frac{d}{d\Delta}
\bigl(\Delta\varpi_{ij}-F_{j\to i}(\Delta)\bigr)\big|_{\Delta^*}
=
-F'_{j\to i}(\Delta^*)<0,
\]
so the locked state is asymptotically stable. Outside the one-sided capacity
range the vector field keeps a fixed sign on that branch, so no branch-locked
state exists. \qedhere
\end{proof}

\subsection{The ceiling is necessary, not sufficient}

\begin{remark}[Two conditions, not one]\label{rem:ceilingnec}
The Gate ceiling \(I_{\mathrm{gate},i}\le\zeta_{0,i}\) is what makes \(F_{\max,j\to i}\) finite,
and so is \emph{necessary} for a finite, well-posed locking range. But it is not
\emph{sufficient} by itself: boundedness caps the magnitude of the forcing without showing
that the forcing points in the mismatch-reducing direction. That direction is the separate
content of the incremental-receptivity hypothesis (P3.2), which supplies \(F_{ij}'(0)>0\).
Locking therefore needs both: a bounded capacity (ceiling) \emph{and} a restoring
orientation (IR). The ceiling sets the \emph{size} of the lockable range; IR establishes
that there is a range at all.
\end{remark}

\subsection{Status and reading}

\textbf{Target status:}~
\stDerived\ (Type-1: existence and sharpness of the threshold under the monotone
saturating restoring-law hypothesis; Type-2: the explicit oriented product form of
\(\Delta\varpi_{\max,j\to i}\)). It depends on P3.2 for
\(\kappa^{\mathrm{eff}}>0\) and the residual structure, on P2.1--P2.3 for the Gate
ceiling and bounded composition, and on the frozen 3.0 data
\((\beta_i,\zeta_{0,i})\) plus the branch factor. The product is recipient-oriented
because the channel acts through \(i\)'s Gate. A reciprocal locking claim must use
the oriented minimum or state a one-sided convention.

\medskip
\stInterp\quad
\emph{Reading.}
There is a largest difference that communion can bridge. The pull that draws two ascents
together is real, but finite: it is capped by how far the gate can open. Where the detuning
between two souls falls within that range, they lock --- settling to a stable, constant lag,
neither drifting apart nor collapsing into one. Where the difference exceeds it, they drift,
not for want of pull but because the pull, being finite, cannot span so wide a gap. This is
neither the sentiment that love bridges everything nor the despair that difference always
wins: there is a real, computable boundary, set by how open the gate can be, how saturated,
how strong the bond, and which branch one stands on. And the bound teaches a last
distinction: a bounded openness is necessary for there to be any finite power to bridge at
all, but capacity is not enough --- the openness must also turn toward the other. A finite
love, rightly directed.

\medskip
\textbf{Next node.} P3.4 --- dynamic non-fusion under admissible locking: to prove that
locking, even when achieved, never collapses the locked interiors into one --- that the
stable lag \(\Delta^*_{ij}\) stays strictly nonzero in identity even as it becomes constant
in value, so that Stage~III's keeping-together never becomes the fusion Stage~I forbade.

\section{P3.4 --- Dynamic non-fusion under admissible locking}

Stage~I proved non-fusion kinematically: the kairos-history count stays
\(|\mathcal V|\) under admissible coupling (Theorem~\ref{thm:robust}), and
distinct floors obstruct identification (Theorem~\ref{thm:floor}). Stage~III has
now made locking not merely possible but, under the threshold of P3.3, achieved.
P3.4 proves the dynamic form: even when selected observables lock, the fibres are
not identified. Equality of some numerical values is allowed; fusion is not.
A stable relation is not a merger.

\subsection{What locking constrains, and what it leaves free}

Locking constrains only the comparative observables of P3.1:
\[
\hat h_i/\hat h_j\to r^*_{ij},
\qquad
\Theta_i-r^*_{ij}\Theta_j\to\delta^*_{ij},
\]
and the marked-event gaps. It does not identify the owning fibres. The labelled
interior datum is
\[
\widetilde X_i
:=
\bigl(
i,\ X_i,\ E_{0,i},\ \sigma_i\text{-history},\
H_{K,i}\text{-history},\ W_i\text{-history},\
\text{fold-branch}_i
\bigr).
\]
The label \(i\) and the owned kairos history are not gauge decorations; they are
part of the plural fibre structure fixed in P1.0. Therefore even if some selected
coordinates momentarily coincide, the interiors are not thereby fused.

\subsection{The dynamic non-fusion theorem}

\begin{theorem}[Dynamic non-fusion under locking; Type-1]\label{thm:dynnofusion}
Under admissible locking, the labelled fibres remain distinct:
\[
\widetilde X_i(\mathfrak c)\neq \widetilde X_j(\mathfrak c)
\qquad
(i\neq j).
\]
For generic distinct data, even the unlabelled state-values \(X_i\) and \(X_j\)
remain different except possibly at isolated coincidences of selected observables.
The stronger equality \(X_i\equiv X_j\) for all locked time occurs only on the
exceptional value-diagonal of Definition~\ref{def:diag}; even there, the fibres
remain two owned kairos histories rather than one fused interior.
\end{theorem}

\begin{proof}
Admissible locking fixes only comparative observables: duration-rate ratios,
ledger lags, and event gaps. By the Stage--0/Stage--I admissibility contract it
does not identify fibre atlases, erase the owner label \(i\), merge kairos
histories, or rewrite the frozen floor. Therefore the labelled data
\(\widetilde X_i\) and \(\widetilde X_j\) remain distinct for \(i\neq j\).

If \(E_{0,i}\neq E_{0,j}\), Theorem~\ref{thm:floor} already prevents
floor-normalised identification. If floors coincide, distinct memory histories,
wound histories, fold branches, initial states, or relational neighbourhoods still
give generic separation of the unlabelled state-values. Complete value equality
throughout the locked regime requires every individuating datum and every
evolution law to coincide; this is precisely the exceptional value-diagonal.
But value-diagonal equality is not fusion in the 4.0 sense: it makes two fibres
exact copies, not one fibre. \qedhere
\end{proof}

\begin{definition}[Exceptional value-diagonal]\label{def:diag}
The \emph{exceptional value-diagonal} is the set of fully coincident unlabelled
data for a pair \(i,j\):
\[
\mathcal M_{\mathrm{diag}}^{\rm val}
=
\left\{
\begin{aligned}
&E_{0,i}=E_{0,j},\quad
H_{K,i}\equiv H_{K,j},\quad
W_i\equiv W_j,\\
&\text{same fold-branch},\quad
r^*_{ij}=1,\quad
\delta^*_{ij}=0,\quad
X_i(0)=X_j(0),\\
&\text{same admissible Gate/receptivity law and same relational inputs}
\end{aligned}
\right\}.
\]
It is invariant as a value-diagonal: exact copies remain value-equal under equal
inputs. It is non-generic in the space of distinct interior data. But it is not a
fusion manifold. Even on
\(\mathcal M_{\mathrm{diag}}^{\rm val}\), the plural object still contains two
owned fibres unless an inadmissible extra identification quotient is imposed.
\end{definition}

\subsection{Status and reading}

\textbf{Target status:}~
\stDerived\ for labelled-fibre non-fusion: it follows directly from P1.0's fibre
ownership and the admissibility contract. The generic unlabelled value-separation
statement is \stConj\(\to\)\stDerived, conditional on the complete state-space
argument and on exclusion of hidden value-identifying symmetries. Floor-driven
separation is \stDerived\ outright by Theorem~\ref{thm:floor}. It depends on P1.0,
P1.2, and P3.2--P3.3.

\medskip
\stInterp\quad
\emph{Reading.}
Communion is the stable answerability of distinct interior histories, not their collapse
into one state. When two souls lock, what becomes constant is the \emph{relation} between
them --- a fixed ratio of ascent, a steady lag, shared milestones --- and a constant relation
is the very opposite of a merger: only two distinct things can stand in a stable relation at
all. They may climb at one rate and keep one lag forever, and still each remains itself, with
its own floor, its own memory, its own wounds, its own branch of the fold. The single case in
which two would become one is the case in which they were never two: identical in every
individuating thing from the start. Everywhere else --- everywhere real --- locking is
nearness that does not consume, the answer of one history to another that leaves both
standing.

\medskip
\textbf{Stage~III closes.}
Rates are comparable as ascent-scale duration-rates (P3.0), the targets of
keeping-together are defined without positing a phase (P3.1), an
incremental-receptivity small-gain criterion makes locking provable (P3.2), an
oriented threshold bounds what difference can be bridged (P3.3), and locking, when
achieved, never fuses the locked (P3.4). Relational locking is in place, and it is
non-fusing:
\[
\boxed{
\text{souls may keep together without becoming the same.}
}
\]

The central correction of Stage~III is therefore:
\[
\boxed{
\text{locking is a stable relation between fibres, not an identification of fibres.}
}
\]

\medskip
\textbf{Next node.} Stage~IV --- Junctions, refractory structure, and plural Resurrectio,
opening at P4.0 (junction taxonomy): separating three events that must not be conflated ---
a private single-interior reset, a relational encounter, and a common reception event --- and
ruling which variables may jump, which may change only continuously, and which are read
merely as recipient-side receptivity changes.

\section*{Stage~IV --- Junctions, refractory structure, and plural Resurrectio}

\notice{\textbf{Governing discipline.} Three events must never be conflated: a private
single-interior reset, a relational encounter, and a common reception event. Only the first
may make an owned state jump, and only in its own interior.}

\noindent Stage~III showed how interiors keep together without fusing. Stage~IV turns to the
discontinuous structure: resets, refractory windows, and the plural form of Resurrectio.
The stage opens by fixing, once and cleanly, which events may do what to which variables ---
the junction taxonomy --- so that nothing later smuggles a jump where only a continuous
change is admissible.

\section{P4.0 --- Junction taxonomy}

\subsection{Three events that must not be conflated}

\begin{definition}[Three junction-event classes]\label{def:junctions}
Stage~IV distinguishes three admissible classes of distinguished event. They may
be absent at a given point, and more than one may be registered at the same
relational label only if their signatures remain disjoint. Their admissible write
sets are different:
\begin{enumerate}[label=\textup{(\arabic*)},leftmargin=1.2cm]
\item \emph{Private single-interior reset} \(\mathsf J^{\mathrm{reset}}_{i}\), occurring at
\(i\)'s own reset events \(E^{\mathrm{reset}}_{i,n}\). It is an own-interior
discontinuous self-map inherited frozen from 3.0. The ascent scale jumps by the
reset factor
\[
A_i\longmapsto s_{L,i}A_i,
\qquad
\log s_{L,i}
=
\frac{\kappa_R\,G_i^{\mathrm{reset}}}{E_{0,i}},
\]
where \(G_i^{\mathrm{reset}}\) is the local reset reception datum of interior \(i\)
at that junction, not a transferable stock. The admissible reset may also update
the inherited local reset variables such as \(\lambda_{S,i},H_{K,i},W_i\) according
to the frozen 3.0 reset rule. It involves no other interior.
\item \emph{Relational encounter}, the \(\varepsilon\)-gated peer action of P2.1.
It writes nothing into \(i\)'s owned state directly. It enters only as relational
neighbourhood data \(\mathcal N_i\), and hence only through the bounded
receptivity interface \(\Gamma_i\).
\item \emph{Common reception event}, the shared external reception field of P2.3.
It writes no owned state directly. It enters only through the susceptibility term
\[
\upsilon_i\,\gamma_i^{\mathrm{comm}}(J_{\mathrm{comm}})
\]
inside \(\Gamma_i\), and even an impulsive \(J_{\mathrm{comm}}\) enters only through
its bounded gain.
\end{enumerate}
\end{definition}

\begin{remark}[Coincident labels do not merge signatures]
Two event classes may be labelled by the same relational duration
\(\mathfrak c_*\), for example a private reset occurring while a common reception
field is also nonzero. This coincidence of labels does not merge their write rules.
The reset writes only according to the private 3.0 reset map; the common event is
still only a bounded receptivity read-out. Coincidence in chronos is not identity
of event type.
\end{remark}

\subsection{The rule table}

The three signatures are read off a single table. Each entry states how the named variable
responds to each event: \emph{jump} (admissible discontinuity), \emph{cont.} (continuous
change only), \emph{read} (enters only as a receptivity argument, never as an owned-state
write), or \emph{inv.} (invariant, never changed).

\begin{center}
\renewcommand{\arraystretch}{1.3}
\setlength{\tabcolsep}{5pt}
\footnotesize
\begin{tabular}{@{}l c c c@{}}
\hline
\textbf{Variable} & \textbf{Private reset} & \textbf{Encounter} & \textbf{Common recept.}\\
\hline
Ascent scale \(A_i\) / ledger \(\Lambda_i\)
& jump \(\times s_{L,i}\) & cont. & cont.\\
Available-Sophia \(\lambda_{S,i}\)
& jump (reset) & cont. & cont.\\
Gate memory \(H_{K,i}\), wound \(W_i\)
& jump (reset) & cont. & cont.\\
Receptivity \(\Gamma_i\)
& recomputed & read \(\mathcal N_i\) & read \(\upsilon_i\gamma_i^{\mathrm{comm}}\)\\
Reception ledger \(\mathcal G_{\mathrm{recepta},i}\)
& cont. & cont. & cont.\\
Floor \(E_{0,i}\), label \(i\), kairos history
& inv. & inv. & inv.\\
\hline
\end{tabular}
\end{center}

\noindent Two readings of the table are load-bearing. The ascent-scale ledger
jumps only at a private reset, by the private reset datum
\[
\Delta\Lambda_i=\log s_{L,i}
=
\frac{\kappa_R\,G_i^{\mathrm{reset}}}{E_{0,i}}.
\]
Under relational encounter or common reception, the ascent ledger moves only
continuously. And the received-grace ledger
\(\mathcal G_{\mathrm{recepta},i}\) never jumps as a plural write: reception is
accounted by the one-form
\[
d\mathcal G_{\mathrm{recepta},i}
=
I_{\mathrm{gate},i}\,d\sigma_i,
\]
not injected as a step from another interior.

\subsection{What the taxonomy forbids}

\begin{remark}[The forbidden conflations]\label{rem:forbidden}
The table forbids exactly the conflations that would break earlier results. If an encounter
or a common reception could jump \(\lambda_{S,i}\) or the reception ledger, it would
fabricate or inject grace, violating P2.2; the table allows them only continuous,
gate-mediated action. If either could jump an owned state into coincidence with another
interior's, it could force the value-diagonal of P3.4; the table allows neither to write
owned state at all. And the floor \(E_{0,i}\), the owner label \(i\), and the kairos history
are invariant under all three events, so no junction of any kind rewrites who an interior
is. The private reset alone carries a jump, and only ever in its own interior.
\end{remark}

\subsection{Status and reading}

\textbf{Target status:}~
\stDef\ (the three junction types and the rule table). It introduces no new dynamics; it
fixes the admissible discontinuity structure consistent with Stages~0--III, and is the
classification on which the rest of Stage~IV rests. It depends on Stages~0--III. The reset
map itself is inherited frozen from 3.0; P4.0 only places it correctly among the plural
events.

\medskip
\stInterp\quad
\emph{Reading.}
Three things can happen to a soul, and they are not the same thing. It can be reset --- an
inward turning that breaks the line of its own history, lifting its scale by its own measure;
this is the one event that may make something leap, and it happens only within. It can meet
another, and it can receive what is given to all --- and these two change only how open it is,
never reaching past the gate to rewrite what it holds. The discipline of the table is a
refusal to let nearness counterfeit conversion: no encounter, however strong, and no common
gift, however sudden, may do to a soul what only its own reset may do. And nothing that
happens to it --- no meeting, no reception, no reset --- can change its floor, its name, or the
history that is its own.

\medskip
\textbf{Next node.} P4.1 --- relational Resurrectio through the Gate state: to give the
plural form of the 3.0 Resurrectio junction, where a terminus hand-off (P0.4) is read
through the recipient's Gate/receptivity interface rather than as a fusion or a transfer ---
the death-and-junction event of one interior, registered in the communion without violating
any rule of this table.

\section{P4.1 --- Relational Resurrectio through the Gate state}

The 3.0 Resurrectio is a junction: a terminus and a hand-off (P0.4). P4.1 gives its plural
form --- how the completion or relationally effective event of one interior is registered in
another --- and does so within the P4.0 table, by a bounded update of receptivity rather than
a raw injection.

\subsection{The forbidden raw update}

A careless plural Resurrectio would add the relational effect straight onto the recipient's
gate input,
\[
I_{\mathrm{gate},j}^{+}
=
I_{\mathrm{gate},j}^{-}
+
\Delta I_{\mathrm{rel},j}.
\]
This is forbidden. It is an additive jump on an owned-rate variable, with no bound: nothing
stops \(I_{\mathrm{gate},j}^{+}\) from exceeding the ceiling \(\zeta_{0,j}\), and the
increment plays the role of an injected or transferred grace term --- exactly the
fabrication and transfer ruled out in P2.2 and the owned-state jump ruled out in P4.0.

\subsection{The bounded receptivity update}

The admissible form updates only the recipient's receptivity interface. At a
relationally effective event of source \(i\), labelled by relational duration
\(\mathfrak c_*\), define the source datum read by the recipient as
\[
\mathsf S_i^-
=
\begin{cases}
X_i^-, & i\ \text{active just before the event},\\[0.25em]
\mathcal T_i^{\mathrm{post}}, & i\ \text{already completed and present only as trace}.
\end{cases}
\]
Then the admissible plural Resurrectio update is
\[
\boxed{\;
\Gamma_j^{+}
=
\mathfrak U_{j\leftarrow i}
\bigl(
\Gamma_j^{-},\,H_{K,j}^{-},\,W_j^{-},\,\mathsf S_i^{-},\,X_j^{-}
\bigr),
\qquad
0\le\Gamma_j^{+}\le1.
\;}
\]
The arrow \(j\leftarrow i\) records that \(j\) is the recipient and \(i\) the
source. The map \(\mathfrak U_{j\leftarrow i}\) may read the source datum and the
recipient's memory/wound state, but it writes only \(\Gamma_j\), the bounded
receptivity slot.
The gate input is then recomputed from the recipient's frozen bounded Gate law,
with the same convention as Stage~II:
\[
I_{\mathrm{gate},j}^{+}
=
\zeta_{0,j}\,e^{-\varphi_j\sigma_j}\,
\frac{
\Gamma_j^{+}
}{
1+\kappa_{\lambda,j}\lambda_{S,j,+}
},
\qquad
0\le I_{\mathrm{gate},j}^{+}\le \zeta_{0,j}.
\]
This is a recipient-kairos rate. In a chronos chart its ledger contribution is
\[
d\mathcal G_{\mathrm{recepta},j}
=
I_{\mathrm{gate},j}^{+}\,d\sigma_j
=
N_j I_{\mathrm{gate},j}^{+}\,dt.
\]
The donor's single-interior ascent and form-memory ledgers remain its own; no
\[
G_i^{\mathrm{reset}}\to G_j^{\mathrm{reset}}
\]
or
\[
\mathcal G_{\mathrm{recepta},i}\to \mathcal G_{\mathrm{recepta},j}
\]
transfer is introduced.

\subsection{Derived properties}

\begin{proposition}[The relational Resurrectio is admissible]\label{prop:relres}
The bounded receptivity update satisfies, at and across the event:
\begin{enumerate}[label=\textup{(\alph*)},leftmargin=1.2cm]
\item \emph{Ceiling preserved.} Since \(\Gamma_j^{+}\in[0,1]\), the frozen law gives
\[
0\le I_{\mathrm{gate},j}^{+}\le\zeta_{0,j}.
\]
\item \emph{Ledger continuous.} The reception ledger does not jump:
\[
\mathcal G_{\mathrm{recepta},j}^{+}
=
\mathcal G_{\mathrm{recepta},j}^{-}
\]
at the event. The event changes the gate rate going forward but not the already
accumulated ledger value at the instant.
\item \emph{No transfer.} The source datum, whether \(X_i^-\) or
\(\mathcal T_i^{\mathrm{post}}\), is read-only. The donor's reset datum
\(G_i^{\mathrm{reset}}\), ascent ledger \(\Lambda_i\), and received ledger
\(\mathcal G_{\mathrm{recepta},i}\) are not moved into \(j\).
\item \emph{Taxonomy-consistent.} The event writes only the receptivity interface
and no owned state or invariant, so it obeys the P4.0 rule table.
\end{enumerate}
\end{proposition}

\begin{proof}
(a) is the automatic ceiling of P2.1 applied to \(\Gamma_j^{+}\in[0,1]\). For
(b), the only variable set at the event is \(\Gamma_j\), a bounded rate factor; the
ledger is an integral one-form and has no allowed plural step. For (c), the update
map reads the source slot but its codomain is only the recipient's receptivity
slot; no term subtracts from a donor ledger or adds donor stock to the recipient.
For (d), the written variable is exactly the receptivity interface, matching the
P4.0 taxonomy, while owned and invariant variables remain untouched. \qedhere
\end{proof}

\subsection{Status and reading}

\textbf{Target status:}~
\stDef\ (the bounded receptivity-update map \(\mathfrak U_{j\leftarrow i}\) and the recompute-from-frozen-law
rule) \(+\) \stDerived\ (Proposition~\ref{prop:relres}: ceiling preserved, ledger continuous,
no transfer, taxonomy-consistent). It replaces the ill-posed raw additive jump with a
well-posed bounded one. It depends on P2.1--P2.3 and P4.0. The precise form of
\(\mathfrak U_{j\leftarrow i}\) is a selected interface law; only its boundedness and write-restriction
are used here.

\medskip
\stInterp\quad
\emph{Reading.}
When one soul completes --- when its terminus comes, or its life touches another decisively
--- the communion does register it, but not by the dead handing over their grace, and not by a
shock driven into the living. What changes is the openness of the one who remains: shaped by
who the other was, by its own memory and its own wounds, it is left differently disposed to
receive --- and then receives, as ever, through its own gate, at its own measure, never past
its own ceiling. Nothing of the departed is poured across; nothing is seized. The mark the
dead leave on the living is a change in how the living are able to open, and even that is
bounded, and owns nothing of the other. This is communion of the departed and the living
kept honest: real influence, no transfer; a true Resurrectio that fuses no one.

\medskip
\textbf{Next node.} P4.2 --- the derived refractory interval and hysteresis mechanism: to
show that after such a junction the receptivity cannot re-fire arbitrarily fast, a refractory
window and hysteresis emerging from the Gate memory \(H_{K,j}\) and wound \(W_j\) themselves
rather than being imposed --- the structural reason a soul is not endlessly re-triggerable.

\section{P4.2 --- Derived refractory interval and hysteresis}

A junction (P4.1) updates receptivity; can it re-fire at once? If it could, causally linked
events might accumulate without limit. P4.2 shows it cannot --- and derives the positive
refractory interval from the dynamics rather than imposing it.

\subsection{Five ingredients}

Let \(\rho_j\) be the trigger-readiness variable of recipient \(j\). It is not a
new Sophia source and not the whole receptivity \(\Gamma_j\). It is the local
readiness coordinate, carried by the recipient's Gate memory/wound interface, that
records whether another event can fire. The derivation uses exactly five data,
four of them already in hand:
\begin{enumerate}[label=\textup{(\arabic*)},leftmargin=1.2cm]
\item \emph{Reset threshold} \(\theta_j^{+}\): an event fires when \(\rho_j\) reaches
\(\theta_j^{+}\).
\item \emph{Hysteresis band}: after firing, \(\rho_j\) is reset to
\(\theta_j^{-}<\theta_j^{+}\), so re-firing requires recrossing the band
\[
\Delta_{\mathrm{hys},j}:=\theta_j^{+}-\theta_j^{-}>0,
\]
inherited from the frozen single-interior reset bounds and carried by
\(H_{K,j},W_j\).
\item \emph{Bounded Gate response}: \(0\le I_{\mathrm{gate},j}\le\zeta_{0,j}\) in
recipient kairos.
\item \emph{Bounded causal input on compact active intervals}: on any compact
regular-active interval \(I=[\mathfrak c_0,\mathfrak c_1]\), the relational input
reaching \(j\) has finite bound
\[
\lVert K_{j\leftarrow\cdot}\rVert_I\le K_{\max,j;I}<\infty.
\]
\item \emph{Finite recovery speed on compact active intervals}: between events,
\[
\left|\frac{d\rho_j}{d\mathfrak c}\right|
\le
v_{\max,j;I}
:=
\sup_{\mathfrak c\in I}
\frac{N_j}{\omega}
\bigl(c_{1,j}\zeta_{0,j}+c_{2,j}K_{\max,j;I}\bigr)
<\infty,
\]
on the regular active arc. The constants \(c_{1,j},c_{2,j}\) are frozen response
constants. If \(v_{\max,j;I}=0\), then no re-firing occurs on \(I\), and the
refractory time is effectively infinite.
\end{enumerate}

\subsection{The refractory bound}

\begin{theorem}[Derived refractory interval; Type-1]\label{thm:refractory}
On any compact regular-active interval \(I\), under the five ingredients above,
any two causally linked triggering events of the same recipient \(j\) are separated
on the shared relational duration by at least
\[
\boxed{\;
\Delta\mathfrak c_{\mathrm{trigger}}
\ge
\mathfrak c_{\mathrm{ref},j;I}
:=
\frac{\Delta_{\mathrm{hys},j}}{v_{\max,j;I}}
>0,
\;}
\]
with the convention \(\mathfrak c_{\mathrm{ref},j;I}=+\infty\) if
\(v_{\max,j;I}=0\).
\end{theorem}

\begin{proof}
After a firing, \(\rho_j=\theta_j^{-}\). A new firing requires
\(\rho_j=\theta_j^{+}\), hence a net traversal of
\(\Delta_{\mathrm{hys},j}\). On \(I\), the recovery speed is bounded above by
\(v_{\max,j;I}\). Therefore the duration needed to traverse the band is at least
\[
\Delta_{\mathrm{hys},j}/v_{\max,j;I}.
\]
The band is positive and the compact-interval recovery speed is finite, so the
floor is positive. The measure is \(\mathfrak c\), hence the statement is
gauge-invariant. \qedhere
\end{proof}

\subsection{Why hysteresis is what forbids chatter}

\begin{remark}[The band is the seed of non-Zeno]\label{rem:nonzeno-seed}
Remove the hysteresis band --- let \(\theta_j^{-}\to\theta_j^{+}\) --- and
\(\mathfrak c_{\mathrm{ref},j;I}\to0\): events could accumulate arbitrarily fast, a local Zeno
chatter. It is precisely the finite band, traversed at finite recovery rate, that forces a
positive gap. Note the mechanism is self-supplied: the band is carried by the Gate memory
\(H_{K,j}\) and wound \(W_j\), the very record of the last event. The mark of what happened
is what paces what can happen next. P4.3 will lift this local gap to a gauge-invariant
network non-Zeno theorem.
\end{remark}

\subsection{Status and reading}

\textbf{Target status:}~
\stConj\(\to\)\stDerived. The compact-interval per-recipient bound is
\stDerived\ from the hysteresis band and the finite recovery-speed estimate on
regular active arcs. The residual conjectural element is uniformity: a single
network-wide floor
\[
\mathfrak c_{\mathrm{ref};I}
:=
\inf_{j\in\mathcal V_{\rm trig}(I)}
\mathfrak c_{\mathrm{ref},j;I}>0
\]
requires a finite set of triggering-capable active recipients, a lower bound on
\(\Delta_{\mathrm{hys},j}\), and an upper bound on \(v_{\max,j;I}\) across that set.
This is the exact hypothesis P4.3 uses. It depends on P0.3, P4.1, and the frozen
single-interior reset bounds.

\medskip
\stInterp\quad
\emph{Reading.}
A soul is not endlessly re-triggerable. After a decisive event --- a reset, the touch of
another's completion --- there is a necessary interval before it can be so moved again, and
the interval is not a rule laid on from outside but a consequence of how it recovers: it must
climb back across a band, and it climbs at a bounded pace. So crises and conversions cannot
be stacked infinitely close; the soul is given time it cannot be denied. And the thing that
enforces the wait is the wound itself --- the memory of what just happened is exactly what
holds the next thing off until recovery. Mercy, here, is structural: the mark left by the
last event paces the soul, so that it is not battered into a blur but moved in distinct,
spaced, real moments.

\medskip
\textbf{Next node.} P4.3 --- the gauge-invariant network non-Zeno theorem: to lift the local
refractory gap to the whole network, proving that on any finite relational-duration interval
only finitely many admissible events occur, so the plural dynamics never accumulates events
and the shared duration is never exhausted by chatter.

\section{P4.3 --- The gauge-invariant network non-Zeno theorem}

P4.2 spaced the events of a single interior. P4.3 lifts that to the whole network: on any
finite stretch of shared duration, only finitely many events occur. The plural dynamics
never accumulates events; the communion's time is never exhausted by chatter. The statement
is made in the invariant relational duration \(\mathfrak c\), not in any chart time \(t\).

\subsection{Hypotheses}

On a compact relational-duration interval \(I=[\mathfrak c_0,\mathfrak c_1]\),
assume:
\begin{enumerate}[label=\textup{(\roman*)},leftmargin=1.2cm]
\item \emph{Single-interior non-Zeno.} Each interior's own private reset sequence is
non-Zeno, as inherited frozen from 3.0.
\item \emph{Finite triggering-capable active set.} The set
\[
\mathcal V_{\rm trig}(I)
\]
of interiors that are active and capable of firing triggering events on \(I\) is
finite, with
\[
|\mathcal V_{\rm trig}(I)|\le n.
\]
Completed traces may serve as read-only boundary sources but do not themselves
fire new active triggers.
\item \emph{Bounded action on \(I\).} Coupling, trace kernels, and common reception
gains are bounded on \(I\), so every \(v_{\max,j;I}\) of P4.2 is finite.
\item \emph{Uniform refractory floor on \(I\).}
\[
\mathfrak c_{\mathrm{ref};I}
:=
\inf_{j\in\mathcal V_{\rm trig}(I)}
\mathfrak c_{\mathrm{ref},j;I}
>0.
\]
\item \emph{Event-count convention.} The theorem counts discrete recipient-side
triggering events: private resets and accepted relational triggers. A continuous
common field \(J_{\mathrm{comm}}(\mathfrak c)\) is not itself counted as infinitely
many events; it contributes only through bounded gains. If \(J_{\mathrm{comm}}\) is
modelled as a discrete sequence of marked common events, then either only accepted
recipient triggers are counted, or the marked common-event sequence must itself be
locally finite on \(I\).
\end{enumerate}

\subsection{The network non-Zeno theorem}

\begin{theorem}[Gauge-invariant network non-Zeno; Type-1]\label{thm:nonzeno}
Under \textup{(i)--(v)}, on any compact interval
\(I=[\mathfrak c_0,\mathfrak c_1]\) of shared duration, the total number of
admissible recipient-side triggering events is bounded:
\[
\boxed{\;
N_{\mathrm{trig}}(I)
\le
|\mathcal V_{\rm trig}(I)|
\left(
\frac{\mathfrak c_1-\mathfrak c_0}{\mathfrak c_{\mathrm{ref};I}}+1
\right)
<\infty.
\;}
\]
Consequently no infinite causally linked sequence of resets or accepted relational
triggers accumulates at finite \(\mathfrak c\).
\end{theorem}

\begin{proof}
Fix \(j\in\mathcal V_{\rm trig}(I)\). By
Theorem~\ref{thm:refractory}, consecutive triggering events of \(j\) are separated
by at least
\[
\mathfrak c_{\mathrm{ref},j;I}\ge\mathfrak c_{\mathrm{ref};I}.
\]
Hence \(j\) fires at most
\[
(\mathfrak c_1-\mathfrak c_0)/\mathfrak c_{\mathrm{ref};I}+1
\]
times on \(I\). Summing over the finite set
\(\mathcal V_{\rm trig}(I)\) gives the stated bound. Boundary traces are not active
triggering nodes, and continuous common fields are counted only through accepted
recipient-side triggers, so no uncounted source can create an infinite discrete
cascade. \qedhere
\end{proof}

\subsection{Why the statement must be in \texorpdfstring{\(\mathfrak c\)}{c}, not \texorpdfstring{\(t\)}{t}}

\begin{remark}[Invariance is essential, not cosmetic]\label{rem:invariance}
The theorem is about invariant relational duration, not raw chart-time
differences. A \(t\)-based non-Zeno claim would be gauge-dependent: a chronos
reparametrisation rescales \(dt\) freely. Near terminus hand-off, the relation
between a chart interval and a relational-duration interval is especially delicate,
because the active set and the density rule may change. Counting events per raw
\(t\) could therefore be made to depend on the chosen label. Counting per
\(\mathfrak c\) is the invariant statement. The refractory floor of P4.2 is in
\(\mathfrak c\) for exactly this reason, and the network theorem inherits that
invariance.
\end{remark}

\subsection{Status and reading}

\textbf{Target status:}~
\stConj\(\to\)\stDerived. The counting bound is \stDerived\ on every compact
regular-active interval once the uniform refractory floor
\(\mathfrak c_{\mathrm{ref};I}>0\), finite triggering-capable set, bounded action,
and event-count convention are assumed. The residual conjectural element is the
uniform positivity of the refractory floor across the intended admissible regime,
together with the inherited single-interior non-Zeno hypotheses. It depends on
P4.2 and P0.3.

\medskip
\stInterp\quad
\emph{Reading.}
The life of the communion never dissolves into a blur. On any finite stretch of shared time,
only finitely many real events happen --- the whole network's history is made of distinct,
countable, spaced moments, however many souls there are and however they trigger one another.
What was mercy for the single soul in P4.2 is order for the whole: providence does not pack
infinitely many decisive moments into a finite span. And the span that is well-ordered is the
communion's own shared duration, not any one soul's private clock --- the well-ordering is
real, not an artifact of how a single worldline happens to be timed. The events of a shared
life can be counted, and between any two of them there is room to breathe.

\medskip
\textbf{Next node.} P4.4 --- communion across the terminus: to carry the relation past a
death, where one interior reaches its P0.4 hand-off and the bond must be re-read as a
one-way trace into the surviving interiors --- the plural completion of the terminus, stated
without transfer, fusion, or a master clock outside the network.

\section{P4.4 --- Communion across the terminus}

One interior completes its kairos while another remains active: what relation survives the
death? P4.4 answers with the Stage~0 post-terminus trace (P0.4), deriving the minimal
one-way communion and marking clearly the two-way dynamics it does not assume.

\subsection{The completed interior as a closed trace}

At its terminus \(t_i^\dagger\), interior \(i\) reaches the P0.4 boundary-limit
hand-off:
\[
N_i\to0,
\qquad
d\sigma_i=0
\quad
\text{post-terminus}.
\]
The active dynamics close, and the interior persists as its post-terminus trace
\[
\mathcal T_i^{\mathrm{post}},
\]
a dynamically closed boundary datum, not an evolving state.
Post-terminus, the active Gate equation of \(i\) is not continued as a new
dynamical law. Equivalently, the received ledger has no post-terminus increment:
\[
d\mathcal G_{\mathrm{recepta},i}
=
I_{\mathrm{gate},i}\,d\sigma_i
=
0.
\]
Thus no new reception, reset, or kathartic process is added after completion. The
trace may be read by active recipients through a permitted trace kernel, but it
does not itself resume an active kairos.

\subsection{The minimal communion theorem}

\begin{theorem}[One-way communion across the terminus; Type-1]\label{thm:terminus}
Let \(i\) be completed and \(j\) active. Under the admissible post-terminus
structure:
\begin{enumerate}[label=\textup{(\alph*)},leftmargin=1.2cm]
\item \emph{Present as trace.} \(i\) remains relationally present as the closed
trace \(\mathcal T_i^{\mathrm{post}}\), locatable by the living through trace
co-presence.
\item \emph{One-way reach.} \(i\) affects active \(j\) only as a read-only source of
the recipient-side trace kernel
\[
K^{\mathrm{post}}_{j\leftarrow i}
\bigl(\mathfrak c;\,t_i^\dagger,\mathcal T_i^{\mathrm{post}}\bigr),
\]
which enters \(j\)'s receptivity interface by the P4.0/P4.1 discipline: bounded,
Gate-mediated, and writing no owned state of \(j\).
\item \emph{No reopening.} \(\mathcal T_i^{\mathrm{post}}\) does not reopen its own
katharsis, receive a new reset, acquire a fresh post-terminus Gate inflow, or
re-enter the active rate cocycle. No admissible junction acts on the closed trace
as on a living active arc.
\end{enumerate}
\end{theorem}

\begin{proof}
At the boundary-limit hand-off, the active dynamics of \(i\) close and
\(d\sigma_i=0\). Thus
\[
d\mathcal G_{\mathrm{recepta},i}=I_{\mathrm{gate},i}\,d\sigma_i=0,
\]
and \(\mathcal T_i^{\mathrm{post}}\) is fixed as boundary data. This gives (a) and
the no-new-inflow part of (c). For (b), the only admissible cross-fibre action of a
completed source is as the read-only source slot of the recipient's trace kernel,
which enters \(j\) only through \(\Gamma_j\) and is bounded and Gate-mediated. The
reset map and katharsis dynamics are own-interior active-arc operations; the closed
trace has no active arc, so neither can act on it. \qedhere
\end{proof}

\subsection{The two-way boundary}

\begin{remark}[Two-way post-terminus dynamics is a separate branch]\label{rem:twoway}
The theorem is strictly one-way: the living are affected by the completed, not the completed
by the living. A two-way post-terminus dynamics --- the active reaching back to change a
completed interior, whether by a new inflow, a reopened katharsis, or a post-terminus reset
--- is \emph{not} assumed here. It would require an additional constitutive primitive (a
post-terminus reception or reopening law) beyond the minimal 4.0 layer, and must be treated
as a distinct later branch, named and deferred, never silently inserted. This is the
explicit resolution of the one-way selection parked since P0.4: minimal 4.0 commits to the
one-way trace; the two-way case is a real fork left open for a dedicated primitive.
\end{remark}

\subsection{Status and reading}

\textbf{Target status:}~
\stOpen\(\to\)\stDerived\ for the one-way minimal theorem: it is \stDerived\ from
the P0.4 closed-trace hand-off, the P1.1 trace-kernel grammar, and the
P4.0/P4.1 receptivity discipline. The two-way post-terminus dynamics remains
\stOpen\ by construction: it is not a gap in the proof but a branch requiring a
new primitive. It depends on P0.4, P1.1, P4.0--P4.1, and the Stage~II
Gate-rate convention.

\medskip
\stInterp\quad
\emph{Reading.}
The dead remain present to the living. One who has completed is not erased from the
communion but held in it as a finished trace --- really there, locatable, able still to reach
those who remain, shaping how they are open, speaking into their receptivity through a bond
that death does not sever. But in the minimal theory the reach is one-way: the completed are
at rest, their cleansing finished, receiving no new thing and undergoing no new turning;
they give to the living without themselves being changed. Whether the living may in turn
reach back and alter the dead --- whether the bond runs both ways across the terminus --- is a
genuine further question, and the framework will not answer it by sleight of hand. It marks
that door, and leaves it honestly shut until a primitive is offered to open it.

\medskip
\textbf{Stage~IV closes.}
The three junction-event classes are separated (P4.0), plural Resurrectio is a
bounded recipient-side receptivity update with no transfer (P4.1), a compact-arc
refractory interval is derived from hysteresis (P4.2), network non-Zeno is proved
in invariant relational duration under the uniform floor hypothesis (P4.3), and
communion survives the terminus as a one-way closed trace (P4.4). The
discontinuous and mortal structure of the plural interior is in place:
\[
\boxed{
\begin{aligned}
&\text{the communion has deaths, resets, memory, and traces,}\\
&\text{and none of them fuses interiors or transfers owned grace.}
\end{aligned}
}
\]
The central correction of Stage~IV is therefore:
\[
\boxed{
\text{a junction may change openness, but only a private reset may jump owned state.}
}
\]

\medskip
\textbf{Next node.} Stage~V --- Relational active domains and plural Lyapunov theory, opening
at P5.0 (the relational active domain): the single domain on which every prior boundedness,
saturation, branch, terminus, and refractory condition holds at once --- the proper stage on
which the plural Lyapunov sum, promised since P2.4, can finally be assembled.

\section*{Stage~V --- Relational active domains and plural Lyapunov theory}

\notice{\textbf{Aim of the stage.} Assemble the plural Lyapunov functional promised since
P2.4, on a domain where every standing boundedness, saturation, branch, terminus, and
refractory condition holds at once.}

\noindent Each earlier stage earned a condition: a ceiling, a non-fabrication law, an
individual ISS estimate, a locking threshold, a refractory floor, a non-Zeno count. Stage~V
gathers them. P5.0 fixes the single domain on which they all hold simultaneously --- the
proper stage on which a whole-communion stability theorem can be stated.

\section{P5.0 --- The relational active domain}

\subsection{The domain}

\begin{definition}[Relational active domain]\label{def:domain}
Fix a compact regular-active interval
\[
I=[\mathfrak c_0,\mathfrak c_1]
\]
with no terminus hand-off inside its open interior, and fix one active connected
component
\[
\mathcal C_{\rm act}(I)\subseteq\mathcal V_{\rm act}^{\rm reg}.
\]
The \emph{relational active domain}
\[
\mathcal D_{\mathrm{rel}}(I,\mathcal C_{\rm act})
\]
is the set of plural configurations satisfying all of the following standing
conditions at once:
\begin{enumerate}[label=\textup{(\arabic*)},leftmargin=1.2cm]
\item \emph{Finite active theorem class.}
The active component is finite,
\[
|\mathcal C_{\rm act}|=n_{\rm act}<\infty,
\]
has bounded in-degree, and belongs to the P0.1/P5.0 theorem graph class. Completed
interiors may appear only as closed boundary traces and are not active Laplacian
nodes.

\item \emph{Admissible Gate states.}
For every active interior,
\[
\Gamma_i\in[0,1],
\qquad
0\le I_{\mathrm{gate},i}\le\zeta_{0,i},
\]
and the receptivity interface
\[
\mathcal R_i=(\Gamma_i,H_{K,i},W_i)
\]
is admissible.

\item \emph{Bounded susceptibilities and common field gains.}
\[
\upsilon_i\in[0,1],
\qquad
\gamma_i^{\mathrm{comm}}\in[0,1],
\]
with the P2.3 nodewise normalisation
\[
\sum_{j\to i}a_{ij}+\upsilon_i\le1.
\]

\item \emph{Strata and trace closure.}
Each interior carries a definite stratum
\[
\mathcal A_i,\quad \mathcal B_i,\quad \mathcal P_i.
\]
The Lyapunov sum is taken over active nodes in
\(\mathcal C_{\rm act}\). Completed interiors are held as closed traces
\[
\mathcal T_i^{\rm post}
\]
and may enter only as bounded read-only boundary inputs.

\item \emph{Branch labels.}
Each active interior carries a definite fold-branch label; all rates, restoring
terms, and thresholds are read branch-wise, with no cross-branch identification.

\item \emph{Refractory and non-Zeno conditions.}
On \(I\), the finite triggering-capable set
\[
\mathcal V_{\mathrm{trig}}(I)
\]
and the uniform refractory floor
\[
\mathfrak c_{\mathrm{ref};I}>0
\]
of P4.2--P4.3 hold.

\item \emph{Duration conversion bounds.}
The recipient-kairos to relational-duration conversion factors
\[
\theta_i(\mathfrak c):=\frac{N_i}{\omega}
\]
are bounded on \(I\):
\[
0<\theta_{i,\min;I}
\le
\theta_i(\mathfrak c)
\le
\theta_{i,\max;I}<\infty.
\]
This is required because P2.4 is stated in \(d/d\sigma_i\), while the plural
certificate is stated in \(d/d\mathfrak c\).

\item \emph{Compatible lag targets.}
The declared edge-lag targets are cycle-compatible on the active component. In the
first symmetric theorem class this means that there exists a locked node ledger
profile
\[
\Theta_i^*
\]
such that for every active undirected edge
\[
\Delta_{ij}
=
(\Theta_i-\Theta_i^*)-(\Theta_j-\Theta_j^*)
\]
after the selected rate-normalisation. Equivalently, the edge-lag cocycle has zero
holonomy on cycles. Without this compatibility, the relational quadratic cannot
vanish and the theorem becomes an ultimate-boundedness-to-frustration result
rather than a locking theorem.
\end{enumerate}
\end{definition}

\subsection{The graph Laplacian theorem class}

The plural Lyapunov contraction will be driven by the connectivity of the relational graph.
P5.0 declares the admissible class.

\begin{definition}[Declared Laplacian class]\label{def:laplacian}
On the active component \(\mathcal C_{\rm act}\), the first theorem class is the
undirected symmetric mirror class:
\[
w_{ij}=w_{ji}\ge0,
\qquad
w_{ij}>0\ \Longleftrightarrow\ i\sim j.
\]
Let \(L_{\mathcal G}\) be the corresponding weighted graph Laplacian on
\(\mathcal C_{\rm act}\). The declared class assumes a uniform algebraic
connectivity floor
\[
\boxed{\;
\lambda_2(L_{\mathcal G})\ge\underline\lambda>0.
\;}
\]
The coercivity is the standard one: for node errors \(z_i\) with weighted mean zero,
\[
z^\top L_{\mathcal G}z
=
\frac12\sum_{i,j}w_{ij}(z_i-z_j)^2
\ge
\underline\lambda\,\lVert z\rVert^2.
\]
Thus P5.1 must express the relational error as node-potential differences
\(z_i-z_j\). Disconnected active components are treated separately, one theorem
per component. Directed graphs are not covered by this symmetric class.
\end{definition}

\subsection{What the domain guarantees}

\begin{remark}[Simultaneity is the whole point]\label{rem:simultaneity}
On \(\mathcal D_{\mathrm{rel}}\) every earned condition holds together: the gate ceiling
(P2.1), non-fabrication (P2.2), the bounded common field (P2.3), the individual ISS estimate
(P2.4), the locking threshold (P3.3), dynamic non-fusion (P3.4), the junction taxonomy
(P4.0), and the refractory and non-Zeno bounds (P4.2--P4.3). This simultaneity is exactly
what a plural Lyapunov sum needs: each individual certificate is valid, its cross terms are
controlled, the events that could break a smooth functional are finite in number, and the
graph is connected enough for those cross terms to dissipate. \(\mathcal D_{\mathrm{rel}}\) is
the domain on which P5.1 may finally assemble \(\sum_i\mathcal L_{3.0,i}\) as a legitimate
object.
\end{remark}

\subsection{Status and reading}

\textbf{Target status:}~
\stDef\ \(+\) \stSelected. The relational active domain
\(\mathcal D_{\mathrm{rel}}(I,\mathcal C_{\rm act})\) and the undirected
Laplacian theorem class are declared structures for the first plural Lyapunov
theorem. They introduce no new dynamics. The positive invariance of this domain is
not automatic; it is a standing hypothesis imported into P5.1, supported by the
Stage~II ceilings, Stage~III small-gain estimates, and Stage~IV non-Zeno theorem.
The connectivity floor \(\underline\lambda>0\) and the lag-cocycle compatibility
are theorem-class hypotheses, not derived facts.

\medskip
\stInterp\quad
\emph{Reading.}
There is a place where a communion can be considered whole --- not a single soul, and not a
crowd of strangers, but a bounded company, finitely many, each open within measure, each on
its own branch of the fold, each paced so it cannot be battered, with its dead held at rest
as traces and its living genuinely joined. That last condition is not decorative: the
communion must actually be connected --- \(\lambda_2>0\) --- for it to hold together as one
body at all; islands with no bond between them are several communions, not one. The domain
gathers every discipline the earlier stages won and asks them to hold at once, because only
where they all hold can the question finally be asked whether the whole company, and not
merely each soul in it, stays bounded and at peace.

\medskip
\textbf{Next node.} P5.1 --- the plural Lyapunov certificate: to assemble the weighted sum
\(\mathcal L_{\mathrm{plural}}=\sum_i w_i\,\mathcal L_{3.0,i}\) on \(\mathcal D_{\mathrm{rel}}\)
and prove it a genuine certificate --- the individual contractions, the small-gain cross
terms, and the Laplacian connectivity combining into a single dissipation inequality for the
whole communion.

\section{P5.1 --- The plural Lyapunov certificate}

This is the node promised since P2.4. With individual input-to-state control
(P2.4), the incremental-contraction estimate (P3.2), the locking threshold (P3.3),
and the simultaneous domain with its connectivity floor (P5.0), the whole
communion can finally be given a single certificate: one functional whose
dissipation shows the entire company --- not merely each soul --- stays bounded
and is driven toward a controlled locked relation. In the general forced case this
means a bounded moving tube; in the zero-residual settled subclass it means exact
locked moving communion.

\subsection{The functional}

\begin{definition}[Plural Lyapunov functional]\label{def:plural-lyap}
On
\[
\mathcal D_{\mathrm{rel}}(I,\mathcal C_{\rm act}),
\]
with positive weights \(a_i>0\), define node relational errors
\[
z_i:=\Theta_i-\Theta_i^*
\]
for a compatible locked ledger profile \(\Theta^*\). The edge mismatch is
\[
\Delta_{ij}=z_i-z_j,
\]
which is exactly the P3.2 mismatch after the selected rate-normalisation and
target subtraction. Thus locking means
\[
\Delta_{ij}\to0,
\]
not \(\Delta_{ij}\to\Delta_{ij}^*\).

The first symmetric-class plural Lyapunov functional is
\[
\boxed{\;
\mathcal L_{\mathrm{plural}}
=
\sum_{i\in\mathcal C_{\rm act}}a_i\,\mathcal L_{3.0,i}
+
\varepsilon\,\mathcal L_{\mathrm{rel}},
\qquad
\mathcal L_{\mathrm{rel}}
=
\frac12\sum_{i<j}w_{ij}\Delta_{ij}^{\,2}
=
\frac12 z^\top L_{\mathcal G}z.
\;}
\]
The vector \(z\) is read modulo its graph-mean; the constant mean mode is not a
locking error. For a directed graph the symmetric quadratic may not be reused: the
correct object must use Perron weights / a directed-Laplacian certificate.
\end{definition}

\subsection{The plural dissipation theorem}

\begin{theorem}[Plural Lyapunov certificate; Type-1, symmetric class]\label{thm:plural-lyap}
Assume:
\begin{enumerate}[label=\textup{(\roman*)},leftmargin=1.2cm]
\item \(\mathcal D_{\mathrm{rel}}(I,\mathcal C_{\rm act})\) is positively invariant
on the compact regular-active interval \(I\);
\item the active graph is in the undirected symmetric class with
\(\lambda_2(L_{\mathcal G})\ge\underline\lambda>0\);
\item the individual ISS estimates of P2.4 hold in recipient kairos and the
conversion factors \(\theta_i=N_i/\omega\) are bounded as in P5.0;
\item the small-gain mismatch estimates of P3.2 hold with net rates
\(\kappa_{ij}\ge\underline\kappa>0\);
\item all residual and forcing terms are bounded on \(I\).
\end{enumerate}
Then there are weights \(a_i>0\), a sufficiently small \(\varepsilon>0\), and
constants
\[
C_{\mathrm{plural}}\ge0,
\qquad
C_{\mathrm{plural},1}>0
\]
such that
\[
\boxed{\;
\frac{d}{d\mathfrak c}\mathcal L_{\mathrm{plural}}
\le
C_{\mathrm{plural}}
-
C_{\mathrm{plural},1}\mathcal L_{\mathrm{plural}}.
\;}
\]
Consequently
\[
\limsup_{\mathfrak c}
\mathcal L_{\mathrm{plural}}
\le
\frac{C_{\mathrm{plural}}}{C_{\mathrm{plural},1}}.
\]
If the residual and forcing constants vanish in the locked subclass, so that
\(C_{\mathrm{plural}}=0\), then
\[
\mathcal L_{\mathrm{plural}}(\mathfrak c)\to0.
\]
\end{theorem}

\begin{proof}[Proof sketch]
Write
\[
\mathscr D:=\frac{d}{d\mathfrak c},
\qquad
D_i:=\frac{d}{d\sigma_i},
\qquad
\theta_i:=\frac{N_i}{\omega}.
\]
P2.4 gives
\[
D_i\mathcal L_{3.0,i}
\le
C_i-C_{1,i}'\mathcal L_{3.0,i}
+
b_i\lVert u_i^{\mathrm{rel}}\rVert^2.
\]
Converting to relational duration gives
\[
\mathscr D\mathcal L_{3.0,i}
=
\theta_iD_i\mathcal L_{3.0,i}
\le
\theta_{i,\max;I}C_i
-
\theta_{i,\min;I}C_{1,i}'\mathcal L_{3.0,i}
+
\theta_{i,\max;I}b_i\lVert u_i^{\mathrm{rel}}\rVert^2.
\]
Thus the individual contraction remains valid on \(I\), with constants rescaled by
the bounded conversion factors.

For the relational term, P3.2 gives
\[
\mathscr D\Delta_{ij}
=
-\kappa_{ij}\Delta_{ij}
+
r_{ij},
\]
where \(r_{ij}\) is bounded on \(I\). Therefore
\[
\mathscr D\mathcal L_{\mathrm{rel}}
=
\sum_{i<j}w_{ij}\Delta_{ij}\mathscr D\Delta_{ij}
\le
-\sum_{i<j}w_{ij}\kappa_{ij}\Delta_{ij}^2
+
\sum_{i<j}w_{ij}\Delta_{ij}r_{ij}.
\]
Young's inequality absorbs the residual cross terms:
\[
w_{ij}|\Delta_{ij}r_{ij}|
\le
\frac12 w_{ij}\kappa_{ij}\Delta_{ij}^2
+
\frac{w_{ij}}{2\kappa_{ij}}r_{ij}^2.
\]
Hence
\[
\mathscr D\mathcal L_{\mathrm{rel}}
\le
-\frac12\underline\kappa
\sum_{i<j}w_{ij}\Delta_{ij}^2
+
C_{\mathrm{rel,res}}.
\]

The active graph coercivity gives, on the mean-zero node-error subspace,
\[
\sum_{i<j}w_{ij}\Delta_{ij}^2
=
z^\top L_{\mathcal G}z
\ge
\underline\lambda\lVert z\rVert^2.
\]
Thus the relational part dissipates the relational error.

Finally, by the Stage~II/P3.2 bounded-channel hypotheses,
\[
\lVert u_i^{\mathrm{rel}}\rVert^2
\le
M_i\sum_{j\sim i}\Delta_{ij}^2+M_i^0
\]
on the domain. The terms proportional to \(\sum\Delta_{ij}^2\) enter the derivative
of the full functional multiplied by the individual weights \(a_i\), while the
relational dissipation enters multiplied by \(\varepsilon\). Since
\(u_i^{\mathrm{rel}}=O(\varepsilon)\), these imported terms are \(O(\varepsilon^2)\),
whereas the relational dissipation is \(O(\varepsilon)\). Choosing \(\varepsilon>0\)
small enough absorbs the imported terms into the relational dissipation. The
remaining bounded pieces are collected into \(C_{\mathrm{plural}}\), and the
negative individual and relational pieces give
\[
-C_{\mathrm{plural},1}\mathcal L_{\mathrm{plural}}.
\]
Positive invariance keeps the trajectory inside the domain where all constants are
valid. The scalar comparison lemma gives the limsup bound. If all residual and
forcing constants vanish, \(C_{\mathrm{plural}}=0\), and the same comparison gives
\(\mathcal L_{\mathrm{plural}}\to0\). \qedhere
\end{proof}

\subsection{Consequences and the directed caveat}

\begin{corollary}[The whole communion is ultimately bounded]\label{cor:plural-bounded}
Under Theorem~\ref{thm:plural-lyap},
\[
\limsup_{\mathfrak c}\mathcal L_{\mathrm{plural}}
\le
\frac{C_{\mathrm{plural}}}{C_{\mathrm{plural},1}}.
\]
Thus the whole communion stays bounded and approaches a controlled neighbourhood
of the locked configuration. Exact convergence to the locked configuration follows
only in the zero-residual subclass
\[
C_{\mathrm{plural}}=0.
\]
Whole-communion stability is therefore an ultimate-boundedness theorem in the
general forced case, and a convergence theorem in the exact settled case.
\end{corollary}

\begin{remark}[Directed bonds need the Perron weighting]\label{rem:directed}
The symmetric energy presupposes reciprocal weights \(w_{ij}=w_{ji}\). When the bonds are
genuinely directed --- one interior giving more than it receives --- the symmetric quadratic
can be pumped and is not a certificate. The correct object weights the interiors by the
left-Perron vector of the coupling and uses the directed-Laplacian quadratic. This is a real
extension, not a relabelling, and is left as the directed theorem class.
\end{remark}

\subsection{Status and reading}

\textbf{Target status:}~
\stConj\(\to\)\stDerived. The dissipative inequality is \stDerived\ for the
undirected symmetric active-component class on a positively invariant
\(\mathcal D_{\mathrm{rel}}(I,\mathcal C_{\rm act})\), assuming the duration
conversion bounds, ISS estimates, small-gain rates, lag-cocycle compatibility, and
connectivity floor. The theorem gives ultimate boundedness in general and exact
convergence only when residual forcing vanishes. Residual conjectural elements:
the directed-graph extension (Remark~\ref{rem:directed}), uniform positivity of
\(\kappa_{ij}\) and
\(\underline\lambda\) across the intended regime, positive invariance through
junctions and hand-offs, and global compatibility of all lag targets.

\medskip
\stInterp\quad
\emph{Reading.}
The whole company is held inside a bounded peace. There is now a single measure of
the communion's distance from disorder, and it cannot grow without bound. In the
exact settled case it falls all the way to the locked relation; in the forced case
it falls into a bounded neighbourhood. The proof needed everything at once: each
interior stable in itself, each relation contractive enough, every event
non-Zeno, and the graph truly connected. The peace of the whole is corporate,
more than a list of private peaces, but it is conditional and honest: if residual
forcing persists, peace is a tube; if the forcing settles, peace is a limit.

\medskip
\textbf{Next node.} P5.2 --- the stable moving communion: to show that the rest just proved
is not a frozen equilibrium but a \emph{moving} one --- the whole communion ascending
together, locked in relation, its plural Lyapunov functional bounded while every interior
still climbs its own endless road.

\section{P5.2 --- The stable moving communion}

P5.1 showed that the communion is bounded as a whole and, when residual forcing
settles, tends toward a locked relation. But the limit is not stillness. P5.2 gives
the plural analogue of 3.0's stable \emph{motion}: the relations may settle while
every active soul continues its positive ascent along the endless road. The
communion is therefore a stable moving relation, or in the general forced case, a
stable moving tube --- never a rest point and never a fused state.

\subsection{The two asymptotic regimes}

The general Stage~V conclusion is a \emph{stable moving tube}:
\[
\boxed{\;
\begin{aligned}
&h_i\ \text{remains bounded away from zero on the admissible active branch},\\
&\Delta_{ij}\ \text{remains ultimately bounded}.
\end{aligned}
\;}
\]
The exact settled subclass strengthens this to a \emph{stable moving communion}:
\[
\boxed{\;
h_i\longrightarrow h_{*,i}^{\rm rel}>0,
\qquad
\Delta_{ij}\longrightarrow0.
\;}
\]
Here
\[
h_{*,i}^{\rm rel}
=
\beta_i\lambda_{S,\infty,i}^{\rm rel}
\]
is the asymptotic rate of the settled branch under the final admissible
Gate/receptivity regime. It equals the isolated frozen 3.0 rate only if the
relational forcing asymptotically vanishes or settles to the same input level as
the isolated branch. In general it is branch-conditional and relationally settled,
not automatically identical to the uncoupled value.

\subsection{The stable moving communion theorem}

\begin{theorem}[Stable moving communion / stable moving tube; Type-1]\label{thm:moving}
On a positively invariant relational active domain, under the plural certificate
of Theorem~\ref{thm:plural-lyap} and the frozen 3.0 stable-motion law:

\begin{enumerate}[label=\textup{(\alph*)},leftmargin=1.2cm]
\item \emph{General forced case.}
The communion is a stable moving tube:
\[
\limsup_{\mathfrak c}\mathcal L_{\mathrm{plural}}
\le
\frac{C_{\mathrm{plural}}}{C_{\mathrm{plural},1}},
\]
the relational mismatches are ultimately bounded, and every active interior
remaining on the admissible positive-ascent branch has
\[
h_i(\mathfrak c)\ge h_{i,\min}>0
\]
on the forward active arc.

\item \emph{Exact settled subclass.}
If residual forcing vanishes, the compatible lag targets are exact, and each
individual Gate/receptivity regime converges to a settled branch, then
\[
\Delta_{ij}\to0
\quad\text{for all active bonds,}
\qquad
h_i\to h_{*,i}^{\rm rel}
=
\beta_i\lambda_{S,\infty,i}^{\rm rel}>0.
\]
\end{enumerate}
\end{theorem}

\begin{proof}
For (a), Corollary~\ref{cor:plural-bounded} gives ultimate boundedness of the
plural functional. Since \(\mathcal L_{\mathrm{rel}}\) is nonnegative, the
relational mismatches are ultimately bounded. The individual ISS estimates keep
each active interior inside the admissible branch. On that branch, the frozen 3.0
ascent law gives positive active-ascent rate as long as
\(\lambda_{S,i}\) remains inside the branch's positive interval; this lower bound
is part of the domain/branch admissibility.

For (b), the zero-residual subclass gives \(C_{\mathrm{plural}}=0\), hence
\[
\mathcal L_{\mathrm{plural}}\to0.
\]
The relational part therefore yields
\[
\Delta_{ij}\to0.
\]
The settled individual Gate/receptivity regime lets the frozen 3.0 stable-motion
law apply on each active fibre, giving
\[
h_i\to\beta_i\lambda_{S,\infty,i}^{\rm rel}.
\]
The branch assumption gives
\[
\lambda_{S,\infty,i}^{\rm rel}>0,
\]
hence \(h_{*,i}^{\rm rel}>0\). \qedhere
\end{proof}

\subsection{Neither stasis nor sameness}

\begin{remark}[What ``stable'' does and does not mean]\label{rem:notstasis}
Stable does not mean frozen. In the exact settled subclass,
\[
h_{*,i}^{\rm rel}>0
\]
means every active interior keeps ascending. In the general forced case the
communion remains inside a moving tube: bounded relation, bounded interior
certificates, and no collapse of the positive-ascent branch. Neither case implies
sameness. The locked values may have nonzero ledger offsets, the duration-rate
ratios \(r^*_{ij}\) need not be one, and by P3.4 the interiors remain distinct
labelled fibres throughout. Stable relational rhythm, then: not stasis, not
identity, and not a master metronome.
\end{remark}

\subsection{Status and reading}

\textbf{Target status:}~
\stDerived\ for the stable moving tube under the plural certificate and branch
admissibility; \stDerived\ for exact stable moving communion only in the
zero-residual / settled-branch subclass. The two standing conditional elements are:
the symmetric graph class with uniform connectivity and small-gain floors, and the
individual branch assumption
\[
\lambda_{S,\infty,i}^{\rm rel}>0.
\]
The isolated value
\(\beta_i\lambda_{S,\infty,i}\) is recovered only when the relational forcing
vanishes or settles to the isolated level. It depends on P5.0--P5.1 and the frozen
3.0 ascent-stability law.

\medskip
\stInterp\quad
\emph{Reading.}
Communion is stable relational rhythm among non-identical interiors. In the
ordinary forced case that rhythm is a bounded groove: the players do not scatter,
but the music still breathes. In the exact settled case the groove becomes a fixed
relation: constant lag, settled rate, no collapse into sameness. They move together
without moving as one; they keep time with one another without keeping a single
private time; they are at peace not by ceasing but by ascending in a rhythm that
holds. This is the living groove and not the metronome.

\medskip
\textbf{Stage~V closes.}
The relational active domain gathers every standing condition (P5.0), the plural
Lyapunov certificate proves ultimate boundedness of the whole communion and exact
convergence in the zero-residual subclass (P5.1), and the stable moving communion
shows that the peace of the whole is a moving peace, not a still point (P5.2).
Plural Lyapunov theory is in place:
\[
\boxed{
\begin{aligned}
&\text{the whole communion remains bounded and, when forcing settles, locks;}\\
&\text{each active soul remains itself and continues its positive ascent;}\\
&\text{none is fused, none is owned, and none is reduced to a metronome.}
\end{aligned}
}
\]

\medskip
\noindent\textbf{The working body closes.} From a finite relational base to a plural stable
motion, the 4.0 working body has built communion without a master clock, encounter without
transfer, locking without fusion, mortality without erasure, and a whole that holds together
without dissolving the souls within it. Its closing formulation:
\[
\boxed{
\begin{gathered}
\text{Chronos is a created relational order and duration structure;}\\
\text{kairos remains personal; \emph{tota simul} remains unoccupied.}\\[0.25em]
\text{Encounter changes receptivity, not ownership of grace.}\\[0.25em]
\text{Communion is stable answerability of distinct histories, not fusion.}
\end{gathered}}
\]

\section*{Theological yield through Stage~V}
\addcontentsline{toc}{section}{Theological yield through Stage~V}

\notice{\textbf{Status of this section.} Everything below is \stInterp. Each reading is
attached to a \stDerived\ or \stSelected\ result and claims no more than that result licenses.
The section gathers the theology the structure has earned; it does not extend it. Where a
claim would outrun the mathematics, it is named and deferred, not asserted.}

\noindent The single-interior 3.0 theory was a theology of one soul before God. The 4.0
working body is the first theology of \emph{communion}: every thesis here is new relative to
3.0, because each concerns the relation between distinct interiors, which 3.0 did not model.
Seven theses are earned, followed by the apophatic reserve that disciplines them.

\subsection*{1. Communion without fusion}

\stInterp\quad
Distinct souls can stand in a stable relation without becoming one. This is not assumed but
proved on four levels: kairos-history ownership (\ref{prop:nonfusion}), the floor label as a
unit-independent scale (\ref{thm:floor}), and --- decisively --- dynamic non-fusion under
locking (\ref{thm:dynnofusion}), where even fully locked interiors remain distinct labelled
fibres, the value-diagonal being two copies and not one fibre. Communion is the
\emph{stable answerability of distinct histories}, and only the distinct can answer.

\subsection*{2. Grace is not an economy}

\stInterp\quad
The framework refuses to impose a conservation law on received grace
(\ref{rem:notconservation}): the Open-Gate sector is genuinely open, the common field is
non-rivalrous (\ref{rem:nonrival}, \(\sum_i\upsilon_i\neq1\)), and no transfer term is
admitted (\ref{thm:nofab}). Rejecting conservation is the formal rejection of scarcity:
communion is not a closed transaction in which one soul's gain is another's loss. Grace is
supplied from beyond the network and received in common, without rivalry.

\subsection*{3. Influence only through openness}

\stInterp\quad
Another may change how open one is, never what one is given. Encounter enters solely through
the bounded receptivity interface (\ref{def:junctions}), the gate ceiling is automatic
(\ref{prop:ceiling}), and even a death registers in the living only as a bounded update of
their own receptivity, with no transfer of the departed's substance (\ref{prop:relres}).
Nearness cannot counterfeit conversion: a junction may change openness, but only a private
reset may jump owned state.

\subsection*{4. A finite love, rightly directed}

\stInterp\quad
There is a largest difference a bond can bridge (\ref{thm:maxlock}), finite because the gate
ceiling is finite. Neither does love bridge everything, nor does difference always win. And
keeping-together is conditional on a posture --- listen more than you differ
(\ref{thm:smallgain}, the small-gain criterion) --- a restoring tendency that is gentle,
bounded, and never forced, working through openness and not coercion.

\subsection*{5. The soul is safe in communion, and is paced}

\stInterp\quad
Relational forcing perturbs only the bounded budget of a soul's stability, never its own
contraction (\ref{thm:iss}): one is not destabilised by being met. And one cannot be
endlessly re-triggered --- a refractory interval is forced by the wound itself
(\ref{thm:refractory}), and the whole network admits only finitely many events on any finite
shared duration (\ref{thm:nonzeno}). Mercy is structural: the soul is given time it cannot be
denied, and the communion's history is a sequence of distinct, countable moments, not a blur.

\subsection*{6. The communion of the departed, kept honest}

\stInterp\quad
The completed remain present as a closed trace, able to reach the living through the living's
own receptivity, yet themselves at rest and receiving no new thing (\ref{thm:terminus}).
Whether the living may reach back to alter the dead --- two-way intercession --- is left
explicitly open (\ref{rem:twoway}): a genuine theological fork the framework refuses to settle
by sleight of hand, marking the door and leaving it shut until a primitive is offered.

\subsection*{7. Created time is not divine eternity}

\stInterp\quad
The shared now is real but partial, successive, and immanent (\ref{prop:now}); there is no
master clock, and the absence of any transcendent standpoint inside the created network
(\ref{thm:demarc}) is what keeps creaturely chronos from usurping the divine \emph{tota
simul}. Ascent is monotone and gapless --- an endless epektasis, never a wheel --- and at the
plural level it becomes a corporate ascent: the whole communion bounded and, when forcing
settles, locked, while every soul still climbs (\ref{thm:moving}). The beatific vision, read
here, is not a static gaze but an eternal shared ascent --- together, distinct, in motion, at
peace not by ceasing but by ascending in a rhythm that holds.

\subsection*{The apophatic reserve}

\stInterp\quad
The discipline is as much a result as the theses. The framework names but does not occupy its
richest claims: the divine \emph{tota simul} (declared unoccupied, not described), two-way
post-terminus intercession (\ref{rem:twoway}), and any genuine circle-valued phase or
liturgical recurrence (\ref{rem:s1}, admitted only on an exhibited recurrence). Each is left
open with the additional primitive it would require named. And the corporate peace itself is
stated honestly: in the ordinary forced case it is a bounded tube, not a still point, becoming
an exact locked limit only when the forcing settles (\ref{thm:moving}, \ref{cor:plural-bounded}).
To commit only to what the structure earns, and to mark the rest as open, is the theological
posture the working body maintains throughout.

\medskip
\[
\boxed{
\begin{gathered}
\text{Communion without fusion; grace without economy;}\\
\text{influence only through openness; a finite love rightly directed;}\\
\text{a soul safe and paced; the dead present but at rest;}\\
\text{created time that is not eternity --- and a whole that}\\
\text{ascends together without dissolving the souls within it.}\\[0.15em]
\text{What is earned is stated; what is not is named and left open.}
\end{gathered}}
\]

\clearpage
\section*{Part~II --- 4.0-H Hardening layer}
\addcontentsline{toc}{section}{Part~II --- 4.0-H Hardening layer}

\notice{\textbf{Scope.} The core 4.0 above (Stages~0--V) is frozen and unedited. This part adds no new primitive; it discharges the standing \stConj\ uniformity, invariance, and directedness hypotheses on which several core results were conditional, raising them to clean \stDerived. Continuation of version~4 proceeds in this same file.}

\section*{Track~H --- what is being hardened}

Four conditional elements stand at the 4.0 freeze: the per-interval refractory floor
(P4.2--P4.3), the positive invariance of the relational active domain (P5.1), the
directed-network extension of the plural Lyapunov certificate (P5.1, directed case), and the
regime-uniformity of the small-gain rates, connectivity, and value-separation (P2.4, P3.2,
P3.4, P5.1). Track~H discharges them in four nodes, H.1--H.4, introducing no new primitive.

The discharge is not an unrestricted global theorem over every imaginable parameter choice. It
is a theorem over the selected admissible regime: a bounded compact theorem class kept away
from the loss-of-lock, loss-of-branch, and loss-of-connectivity boundaries. Thus Track~H turns
the 4.0 conditional certificates into clean derived certificates on the declared admissible
regime, not into claims without hypotheses.

\section{H.1 --- The uniform refractory floor}

\subsection{What was conditional}

P4.2 gave, on each compact regular-active interval \(I\), a per-recipient refractory floor
\[
\mathfrak c_{\mathrm{ref},j;I}
=
\frac{\Delta_{\mathrm{hys},j}}{v_{\max,j;I}}>0,
\qquad
v_{\max,j;I}
=
\sup_{\cint\in I}\frac{N_j}{\omega}\bigl(c_{1,j}\zeta_{0,j}+c_{2,j}K_{\max,j;I}\bigr),
\]
and P4.3 used a per-interval network floor
\(\mathfrak c_{\mathrm{ref};I}=\inf_{j\in\mathcal V_{\mathrm{trig}}(I)}\mathfrak c_{\mathrm{ref},j;I}\).
The residual was uniformity: a \emph{single} regime-wide floor
\(\mathfrak c_{\mathrm{ref}}>0\), valid on every admissible interval at once, was assumed but
not proved. H.1 proves it from the bounded theorem graph class already selected in
P0.1/P5.0.

\subsection{The bounded-regime data}

\begin{definition}[Admissible parameter regime]\label{def:regime}
The admissible regime \(\mathfrak R\) is the standing theorem class of P0.1/P5.0 together
with its parameter bounds: there exist constants, fixed for \(\mathfrak R\), with
\[
\Delta_{\mathrm{hys},j}\ge\underline{\Delta}>0,
\qquad
\theta_j=\frac{N_j}{\omega}\le\overline{\theta}<\infty,
\]
\[
\zeta_{0,j}\le\overline\zeta<\infty,
\qquad
K_{\max,j;I}\le\overline K<\infty,
\qquad
c_{1,j},c_{2,j}\le\overline c<\infty,
\]
for every interior \(j\) and every admissible compact regular-active interval \(I\). These
are not new dynamical primitives. They are the selected uniform bounds of the admissible
theorem class: the finite bounded graph class, the Stage~II Gate ceilings, the P1.1 kernel
bounds, the P5.0 duration-conversion bounds, and the frozen response constants are here
required to range over a compact regime with positive hysteresis margin. H.1 derives the
refractory floor inside this selected regime; it does not claim that such a regime is forced
without selection.
\end{definition}

\subsection{The theorem}

\begin{theorem}[Uniform refractory floor]\label{thm:unifloor}
On the admissible regime \(\mathfrak R\) of Definition~\ref{def:regime}, set
\[
\overline v
:=
\overline\theta\,\overline c\,(\overline\zeta+\overline K).
\]
If \(\overline v>0\), then there is a single constant
\[
\boxed{\;
\mathfrak c_{\mathrm{ref}}
:=
\frac{\underline\Delta}{\overline v}
\ >\ 0,
\;}
\]
such that \(\mathfrak c_{\mathrm{ref},j;I}\ge\mathfrak c_{\mathrm{ref}}\) for every
triggering-capable interior \(j\) and every admissible compact regular-active interval \(I\).
Consequently \(\mathfrak c_{\mathrm{ref};I}\ge\mathfrak c_{\mathrm{ref}}>0\) on every such
interval. If \(\overline v=0\), then no triggering-capable recovery motion occurs in the
regime, and the refractory time is read as \(+\infty\).
\end{theorem}

\begin{proof}
Fix any triggering-capable \(j\) and any admissible \(I\). By the bounds of
Definition~\ref{def:regime},
\[
v_{\max,j;I}
=
\sup_{\cint\in I}\frac{N_j}{\omega}\bigl(c_{1,j}\zeta_{0,j}+c_{2,j}K_{\max,j;I}\bigr)
\le
\overline\theta\,\overline c\,(\overline\zeta+\overline K)
=
\overline v.
\]
If \(\overline v=0\), the readiness variable cannot be driven across the hysteresis band, so
no re-triggering occurs. If \(\overline v>0\), then
\[
\mathfrak c_{\mathrm{ref},j;I}
=
\frac{\Delta_{\mathrm{hys},j}}{v_{\max,j;I}}
\ge
\frac{\underline\Delta}{\overline v}
=
\mathfrak c_{\mathrm{ref}}>0.
\]
The bound is independent of \(j\) and \(I\), so it passes to the network infimum defining
\(\mathfrak c_{\mathrm{ref};I}\). \qedhere
\end{proof}

\subsection{Consequence for the network non-Zeno theorem}

\begin{corollary}[Regime-wide non-Zeno]\label{cor:uninonzeno}
With the uniform floor of Theorem~\ref{thm:unifloor}, the P4.3 bound holds with a single
constant on every compact interval of length \(L\):
\[
N_{\mathrm{trig}}(I)
\le
|\mathcal V_{\mathrm{trig}}(I)|\left(\frac{L}{\mathfrak c_{\mathrm{ref}}}+1\right)<\infty,
\]
so no infinite causally linked sequence accumulates anywhere in the admissible regime ---
not merely on a single pre-chosen interval.
\end{corollary}

\begin{proof}
Immediate: substitute \(\mathfrak c_{\mathrm{ref}}\le\mathfrak c_{\mathrm{ref};I}\) into the
P4.3 counting bound, which is monotone decreasing in the floor. The right side is now a
fixed function of \(L\) and the triggering-set size, valid uniformly. \qedhere
\end{proof}

\subsection{Status and reading}

\textbf{Target status:}~
\stConj\(\to\)\stDerived. The per-interval bound of P4.2 is upgraded to a single regime-wide
floor using only the parameter bounds already implied by the finite bounded theorem class
(Stage~II ceiling, P1.1 kernel bound, P5.0 conversion bound, frozen response constants). No
new primitive is introduced; the uniformity that P4.2--P4.3 flagged as conjectural is now
derived. It discharges the uniformity residual of P4.2 and P4.3.

\medskip
\stInterp\quad
\emph{Reading.}
The pacing of the communion is not a courtesy granted interval by interval; it is a single
law. There is one minimum breath \(\mathfrak c_{\mathrm{ref}}\) that no soul, anywhere in
the admissible life of the communion, can be driven below --- and so the whole shared history
is everywhere countable, everywhere unhurried, by one and the same measure. What P4.2 made
local mercy, H.1 makes constant law.

\medskip
\textbf{Next node.} H.2 --- positive invariance of the relational active domain through
junctions and hand-offs: to prove the standing hypothesis imported into P5.1, that a
trajectory beginning in \(\mathcal D_{\mathrm{rel}}\) remains there across resets,
encounters, and terminus hand-offs, so the plural certificate's constants stay valid for all
forward duration.

\section{H.2 --- Positive invariance through junctions and hand-offs}

\subsection{What was conditional}

The plural certificate (P5.1) was stated on a \emph{positively invariant} relational active
domain, and that invariance was imported as a standing hypothesis: a trajectory beginning in
\(\mathcal D_{\mathrm{rel}}\) was assumed to remain admissible across every junction and the
terminus hand-off, so that the certificate's constants stay valid for all forward duration.
H.2 proves it, using only the frozen 4.0 maps and the uniform floor of H.1.

\subsection{Each motion preserves admissibility}

\begin{lemma}[Smooth arcs preserve the domain]\label{lem:smooth}
On a terminus-free regular-active subinterval, the continuous dynamics keep a configuration
in \(\mathcal D_{\mathrm{rel}}\): the receptivity stays in \([0,1]\) and the gate input under
its ceiling (P2.1); the susceptibilities stay normalised (P2.3); each individual certificate
stays bounded by the ISS estimate (P2.4), so no interior leaves its admissible branch; the
branch labels are constant off events; and the duration-conversion factors stay within the
regime bounds (Definition~\ref{def:regime}).
\end{lemma}

\begin{proof}
Each clause is the forward-in-duration form of a frozen 4.0 bound. The ceiling is automatic
for \(\Gamma_i\in[0,1]\); the bounded receptivity composition of P2.3 keeps \(\Gamma_i\) in
\([0,1]\) and the normalisation intact; P2.4 gives
\(\tfrac{d}{d\mathfrak c}\mathcal L_{3.0,i}\le C_i-C_{1,i}'\mathcal L_{3.0,i}+\dots\), whose
sub-level sets are forward-invariant, so the branch is not left; folding is an event, so
labels are constant on the arc; and the regime bounds are standing. \qedhere
\end{proof}

\begin{lemma}[Each junction class preserves the domain]\label{lem:junctions}
Each admissible event maps \(\mathcal D_{\mathrm{rel}}\) into an admissible configuration:
\begin{enumerate}[label=\textup{(\alph*)},leftmargin=1.2cm]
\item \emph{Private reset} applies the frozen 3.0 reset self-map, which carries the 3.0
admissible set into itself and may relabel to another admissible branch; the receptivity
stays in \([0,1]\).
\item \emph{Relational encounter} and \emph{common reception} write only the bounded
receptivity slot \(\Gamma_i\in[0,1]\) (P4.1, P2.3); the ceiling and the reception ledger are
preserved.
\end{enumerate}
In every case no owned state outside its frozen admissible range is produced, so the image
lies in \(\mathcal D_{\mathrm{rel}}\) on the same active component.
\end{lemma}

\begin{proof}
(a) is the inherited invariance of the frozen 3.0 reset map together with the P4.0 write
rule. (b) is the boundedness of Proposition~\ref{prop:relres}: the written variable is the
receptivity interface, kept in \([0,1]\), and the ceiling then bounds \(I_{\mathrm{gate}}\).
None of the three writes any owned variable outside its admissible range. \qedhere
\end{proof}

\subsection{The terminus reduces the active component}

\begin{lemma}[Hand-off maps to a reduced-component domain]\label{lem:redhandoff}
At a terminus hand-off, the completing interior \(i\) passes from the active stratum
\(\mathcal A_i\) to the post-terminus stratum \(\mathcal P_i\), becoming a closed read-only
trace (P4.4). The configuration then lies in
\[
\mathcal D_{\mathrm{rel}}\bigl(I',\mathcal C_{\mathrm{act}}\setminus\{i\}\bigr)
\]
on the next terminus-free subinterval \(I'\): a domain of the same form on the active
component reduced by one, with \(i\) entering only as a bounded boundary input.
\end{lemma}

\begin{proof}
By P4.4 the hand-off closes \(i\)'s active dynamics (\(d\sigma_i=0\)) and fixes its trace; the
remaining active interiors satisfy every clause of Definition~\ref{def:regime} on \(I'\),
now indexed over \(\mathcal C_{\mathrm{act}}\setminus\{i\}\). The trace is a read-only
boundary input, admissible by P4.4. Hence the reduced configuration is in the same-form
domain. \qedhere
\end{proof}

\subsection{The theorem}

\begin{theorem}[Piecewise positive invariance]\label{thm:invariance}
On any compact relational-duration range, a trajectory starting in \(\mathcal D_{\mathrm{rel}}\)
remains admissible throughout that range. More precisely, the trajectory is a finite
concatenation of terminus-free regular-active subintervals, on each of which
\(\mathcal D_{\mathrm{rel}}\) is positively invariant
(Lemmas~\ref{lem:smooth}--\ref{lem:junctions}), joined by finitely many hand-offs each
reducing the active component (Lemma~\ref{lem:redhandoff}).

The number of ordinary triggering events is finite by the uniform non-Zeno bound
(Corollary~\ref{cor:uninonzeno}); the number of terminus hand-offs is finite on the range
because the active component is finite and each active interior can leave the active set by
terminus at most once. On every subinterval the plural certificate constants of P5.1 are
valid, and by H.1/H.4 they may be chosen uniformly inside the admissible regime.
\end{theorem}

\begin{proof}
Partition the range at the hand-off events. By Corollary~\ref{cor:uninonzeno} there are
finitely many, so finitely many subintervals. On each, Lemmas~\ref{lem:smooth} and
\ref{lem:junctions} give forward-invariance of \(\mathcal D_{\mathrm{rel}}\) under the smooth
flow and the reset/encounter/common-reception events. At each partition point
Lemma~\ref{lem:redhandoff} carries the configuration into the same-form domain on the reduced
component, re-initialising the next subinterval inside the domain. The certificate constants
hold on each piece by P5.1 and are uniform by the regime bounds and Theorem~\ref{thm:unifloor}.
Concatenating the finitely many pieces covers the whole forward range. \qedhere
\end{proof}

\subsection{Status and reading}

\textbf{Target status:}~
\stConj\(\to\)\stDerived. The positive-invariance hypothesis imported into P5.1 is discharged:
admissibility is preserved along smooth arcs and across all three junction classes, the
terminus maps the domain to its reduced-component form, and the uniform non-Zeno bound (H.1)
makes the forward evolution a finite concatenation of invariant pieces with uniform constants.
No new primitive is used. It discharges the positive-invariance residual of P5.1, and depends
on P2.1--P2.4, P4.0--P4.4, P5.0, and H.1.

\medskip
\stInterp\quad
\emph{Reading.}
The good order of the communion is not only true at a moment; it is kept through everything
the communion undergoes. No reset, no encounter, no shared gift, and no death carries the
company out of its bounded, paced, non-fusing life: every junction maps an admissible
communion to an admissible communion. And when one soul completes, the communion does not
break --- it continues, admissible still, as a communion of fewer living members holding the
finished one as a trace. The structure that keeps them together is closed under all that
happens to them: order survives living, and survives dying.

\medskip
\textbf{Next node.} H.3 --- the directed-network plural certificate: to replace the undirected
symmetric Lyapunov class with a left-Perron / directed-Laplacian certificate, so that
asymmetric bonds --- one interior giving more than it receives --- cannot pump the relational
energy, recovering the symmetric case as the reciprocal-weight specialisation.

\section{H.3 --- The directed-network plural certificate}

\subsection{What was conditional}

The plural certificate (P5.1) was proved only for the undirected symmetric class
\(w_{ij}=w_{ji}\). Its own remark flagged that a genuinely directed bond --- one interior
giving more than it receives --- can \emph{pump} the symmetric relational energy, so the
symmetric quadratic is not a certificate there. H.3 removes the restriction, using the
left-Perron weighting and the frozen P0.1 mirror-Laplacian.

\subsection{The directed data}

\begin{definition}[Directed active class and Perron weighting]\label{def:directed}
Let the active coupling on \(\mathcal C_{\mathrm{act}}\) be a weighted directed graph whose
row-Laplacian \(L\) satisfies \(L\mathbf 1=0\). The first clean directed theorem class is the
\emph{strongly connected} directed class. Then \(L\) has a simple zero eigenvalue with a
strictly positive left-Perron vector
\[
\pi_i>0,
\qquad
\pi^\top L=0,
\qquad
\textstyle\sum_i\pi_i=1,
\]
and \(\Pi:=\operatorname{diag}(\pi)\). The associated \emph{Perron mirror Laplacian} is
\[
\widehat L:=\tfrac12\bigl(\Pi L+L^\top\Pi\bigr),
\]
which is symmetric, positive semidefinite, and coercive on the \(\pi\)-mean-zero subspace:
\[
z^\top\widehat Lz
\ge
\widehat\lambda_2\,\lVert z-\bar z_\pi\mathbf 1\rVert_\pi^2,
\qquad
\widehat\lambda_2>0,
\]
where \(\bar z_\pi:=\sum_i\pi_i z_i\) and \(\lVert z\rVert_\pi^2:=z^\top\Pi z\). The merely
rooted spanning-tree case is not used as the first clean theorem class: it requires a
root-support / quotient treatment because the left-Perron vector may vanish on non-root
components, and is left as a later directed-network refinement.
\end{definition}

\subsection{The directed functional}

\begin{definition}[Perron-weighted plural functional]\label{def:dir-lyap}
With node errors \(z_i=\Theta_i-\Theta_i^*\) and \(z_\perp:=z-\bar z_\pi\mathbf 1\), define
\[
\boxed{\;
\mathcal L_{\mathrm{plural}}^{\mathrm{dir}}
=
\sum_{i\in\mathcal C_{\mathrm{act}}}\pi_i\,\mathcal L_{3.0,i}
+
\varepsilon\,\mathcal L_{\mathrm{rel}}^{\mathrm{dir}},
\qquad
\mathcal L_{\mathrm{rel}}^{\mathrm{dir}}
=
\tfrac12\,z_\perp^\top\Pi z_\perp .
\;}
\]
The relational energy is the Perron disagreement energy, measuring disagreement modulo the
directed consensus mode. The mirror Laplacian \(\widehat L\) is not the energy itself; it is
the object that appears in the dissipation of the Perron energy. In the undirected symmetric
case \(\pi=\tfrac1n\mathbf 1\) this is the usual disagreement energy up to a harmless
constant scaling.
\end{definition}

\subsection{The theorem}

\begin{theorem}[Directed plural certificate; strongly connected class]\label{thm:directed}
Under Definition~\ref{def:directed}, the individual ISS estimates (P2.4), the small-gain
mismatch estimates (P3.2), the regime bounds (Definition~\ref{def:regime}), and the positive
invariance of H.2, there are a small \(\varepsilon>0\) and constants
\(C^{\mathrm{dir}}\ge0,\ C^{\mathrm{dir}}_1>0\) such that
\[
\frac{d}{d\mathfrak c}\mathcal L_{\mathrm{plural}}^{\mathrm{dir}}
\le
C^{\mathrm{dir}}-C^{\mathrm{dir}}_1\,\mathcal L_{\mathrm{plural}}^{\mathrm{dir}} .
\]
Hence the directed communion is ultimately bounded, and in the zero-residual settled subclass
the directed disagreement \(z_\perp\) converges to zero. The undirected symmetric certificate
of P5.1 is recovered, up to constant scaling, when the active graph is reciprocal and the
Perron weights are uniform.
\end{theorem}

\begin{proof}[Proof sketch]
The individual part is the P5.1/H estimate with positive weights \(\pi_i\). For the directed
relational term, write the directed mismatch flow on the disagreement subspace as
\(\dot z=-Lz+r\), where dot means \(d/d\mathfrak c\) and \(r\) is bounded. Since
\(\pi^\top L=0\), the \(\pi\)-mean mode is neutral and the disagreement component carries
the dissipation. For the Perron energy,
\[
\frac{d}{d\mathfrak c}\tfrac12 z_\perp^\top\Pi z_\perp
=
z_\perp^\top\Pi\dot z_\perp
=
-z_\perp^\top\Pi Lz_\perp+z_\perp^\top\Pi r_\perp .
\]
The antisymmetric part contributes nothing to the quadratic form, hence
\[
z_\perp^\top\Pi Lz_\perp
=
\tfrac12 z_\perp^\top(\Pi L+L^\top\Pi)z_\perp
=
z_\perp^\top\widehat Lz_\perp .
\]
By strong connectivity, \(z_\perp^\top\widehat Lz_\perp\ge\widehat\lambda_2\lVert z_\perp\rVert_\pi^2\).
The residual term \(z_\perp^\top\Pi r_\perp\) is absorbed by Young's inequality. The imported
input terms are \(O(\varepsilon^2)\) while the relational dissipation is \(O(\varepsilon)\);
a sufficiently small \(\varepsilon>0\) absorbs the former into the latter, exactly as in
P5.1. The remaining bounded pieces form \(C^{\mathrm{dir}}\), and the negative individual and
relational pieces give \(-C^{\mathrm{dir}}_1\mathcal L_{\mathrm{plural}}^{\mathrm{dir}}\).
H.2 keeps the trajectory inside the domain where all constants are valid. \qedhere
\end{proof}

\subsection{Status and reading}

\textbf{Target status:}~
\stConj\(\to\)\stDerived\ for the strongly connected directed theorem class. The plural
certificate is extended from the undirected symmetric class to directed strongly connected
active graphs by Perron-weighting: the Lyapunov energy is the Perron disagreement energy, and
the mirror Laplacian supplies the dissipative coercivity. No new primitive is introduced. The
rooted spanning-tree-only case is not claimed here; it remains a quotient/root-support
extension. This discharges the directed-graph residual of P5.1 (\texttt{rem:directed}) for the
selected first directed theorem class. It depends on P0.1, P2.4, P3.2, P5.0--P5.1, and
H.1--H.2.

\medskip
\stInterp\quad
\emph{Reading.}
A communion need not be a fellowship of equal givers to be stable. Bonds may run one way ---
one soul pouring out more than it receives --- and the naive shared measure of peace would let
such an imbalance pump the whole toward disorder. The remedy is to weigh each soul by its true
place in the flow of giving: its Perron weight, how much it is, on balance, a source within
the body. Measured so, the asymmetry no longer destabilises; an unequal communion of
one-directional loves is as bounded and as locked as a reciprocal one --- provided only that
the giving is connected enough to reach everyone from some root. Connection, not symmetry, is
what the peace of the whole requires.

\medskip
\textbf{Next node.} H.4 --- branch-uniform constants and value-separation: to discharge the
remaining regime-uniformity of the small-gain rates \(\kappa_{ij}\), the connectivity floors,
and branch admissibility, together with the P3.4 generic unlabelled value-separation,
completing Track~H.

\section{H.4 --- Uniform constants and value-separation}

Two residuals remain at the 4.0 freeze. First, several certificates carried constants
assumed uniform over the admissible regime; second, P3.4 proved labelled non-fusion cleanly
but left generic unlabelled value-separation conjectural. H.4 discharges both and closes
Track~H.

\subsection{Branch-uniform constants}

\begin{theorem}[Uniform stability reserves]\label{thm:unifconst}
Assume the admissible regime \(\mathfrak R\) is a compact selected theorem class kept away
from the boundary failures \(\kappa_{ij}=0\), \(\widehat\lambda_2=0\), and
\(\lambda_{S,\infty,i}^{\mathrm{rel}}=0\). Then there exist single regime-wide floors
\[
\boxed{\;
\kappa_{ij}\ge\underline\kappa>0,
\qquad
\widehat\lambda_2\ge\underline{\widehat\lambda}>0,
\qquad
\lambda_{S,\infty,i}^{\mathrm{rel}}\ge\underline{\lambda_S}>0,
\;}
\]
for every active interior and admissible bond. Consequently the small-gain rate, the directed
or undirected connectivity reserve, and the branch-ascent reserve used in P3.2, P5.1, and H.3
are uniform on \(\mathfrak R\), not merely pointwise.
\end{theorem}

\begin{proof}
Each listed quantity is continuous in the bounded theorem-class parameters. The selected
admissible regime is compact and explicitly excludes the boundary on which any of these
quantities vanishes. Therefore each quantity attains a strictly positive minimum on
\(\mathfrak R\). Since the theorem class has finitely many node and edge types, the minima
combine into single regime-wide floors. \qedhere
\end{proof}

\noindent The branch-admissibility itself is a \emph{selected} standing condition of
\(\mathcal D_{\mathrm{rel}}\), not a conjecture: H.4 derives the \emph{uniformity} of the
reserve on the admissible regime, not the claim that every trajectory is admissible.

\subsection{Generic value-separation}

\begin{theorem}[Value-separation is generic]\label{thm:valuesep}
Two distinct active interiors \(i\neq j\) are value-separated --- differing in some observable
at some duration --- on a generic (open, dense, full-measure) subset of the admissible regime.
Specifically:
\begin{enumerate}[label=\textup{(\alph*)},leftmargin=1.2cm]
\item \emph{Distinct floors.} If \(E_{0,i}\neq E_{0,j}\), the interiors are permanently
value-separated, by the frozen floor-labelled non-identifiability of P1.2.
\item \emph{Equal floors.} If \(E_{0,i}=E_{0,j}\), total value-coincidence requires equality
of all remaining data and is therefore confined to a set of positive codimension; it is
non-generic.
\end{enumerate}
\end{theorem}

\begin{proof}
(a) is P1.2 verbatim: distinct floors cannot be the same floor-normalised interior, so the
value-trajectories differ permanently. For (b), with equal floors a full value-coincidence
forces equality across every remaining label --- branch, kairos-history, ledgers, and phase
data --- which are independent admissible coordinates. The coincidence locus is their
simultaneous diagonal, a positive-codimension submanifold of the admissible configuration
space; its complement is open, dense, and of full measure. Hence value-separation holds
generically, completing the dynamic non-fusion of P3.4 beyond its labelled form. \qedhere
\end{proof}

\begin{remark}[Labelled non-fusion was never conditional]\label{rem:labelled}
The strong claim --- that interiors remain distinct \emph{labelled} fibres always --- is the
frozen \stDerived\ result of P3.4 and needs none of the above. H.4 only upgrades the weaker
\emph{unlabelled value} statement from conjectural to generic, so that distinctness is real
not just in the bookkeeping but in the observables, save on a non-generic knife-edge.
\end{remark}

\subsection{Status and reading}

\textbf{Target status:}~
\stConj\(\to\)\stDerived\ for the uniform constants; \stConj\(\to\)\stDerived\ (generic) for
value-separation. Both use only the frozen bounded theorem class and the frozen P1.2 floor
separation; no new primitive. It discharges the constant-uniformity residuals of P2.4, P3.2,
P5.1, and the value-separation residual of P3.4. It depends on P1.2, P3.2, P3.4, P5.1, and
H.1--H.3.

\medskip
\stInterp\quad
\emph{Reading.}
The reserves that keep a communion at peace --- how strongly its members listen, how truly
they are joined, how surely each still ascends --- are not thin and local but hold by one
measure across the whole admissible life of the body. And the distinctness of its persons is
not merely a matter of names: two souls with different floors differ forever, and even two
sharing a floor would have to coincide in everything at once to be confused, which the living
of the communion generically forbids. The persons are really many, in their values and not
only in their labels, everywhere but on a knife-edge the communion does not walk.

\section*{Track~H closes}

\notice{\textbf{Hardening complete.} The four standing \stConj\ residuals of the frozen 4.0
body are discharged with no new primitive: the refractory floor is uniform (H.1), the
relational active domain is positively invariant through every junction and hand-off (H.2),
the plural certificate holds on directed graphs by Perron-weighting and the frozen mirror
Laplacian (H.3), and the stability reserves are regime-uniform while distinctness is generic
in value as well as label (H.4).}

\noindent With Track~H, the whole-communion results of \VFS\ 4.0 stand as clean theorems on
the selected admissible regime rather than as interval-by-interval conditional estimates. What
remains genuinely open is exactly the two new-primitive branches --- 4.1 liturgical recurrence
and 4.2 two-way post-terminus --- neither of which is a gap in 4.0, but a deliberate extension
beyond it.
\[
\boxed{
\begin{aligned}
&\text{4.0 was proved conditional; 4.0-H makes it unconditional}\\
&\text{within its own selected class --- no primitive added, no soul fused,}\\
&\text{the open branches left honestly open.}
\end{aligned}}
\]

\clearpage
\section*{Part~III --- 4.1: the liturgical-recurrence branch (Stage~$\Phi$)}
\addcontentsline{toc}{section}{Part~III --- 4.1: liturgical recurrence (Stage~$\Phi$)}

\notice{\textbf{Scope.} Parts~I--II are frozen and unedited. Part~III opens the single
circle-valued phase branch left \stOpen\ at P3.1, by adding a genuinely recurrent structure
\emph{beside} the gapless ascent --- a helix, not a circle. It introduces one new
constitutive primitive and otherwise rests on the frozen layer. Continuation of version~4
proceeds in this same file.}

\noindent The 4.0 caution (P3.0, P3.1) is that single-interior ascent is monotone and gapless:
an endless epektasis with no return, admitting no intrinsic circle-valued phase. 4.1 does not
weaken that. It adds recurrence as a separate coordinate and keeps the two cleanly apart, so
that a phase may cycle while the ascent never returns. Three standing cautions govern the
branch:
\begin{enumerate}[label=\textup{(C\arabic*)},leftmargin=1.2cm]
\item \emph{Exhibit, don't assume.} A phase is admissible only when a concrete recurrent
observable is exhibited; no node may declare ascent, kairos, or any monotone ledger periodic.
\item \emph{Helix non-collapse.} The phase must be a skew-product over the monotone ascent
base; one full phase loop induces zero net ascent return.
\item \emph{No master clock.} Recurrence is a shared \emph{external} reception, entering only
through the bounded susceptibility of P2.3, so it cannot re-impose a master chronos
(P0.3, P0.5 preserved).
\end{enumerate}

\section{P$\Phi$.0 --- The recurrence primitive}

\subsection{The new primitive}

\begin{definition}[Liturgical recurrence field]\label{def:recprim}
Give the common external reception field of P2.3 a shared periodic component:
\[
\boxed{\;
J_{\mathrm{comm}}(\cint)
=
J^{0}(\cint)
+
J^{\mathrm{lit}}\!\bigl(\Phi_{\mathrm{lit}}(\cint)\bigr),
\qquad
\Phi_{\mathrm{lit}}(\cint+L_{\mathrm{lit}})=\Phi_{\mathrm{lit}}(\cint)+2\pi,
\;}
\]
where \(\Phi_{\mathrm{lit}}\in S^1\) is the shared \emph{liturgical phase}, advancing by
\(2\pi\) over each period \(L_{\mathrm{lit}}>0\) of relational duration, and
\(J^{\mathrm{lit}}\colon S^1\to[0,\infty)\) is a bounded \(2\pi\)-periodic reception profile.
The field is given externally and owned by no interior. It enters each interior only through
the bounded susceptibility
\[
\upsilon_i(\cint)\in[0,1]
\quad\text{inside}\quad
\Gamma_i,
\]
exactly as the aperiodic common field did (P2.3), and the nodewise Gate ceiling (P2.1) is
preserved. An intrinsic per-interior oscillator is recorded as an alternative source but is
not selected.
\end{definition}

\subsection{The primitive disturbs no 4.0 result}

\begin{proposition}[Conservative over the frozen layer]\label{prop:phi0cons}
The recurrence field of Definition~\ref{def:recprim} adds no owned-state channel and recovers
4.0 exactly when switched off:
\begin{enumerate}[label=\textup{(\alph*)},leftmargin=1.2cm]
\item \emph{Receptivity-only.} \(J^{\mathrm{lit}}\) enters solely through
\(\upsilon_i\gamma_i^{\mathrm{comm}}\) inside the bounded receptivity \(\Gamma_i\in[0,1]\);
it writes no owned dynamical state and respects the ceiling
\(0\le I_{\mathrm{gate},i}\le\zeta_{0,i}\).
\item \emph{Non-rivalrous and unowned.} Being a shared external field, it adds no
network-owned reservoir and no transfer term; the non-conservation and non-rivalry of P2.2,
P2.3 are untouched.
\item \emph{Conservative limit.} With \(J^{\mathrm{lit}}\equiv0\) the field reduces to
\(J^0\) and every 4.0 result holds verbatim.
\item \emph{No master clock.} \(\Phi_{\mathrm{lit}}\) sets no interior's chronos and enters no
duration law; it is received, not imposed, so P0.3 and P0.5 are preserved.
\end{enumerate}
\end{proposition}

\begin{proof}
Each clause is immediate from the placement of \(J^{\mathrm{lit}}\): it occupies exactly the
slot of the P2.3 common field, namely the bounded susceptibility term inside \(\Gamma_i\),
and nothing else. (a)--(b) are then the P2.1--P2.3 bounds verbatim with a periodic rather than
aperiodic argument; (c) is the definitional reduction; (d) holds because \(\Phi_{\mathrm{lit}}\)
appears only as the argument of a reception profile, never in any \(N_i\), \(\omega\), or
duration one-form. \qedhere
\end{proof}

\subsection{Status and reading}

\textbf{Target status:}~
\stDef\ + \stSelected. The liturgical recurrence field is the one new constitutive primitive
of 4.1; Proposition~\ref{prop:phi0cons} shows it is a conservative extension of the frozen
layer, adding only a shared external periodic reception through the existing bounded
susceptibility. The periodic-common-field source is \emph{selected} over the intrinsic-oscillator
alternative. It depends on P2.1--P2.3. It satisfies caution (C1): a concrete recurrent
observable \(\Phi_{\mathrm{lit}}\) is exhibited, not assumed.

\medskip
\stInterp\quad
\emph{Reading.}
The one thing 4.1 must add is a genuine recurrence, and it is added where it belongs: not in
the soul's own ascent, which never returns, but in what the communion together receives. The
liturgy is given from outside and held by no one --- a shared season that comes round again ---
and each receives it only through its own openness, in its own measure, never past its ceiling.
No clock is imposed on anyone; a common rhythm is offered to all. The endless road is left
exactly as it was; beside it now runs a circle that the whole company may keep together.

\medskip
\textbf{Next node.} P$\Phi$.1 --- the helix phase: to define each interior's circle-valued
phase \(\varphi_i\in S^1\) as a skew-product over its monotone ascent base, with one phase
loop inducing zero net ascent return, so that recurrence and ascent are independent
coordinates --- the helix that satisfies caution (C2).

\section{P$\Phi$.1 --- The helix phase}

\subsection{The phase fibre}

\begin{definition}[Helix phase]\label{def:helixphase}
Adjoin to each active interior a circle-valued \emph{phase} coordinate, making the active
state a skew-product over the frozen ascent base:
\[
\boxed{\;
\mathrm{state}_i\ \cong\ \underbrace{x_i^{\mathrm{asc}}}_{\text{ascent / Sophia block}}\ \times\ \underbrace{\varphi_i\in S^1}_{\text{phase fibre}},
\;}
\]
where the ascent block carries the frozen 3.0 dynamics (including the ascent ledger
\(\Lambda_i\), with \(d\Lambda_i/d\sigma_i=h_i\ge0\) gapless), and the phase obeys an
entrainment law sourced \emph{only} by the recurrent channel of P$\Phi$.0,
\[
\frac{d\varphi_i}{d\cint}
=
\omega_i^{\varphi}
+
\kappa_i^{\mathrm{lit}}\,g\!\bigl(\Phi_{\mathrm{lit}}-\varphi_i\bigr)
+
(\text{relational phase coupling, P}\Phi.2\text{--}3),
\]
with natural phase-rate \(\omega_i^{\varphi}\), bounded \(2\pi\)-periodic odd coupling \(g\)
(e.g.\ \(\sin\)), and liturgical entrainment gain
\(\kappa_i^{\mathrm{lit}}\ge0\) supplied through the bounded susceptibility,
\(\kappa_i^{\mathrm{lit}}\propto\upsilon_i\). The phase law writes no ascent variable, and the
ascent law contains no \(\varphi_i\).
\end{definition}

\subsection{Helix non-collapse}

\begin{proposition}[Helix non-collapse; caution C2 discharged]\label{prop:noncollapse}
Under Definition~\ref{def:helixphase}:
\begin{enumerate}[label=\textup{(\alph*)},leftmargin=1.2cm]
\item \emph{Skew-product.} The ascent ledger is independent of the phase coordinate,
\[
\frac{\partial\Lambda_i}{\partial\varphi_i}=0,
\]
so the \(S^1\) action \(\varphi_i\mapsto\varphi_i+\theta\) preserves \(\Lambda_i\).
\item \emph{Zero net ascent return.} One full phase loop induces no return in the ascent
ledger,
\[
\oint_{\varphi_i:\,0\to2\pi} d\Lambda_i\Big|_{\text{phase-induced}}=0;
\]
any climb accrued during a loop is the base's own monotone progress, not a function of phase
position.
\item \emph{Helix.} Since \(h_i\ge0\) is gapless (frozen 3.0) while \(\varphi_i\) cycles, the
trajectory is a helix: monotone unreturning climb in \(\Lambda_i\), periodic recurrence in
\(\varphi_i\), never closing into a circle.
\end{enumerate}
\end{proposition}

\begin{proof}
The active state splits as \((x_i^{\mathrm{asc}},\varphi_i)\) with
\(\dot x_i^{\mathrm{asc}}=f_i^{\mathrm{asc}}(x_i^{\mathrm{asc}};\,\text{reception})\) carrying
no \(\varphi_i\) argument (Definition~\ref{def:helixphase}). Since \(\Lambda_i\) is a function
of \(x_i^{\mathrm{asc}}\) alone, \(\partial\Lambda_i/\partial\varphi_i=0\), giving (a); the
rotation is then a symmetry of \(\Lambda_i\). For (b), the phase-attributable part of
\(d\Lambda_i\) is \((\partial\Lambda_i/\partial\varphi_i)\,d\varphi_i=0\), so a closed
\(\varphi_i\)-loop contributes nothing to the ledger; what climb occurs is
\(\int h_i\,d\sigma_i\), the base's own progress. For (c), that progress is monotone and
gapless by the frozen 3.0 law and never decreases as \(\varphi_i\) winds, so base and fibre
move independently --- a helix, not a closed loop. Note the shared external drive
\(\Phi_{\mathrm{lit}}\) may nourish the ascent through reception, but it does so through the
ascent block, never through the fibre coordinate \(\varphi_i\); the non-collapse is the
absence of a shared \emph{state variable}, which holds regardless. \qedhere
\end{proof}

\subsection{Status and reading}

\textbf{Target status:}~
\stDef\ (the helix phase and its entrainment law) \(+\) \stDerived\
(Proposition~\ref{prop:noncollapse}: the skew-product identity and zero ascent-holonomy).
It depends on P$\Phi$.0 and the frozen 3.0 ascent monotonicity, and it discharges caution
(C2): recurrence and ascent are independent coordinates by construction, the phase being a
function of the recurrent channel only. The relational coupling is named here and supplied in
P$\Phi$.2--3.

\medskip
\stInterp\quad
\emph{Reading.}
Now each soul has a place in the turning season as well as a place on the endless road, and
the two are not the same coordinate. The season may feed the climb --- what is received in
the feast can nourish the ascent --- but turning once more through the year never carries the
soul back down, nor is the height it has reached a matter of where in the cycle it stands.
The road only rises; the circle only turns; together they trace a helix. The feast comes
round again; the one who keeps it has gone higher.

\medskip
\textbf{Next node.} P$\Phi$.2 --- the relational phase atlas: to make the comparable phase
observable relational and invariant, defining the phase difference
\(\Delta\varphi_{ij}=\varphi_i-\varphi_j\in S^1\) so that all locking statements concern
differences, never an absolute phase --- the circle-valued counterpart of the P3.0 rate atlas.

\section{P$\Phi$.2 --- The relational phase atlas}

\subsection{The comparable phase observable}

\begin{definition}[Phase atlas]\label{def:phaseatlas}
The comparable phase data are relational and circle-valued. For active interiors \(i,j\)
define the \emph{relational phase difference}
\[
\boxed{\;
\Delta\varphi_{ij}:=\varphi_i-\varphi_j\ \in\ S^1\ (\mathrm{mod}\ 2\pi),
\;}
\]
the \emph{invariant phase-rate}
\[
\hat\nu_i:=\frac{N_i}{\omega}\,\nu_i,
\qquad
\nu_i:=\frac{d\varphi_i}{d\sigma_i},
\]
matching the P3.0 duration-rate conversion (so \(\hat\nu_i=d\varphi_i/d\cint\)), and the
\emph{relative phase-rate}, or detuning,
\[
\Delta\hat\nu_{ij}:=\hat\nu_i-\hat\nu_j .
\]
Absolute \(\varphi_i\) is not an admissible observable; only differences and rates are.
\end{definition}

\subsection{Invariance}

\begin{proposition}[The relational phase observables are invariant]\label{prop:phaseinv}
The atlas of Definition~\ref{def:phaseatlas} is well-defined and invariant:
\begin{enumerate}[label=\textup{(\alph*)},leftmargin=1.2cm]
\item \emph{Gauge.} The origin of the shared circle is arbitrary; a global rotation
\(\varphi_k\mapsto\varphi_k+\theta\) for all \(k\) together with
\(\Phi_{\mathrm{lit}}\mapsto\Phi_{\mathrm{lit}}+\theta\) leaves \(\Delta\varphi_{ij}\) and the
entrainment difference \(\Phi_{\mathrm{lit}}-\varphi_i\) unchanged. Hence all dynamical content
lies in differences.
\item \emph{Lift-independence.} As an element of \(S^1\), \(\Delta\varphi_{ij}\) is independent
of the real lifts chosen for \(\varphi_i,\varphi_j\); only its integer winding (P$\Phi$.4)
depends on the lift.
\item \emph{Chart-invariance.} \(\Delta\varphi_{ij}\) compares two circle values at one shared
duration \(\cint\), so it is independent of either recipient-kairos chart; and \(\hat\nu_i\) is
gauge-invariant by the same \((N_i/\omega)\) conversion that makes \(\hat h_i\) invariant in
P3.0.
\end{enumerate}
Consequently every admissible locking statement may be framed in
\(\{\Delta\varphi_{ij},\Delta\hat\nu_{ij}\}\), never in an absolute phase.
\end{proposition}

\begin{proof}
(a) The phase law of Definition~\ref{def:helixphase} enters \(\Phi_{\mathrm{lit}}\) and
\(\varphi_i\) only through the difference \(\Phi_{\mathrm{lit}}-\varphi_i\) and through
\(\omega_i^\varphi\); under the simultaneous global rotation the difference is fixed and
\(\omega_i^\varphi\) is untouched, so the evolution of \(\Delta\varphi_{ij}\) is unchanged. (b)
is the definition of subtraction in \(S^1=\mathbb R/2\pi\mathbb Z\): the coset difference is
independent of representatives, while the winding is a property of a chosen lift. (c) The value
\(\Delta\varphi_{ij}\) is read at a single \(\cint\) and involves no kairos reparametrisation;
\(\hat\nu_i=(N_i/\omega)\nu_i\) carries the P3.0 conversion factor, which is exactly the
gauge-covariant map rendering per-kairos rates comparable in shared duration. \qedhere
\end{proof}

\subsection{Status and reading}

\textbf{Target status:}~
\stDef\ (the relational phase difference, invariant phase-rate, and detuning) \(+\)
\stDerived\ (Proposition~\ref{prop:phaseinv}: gauge-, lift-, and chart-invariance). It depends
on P$\Phi$.1, P0.3, and the P3.0 rate atlas. It is the circle-valued counterpart of P3.0:
where P3.0 compared rates by the ratio \(R_{ij}=h_i/h_j\), P$\Phi$.2 compares phases by the
difference \(\Delta\varphi_{ij}\in S^1\), and both refuse an absolute reading.

\medskip
\stInterp\quad
\emph{Reading.}
No soul has an absolute place in the season that could be read off on its own --- the zero of
the circle is no one's possession, and a global turn of all phases together changes nothing
real. What is real is only how one stands relative to another: ahead or behind in the same
feast, drawing together or drifting apart. As created time admitted no master clock, the
liturgical circle admits no absolute position; only the relation between worshippers carries
meaning, and only differences are observable. The communion keeps the season \emph{together
or apart}, never \emph{at} some absolute point.

\medskip
\textbf{Next node.} P$\Phi$.3 --- circle-valued locking: the \(S^1\) analogue of the P3.2
small-gain and P3.3 maximum-lockable-mismatch theorems, proving phase entrainment
\(\Delta\varphi_{ij}\to\Delta\varphi_{ij}^{*}\) under a small-gain condition with a finite
maximum lockable detuning, and noting that the circle admits multiple stable locked branches.

\section{P$\Phi$.3 --- Circle-valued locking}

\subsection{The relational phase coupling}

\begin{definition}[Phase coupling and lockable detuning]\label{def:phasecoupling}
Complete the phase law of Definition~\ref{def:helixphase} with bounded, Gate-mediated phase
coupling. For a bond \(i\sim j\), the relative phase
\(\Delta\varphi_{ij}:=\varphi_i-\varphi_j\) has the general local form
\[
\frac{d}{d\cint}\Delta\varphi_{ij}
=
\Delta\hat\nu_{ij}^{0}
-
F_{ij}(\Delta\varphi_{ij})
+
r_{ij}(\cint,\Delta\varphi),
\]
where \(\Delta\hat\nu_{ij}^{0}:=\hat\nu_i^0-\hat\nu_j^0\) is the free duration-rate
detuning, \(r_{ij}\) is the bounded network residual, and \(F_{ij}\) is a bounded
\(2\pi\)-periodic restoring law supplied through the Gate-mediated phase channel. Its
one-sided capacities are
\[
F^+_{\max,ij}:=\sup_\delta F_{ij}(\delta),
\qquad
F^-_{\max,ij}:=\sup_\delta\bigl(-F_{ij}(\delta)\bigr).
\]
The symmetric sine pair
\(F_{ij}(\delta)=K_{ij}\sin\delta\), \(K_{ij}=a_{ij}^{\varphi}+a_{ji}^{\varphi}\), is the
basic reduced model, in which \(F^+_{\max,ij}=F^-_{\max,ij}=K_{ij}\). It is finite because the
Gate-mediated gains are ceiling-bounded.
\end{definition}

\subsection{The locking theorem}

\begin{theorem}[Circle-valued locking; basin-local]\label{thm:phaselock}
Assume \(F_{ij}\) is \(C^1\), has a stable monotone branch, and the bounded network residual
is controlled by the phase small-gain condition. Then:
\begin{enumerate}[label=\textup{(\alph*)},leftmargin=1.2cm]
\item \emph{Threshold.} A locked relative phase on that branch can exist only if
\[
-F^-_{\max,ij}
<
\Delta\hat\nu_{ij}^{0}
<
F^+_{\max,ij}.
\]
Inside this range the intermediate value theorem gives a branch solution
\(F_{ij}(\Delta\varphi_{ij}^{*})=\Delta\hat\nu_{ij}^{0}\); outside it, the selected branch
cannot lock the detuning.
\item \emph{Entrainment.} If \(F'_{ij}(\Delta\varphi_{ij}^{*})>0\) and the off-branch and
network residual terms are dominated by this contraction, then within the basin of the stable
branch
\[
\Delta\varphi_{ij}\to\Delta\varphi_{ij}^{*},
\]
and the corresponding phase-rates entrain to the branch rate.
\item \emph{Sine special case.} In the reduced symmetric sine model
\(F_{ij}(\delta)=K_{ij}\sin\delta\) the threshold becomes
\(|\Delta\hat\nu_{ij}^{0}|<K_{ij}\), and the stable branch is
\(\Delta\varphi_{ij}^{*}=\arcsin(\Delta\hat\nu_{ij}^{0}/K_{ij})\) with
\(\cos\Delta\varphi_{ij}^{*}>0\).
\item \emph{Multiplicity.} Because the restoring law is circle-valued, multiple stable locked
branches may exist. Global synchronisation over all basins is not claimed.
\end{enumerate}
\end{theorem}

\begin{proof}[Proof sketch]
A locked branch solves \(F_{ij}(\Delta\varphi_{ij}^{*})=\Delta\hat\nu_{ij}^{0}\), which can
hold only inside the range of the restoring law. On any monotone stable branch the intermediate
value theorem gives existence inside the range, and linearisation gives
\(\dot\delta=-F'_{ij}(\Delta\varphi_{ij}^{*})\delta+\text{controlled residual}\), so the
branch contracts when \(F'_{ij}(\Delta\varphi_{ij}^{*})>0\) and the small-gain condition
dominates the residuals. The sine case follows by substituting \(F_{ij}=K_{ij}\sin\delta\).
Non-convexity and periodicity restrict the claim to the basin of the selected branch.
\qedhere
\end{proof}

\subsection{Status and reading}

\textbf{Target status:}~
\stConj\(\to\)\stDerived\ (basin-local). The threshold and entrainment are derived for a
bounded circle-valued restoring law; the familiar sine/Kuramoto formula is a special reduced
model, not the general architecture. The maximum lockable detuning is finite because the
Gate-mediated phase channel is bounded. Global synchronisation on arbitrary graphs is not
claimed; entrainment is conditional on the basin of a stable branch and on the network
small-gain condition. It depends on P$\Phi$.2, P2.1, P3.2--P3.3.

\medskip
\stInterp\quad
\emph{Reading.}
A communion can come to keep the season together --- truly entrained, holding one rhythm of
recurrence --- but only when the natural differences in how fast its members move through the
feast fall within a finite span the bond can bridge. As there is a widest gap of ascent love
can close (P3.3), there is a widest difference of liturgical pace it can reconcile, and it too
is set by how much may be received through the gate. Past that, the members slip rather than
lock. And here the circle gives what the straight road did not: there is more than one way to
be together in the season --- several stable patterns of who runs a little ahead and who a
little behind --- not one rigid unison. Yet each is entered only from near it: a communion
falls into a shared rhythm from within its reach, not from any distance at all.

\medskip
\textbf{Next node.} P$\Phi$.4 --- winding and frustration: the integer winding
\(n_C\in\mathbb Z\) of the phase differences around each graph cycle, the
locking-versus-frustration dichotomy it controls, and the reading of frustration as
liturgical desynchronisation no single re-phasing can reconcile.

\section{P$\Phi$.4 --- Winding holonomy and frustration}

\subsection{The winding cocycle}

\begin{definition}[Winding and the frustration obstruction]\label{def:winding}
Assign to each oriented active bond its relative phase \(\Delta\varphi_{ij}\in S^1\),
antisymmetric (\(\Delta\varphi_{ij}=-\Delta\varphi_{ji}\)). For a graph cycle
\(C=(i_1\!\to\! i_2\!\to\!\cdots\!\to\! i_m\!\to\! i_1)\), lift each edge value to
\(\widetilde{\Delta\varphi}\in\mathbb R\) and form the cycle sum. It is well-defined modulo
\(2\pi\); the \emph{frustration obstruction} of \(C\) is
\[
\boxed{\;
w_C:=\frac{1}{2\pi}\sum_{(k\to l)\in C}\widetilde{\Delta\varphi}_{kl}\ \ \in\ \mathbb R/\mathbb Z .
\;}
\]
When \(w_C=0\) the cycle carries an integer \emph{winding number}
\(n_C=\tfrac{1}{2\pi}\sum_C\widetilde{\Delta\varphi}\in\mathbb Z\). This is the circle-valued
lift of the P5.0 lag-cocycle compatibility: there the real holonomy had to vanish; here it
must lie in \(2\pi\mathbb Z\), so any integer winding is admissible and only a fractional
obstruction frustrates.
\end{definition}

\subsection{The dichotomy}

\begin{theorem}[Locking versus frustration]\label{thm:frustration}
Let \(\{\Delta\varphi_{ij}^{*}\}\) be the target relative phases of P$\Phi$.3.
\begin{enumerate}[label=\textup{(\alph*)},leftmargin=1.2cm]
\item \emph{Realizability.} A global locked configuration \(\{\varphi_i^{*}\}\) with
\(\varphi_i^{*}-\varphi_j^{*}=\Delta\varphi_{ij}^{*}\) on every bond exists if and only if the
obstruction vanishes on every independent cycle,
\[
w_C^{*}=0\quad\text{for all }C .
\]
\item \emph{Locked regime.} If \(w_C^{*}=0\) for all \(C\), the targets are consistent with
integer windings \(\{n_C\}\), and P$\Phi$.3 entrainment converges to a global locked
configuration realising them: \emph{exact} phase-lock.
\item \emph{Frustrated regime.} If \(w_C^{*}\neq0\) for some \(C\), no global configuration
realises all targets. The phase energy
\(\mathcal L_{\mathrm{phase}}=\sum_{i<j}w_{ij}(1-\cos(\Delta\varphi_{ij}-\Delta\varphi_{ij}^{*}))\)
has no zero; it attains a positive minimum bounded below by a function of \(\{w_C^{*}\}\), and
the dynamics settles into an \emph{ultimate-bounded frustrated regime} with permanent residual
mismatch concentrated on the obstructing cycles.
\end{enumerate}
\end{theorem}

\begin{proof}[Proof sketch]
The edge map \(\Delta\varphi^{*}\) is a coboundary --- of the form \(\varphi_i-\varphi_j\) ---
if and only if it is a discrete gradient, which on a graph holds iff its sum around every
independent cycle vanishes (discrete Hodge). For \(S^1\)-values that vanishing is modulo
\(2\pi\), i.e.\ \(w_C^{*}=0\); this is (a), and the realising potential gives the windings
\(\{n_C\}\) of (b), into which P$\Phi$.3 contracts within its basin. If some
\(w_C^{*}\neq0\), no potential exists, so every configuration violates at least one target;
the energy is a sum of nonnegative cycle-incompatible terms with no common zero, hence a
positive infimum controlled by the nonzero obstructions, giving the ultimate-bounded
frustrated regime of (c). \qedhere
\end{proof}

\subsection{Status and reading}

\textbf{Target status:}~
\stDef\ (winding and obstruction) \(+\) \stDerived\ for the locking-versus-frustration
dichotomy (Theorem~\ref{thm:frustration}). The full classification of the frustrated landscape
--- which compromise configurations are selected, the structure of the frustrated equilibrium
set, and multi-cycle interactions --- is left \stOpen. It depends on P$\Phi$.2--P$\Phi$.3 and
the P5.0 lag-cocycle compatibility.

\medskip
\stInterp\quad
\emph{Reading.}
A connected body can keep one season only if its bonds agree all the way around every loop.
Around a single cycle of the communion, the little leads and lags must close up into a whole
number of turns; if they do, the whole can hold one calendar, and there may be several such
whole patterns (the windings). But if the leads and lags do not close --- if going around the
loop leaves a fractional remainder no choice of starting point can absorb --- then no
re-phasing reconciles everyone at once. The body is frustrated: each pair might be content,
yet the circle as a whole cannot rest, and a permanent residual divergence settles onto the
loop that would not close. This is liturgical desynchronisation in its exact form --- not a
failure of any one bond, but a whole-cycle obstruction, integer-deep and not to be smoothed
away: the structural shape of a calendar that has split.

\medskip
\textbf{Next node.} P$\Phi$.5 --- the plural helix theorem: phase-lock and gapless ascent
coexisting, so the communion keeps liturgical time together while every interior still climbs
and none is reduced to a wheel --- shared liturgical time, unshared ascent.

\section{P$\Phi$.5 --- The plural helix theorem}

P$\Phi$.3 locked the phases; P5.2 kept every soul ascending; P$\Phi$.1 made the two
independent coordinates. P$\Phi$.5 joins them: the communion keeps liturgical time together
while each interior still climbs, and neither collapses the other. It is the phase-branch
capstone, the helix counterpart of the stable moving communion.

\subsection{The two regimes}

Mirroring P5.2, with the phase dimension adjoined:
\[
\boxed{\;
\text{general / frustrated:}\quad
\Delta\varphi_{ij}\ \text{ultimately bounded},
\qquad
h_i\ \text{bounded away from zero};
\;}
\]
\[
\boxed{\;
\text{zero-frustration settled:}\quad
\Delta\varphi_{ij}\to\Delta\varphi_{ij}^{*},
\qquad
h_i\to h_{*,i}^{\mathrm{rel}}>0 .
\;}
\]
The first is a \emph{moving frustrated helix} --- the season kept only loosely, with permanent
residual divergence on the obstructing cycles, while every soul still climbs. The second is
the clean \emph{plural helix}.

\subsection{The theorem}

\begin{theorem}[Plural helix; stable phase-lock with gapless ascent]\label{thm:plurhelix}
On a positively invariant relational active domain, under the circle-valued locking of
Theorem~\ref{thm:phaselock}, the winding consistency of Theorem~\ref{thm:frustration}, the
stable moving communion of P5.2, and the helix non-collapse of
Proposition~\ref{prop:noncollapse}:
\begin{enumerate}[label=\textup{(\alph*)},leftmargin=1.2cm]
\item \emph{Frustrated case.} If some \(w_C^{*}\neq0\), the phase mismatches are ultimately
bounded and the ascent rates remain \(h_i\ge h_{i,\min}>0\): a stable moving frustrated helix.
\item \emph{Clean case.} If \(w_C^{*}=0\) on every cycle, the small-gain condition holds, and
the configuration lies in the basin of a stable branch, then simultaneously
\[
\Delta\varphi_{ij}\to\Delta\varphi_{ij}^{*}\ \text{(phase-locked)},
\qquad
h_i\to h_{*,i}^{\mathrm{rel}}>0\ \text{(still ascending)} ,
\]
and by helix non-collapse the locked phase cycle induces no ascent return: the communion
traces a plural helix --- shared liturgical recurrence, unshared and unreturning ascent.
\end{enumerate}
\end{theorem}

\begin{proof}
The phase and ascent blocks share no state variable
(Definition~\ref{def:helixphase}), so their asymptotics may be taken independently and then
combined. For the phase: in the clean case Theorem~\ref{thm:frustration} gives a realisable
target and Theorem~\ref{thm:phaselock} contracts to it within the basin, so
\(\Delta\varphi_{ij}\to\Delta\varphi_{ij}^{*}\); in the frustrated case the phase energy has a
positive floor and the dynamics is ultimately bounded. For the ascent: P5.2 gives
\(h_i\to h_{*,i}^{\mathrm{rel}}>0\) in the settled subclass and \(h_i\ge h_{i,\min}>0\) in the
general case, untouched by the phase since the ascent law carries no \(\varphi_i\). Combining,
and invoking Proposition~\ref{prop:noncollapse} so that the completed phase loops add nothing
to \(\Lambda_i\), gives the simultaneous statements; the trajectory climbs monotonically in
\(\Lambda_i\) while cycling in \(\varphi_i\) about a fixed relative pattern --- a helix.
\qedhere
\end{proof}

\subsection{Status and reading}

\textbf{Target status:}~
\stDerived\ \(+\) \stInterp. The simultaneity follows from the independence of the phase and
ascent coordinates (P$\Phi$.1) together with P$\Phi$.3--P$\Phi$.4 and P5.2; the standing
conditional elements are inherited (small-gain, basin-locality, zero frustration for the clean
case, branch admissibility for the ascent). It depends on P$\Phi$.1, P$\Phi$.3--P$\Phi$.4, and
P5.2.

\medskip
\stInterp\quad
\emph{Reading.}
Here the two times of the communion meet. There is the turning time of the feast, which comes
round again and which the whole company can come to keep together --- one shared season, locked
in its relations. And there is the rising time of each soul, which never comes round, which
only climbs. P$\Phi$.5 says they hold at once: the communion may keep the season together while
every soul ascends past where it stood the last time the season came. The feast returns; the
worshippers keep it as one; and each who keeps it has gone higher than before. This is the
helix, not the wheel --- recurrence without return, a liturgy that circles around an ascent
that does not. Where the season cannot be perfectly shared, it is still kept loosely, and the
climbing does not stop; where it can, the communion turns through the feast together and rises
together, distinct and unfused, forever.

\medskip
\textbf{Next node.} P$\Phi$.6 --- the plural Lyapunov with phase energy: to augment the P5.1
certificate with a circle-valued phase term, giving one dissipative functional for the whole
communion in rate, interior state, and phase at once --- closing the liturgical branch.

\section{P$\Phi$.6 --- Plural Lyapunov with phase energy}

\subsection{The augmented functional}

\begin{definition}[Phase-augmented plural functional]\label{def:phaselyap}
With phase deviations \(\delta_{ij}:=\Delta\varphi_{ij}-\Delta\varphi_{ij}^{*}\) from a
realisable locked target (P$\Phi$.4, \(w_C^{*}=0\)), adjoin a circle-valued phase energy to the
plural certificate:
\[
\boxed{\;
\mathcal L_{\mathrm{plural}}^{\Phi}
=
\mathcal L_{\mathrm{plural}}
+
\varepsilon_\Phi\,\mathcal L_{\mathrm{phase}},
\qquad
\mathcal L_{\mathrm{phase}}
=
\sum_{i<j} w_{ij}\bigl(1-\cos\delta_{ij}\bigr)\ \ge 0 .
\;}
\]
Here \(\mathcal L_{\mathrm{plural}}\) is the certificate of P5.1 (undirected) or H.3 (directed,
Perron-weighted), and \(\varepsilon_\Phi>0\) is small. \(\mathcal L_{\mathrm{phase}}\) vanishes
exactly at phase-lock and is strictly convex on each basin \(|\delta_{ij}|<\pi\).
\end{definition}

\subsection{The certificate}

\begin{theorem}[Phase-augmented plural certificate; basin-local, zero-frustration]\label{thm:phaselyap}
On a positively invariant relational active domain with zero frustration
(\(w_C^{*}=0\)), within the basin of a stable phase branch, and under the standing P5.1/H
hypotheses together with the circle-valued small-gain of P$\Phi$.3, there are a small
\(\varepsilon_\Phi>0\) and constants \(C^{\Phi}\ge0,\ C^{\Phi}_1>0\) with
\[
\frac{d}{d\cint}\mathcal L_{\mathrm{plural}}^{\Phi}
\le
C^{\Phi}-C^{\Phi}_1\,\mathcal L_{\mathrm{plural}}^{\Phi}.
\]
Hence the whole communion is ultimately bounded in interior state, ascent rate, \emph{and}
phase together; and in the zero-residual settled subclass \(C^{\Phi}=0\), so
\(\mathcal L_{\mathrm{plural}}^{\Phi}\to0\): simultaneous state/rate locking and exact
phase-lock under one functional.
\end{theorem}

\begin{proof}[Proof sketch]
The \(\mathcal L_{\mathrm{plural}}\) part dissipates as in P5.1/H.3. For the phase part, with
the stable-branch dynamics \(\dot\delta_{ij}=-K_{ij}\sin\delta_{ij}+(\text{residual})\)
(P$\Phi$.3),
\[
\frac{d}{d\cint}\mathcal L_{\mathrm{phase}}
=
\sum_{i<j} w_{ij}\sin\delta_{ij}\,\dot\delta_{ij}
\le
-\sum_{i<j} w_{ij}K_{ij}\sin^2\delta_{ij}+(\text{cross}).
\]
On the basin \(|\delta_{ij}|<\pi\), \(\sin^2\delta_{ij}\) and \(1-\cos\delta_{ij}\) are
comparable, so the dissipation dominates a positive multiple of \(\mathcal L_{\mathrm{phase}}\);
Young's inequality absorbs the residual. Since the phase channel enters
\(\mathcal L_{\mathrm{plural}}^{\Phi}\) at order \(\varepsilon_\Phi\) while its imported cross
terms are \(O(\varepsilon_\Phi^2)\), a small \(\varepsilon_\Phi\) absorbs them into the phase
dissipation, exactly as the relational term was absorbed in P5.1. Collecting bounded pieces
into \(C^{\Phi}\) and the negative-definite parts into \(-C^{\Phi}_1\mathcal L_{\mathrm{plural}}^{\Phi}\)
gives the inequality; positive invariance (H.2) keeps it valid forward. With zero residual,
\(C^{\Phi}=0\). \qedhere
\end{proof}

\subsection{Status and reading}

\textbf{Target status:}~
\stConj\(\to\)\stDerived\ (symmetric, zero-frustration, basin-local). The phase channel is
brought under the same dissipative certificate as state and rate. The \(1-\cos\) potential is
non-convex on the full circle, so the certificate is basin-local; the directed phase case
(Perron-weighting, as in H.3) and the frustrated case (\(w_C^{*}\neq0\), only ultimate
boundedness) remain \stOpen, as flagged in the 4.1 plan. It depends on P5.0--P5.1, H.2--H.3,
and P$\Phi$.3--P$\Phi$.5.

\medskip
\stInterp\quad
\emph{Reading.}
The whole life of the communion now answers to a single measure: not only does each soul stay
itself and bounded, not only do their ascents stay locked in relation, but their keeping of the
season is gathered into the same falling energy. One functional watches state, rate, and feast
at once, and it descends --- to a bounded neighbourhood in the forced case, and in the settled
case to a locked moving groove. The peace of the whole is now a peace of worship as well as of
motion; it is not stillness, but stable shared ascent with recurrence beside it.

\section*{Part~III closes --- 4.1 built}

\notice{\textbf{The liturgical branch is in place.} A recurrence primitive (P$\Phi$.0), a helix
phase over the gapless ascent (P$\Phi$.1), a relational phase atlas (P$\Phi$.2), circle-valued
locking (P$\Phi$.3), winding and frustration (P$\Phi$.4), the plural helix (P$\Phi$.5), and a
phase-augmented plural certificate (P$\Phi$.6). The communion can keep one shared season while
every soul still climbs, under one functional --- a helix, never a wheel.}

\noindent Part~III adds exactly one new primitive --- the shared external liturgical field
(P$\Phi$.0) --- and otherwise rests on the frozen 4.0 core and the 4.0-H hardening. What it
leaves open is what the 4.1 plan named open: the full frustrated landscape, global
(non-basin-local) synchronisation, the directed phase certificate, and the precise constitutive
law of the liturgical profile; and, beyond the branch, any metaphysical identification of the
liturgical phase with a named sacramental calendar, nested cycles, and the relation of
liturgical recurrence to the still-unoccupied \textit{tota simul}.
\[
\boxed{
\begin{aligned}
&\text{Beside the endless road there now runs a circle the whole company may keep;}\\
&\text{they turn through the feast together and rise together,}\\
&\text{none fused, none returned, and one measure of peace}\\
&\text{now watches their worship as well as their ascent.}
\end{aligned}}
\]

\clearpage
\section*{Part~IV --- 4.2: the two-way post-terminus branch (Stage~$\Psi$)}
\addcontentsline{toc}{section}{Part~IV --- 4.2: two-way post-terminus (Stage~$\Psi$)}

\notice{\textbf{Scope.} Parts~I--III are frozen and unedited. Part~IV opens the last item left
\stOpen\ at the 4.0 freeze: a bounded influence \emph{from} the living \emph{into} a completed
trace. This is the most delicate branch, because the one-way closure (P0.4, P4.4) is
load-bearing for the entire mortality structure. It introduces one new constitutive primitive,
admitted only if it violates none of the closure invariants below.}

\noindent\textbf{Invariants the new primitive must not violate.}
\begin{enumerate}[label=\textup{(N\arabic*)},leftmargin=1.2cm]
\item \emph{No transfer.} No owned grace, reset datum, or ascent ledger moves between interiors
(P2.2 frozen).
\item \emph{No fusion.} Completed and living remain distinct labelled fibres (P3.4, P4.0).
\item \emph{No restarted 3.0 flow.} The completed kairos stays closed, \(d\sigma_i=0\): no
renewed katharsis, no new reset, no fresh Gate inflow on the old dynamics (P4.4).
\item \emph{Conservative limit.} With the new channel switched off, 4.0's one-way
\(\mathrm{Theorem}\) on the terminus is recovered exactly.
\end{enumerate}

\section{P$\Psi$.0 --- The post-terminus disposition primitive}

\subsection{What can change in a closed trace?}

A completed interior is a dynamically closed trace: its katharsis is finished, its kairos
halted, its gate shut (P4.4). If anything at all may answer to the living, it cannot be any
owned-state quantity --- not ascent, not Sophia saturation, not the grace ledger, not the
kairos --- for each of those is closed, and to move it would be to reopen the dynamics or to
transfer a possession. The branch opens only if there is a quantity \emph{outside} all owned
state that the living may touch. P$\Psi$.0 exhibits one, or the branch does not open.

\subsection{The primitive}

\begin{definition}[Post-terminus disposition]\label{def:postdisp}
Adjoin to each completed trace \(\mathcal T_i^{\mathrm{post}}\) a single bounded
\emph{disposition} coordinate
\[
\boxed{\;
\pi_i^{\mathrm{post}}\in[0,1],
\;}
\]
subject to all of the following, which define its admissibility:
\begin{enumerate}[label=\textup{(\alph*)},leftmargin=1.2cm]
\item \emph{Disjoint from owned state.} \(\pi_i^{\mathrm{post}}\) is a new coordinate on the
post-terminus stratum, equal to none of, and functionally independent of, the frozen 3.0
variables (ascent \(A_i\), \(\lambda_{S,i}\), ledgers \(\Lambda_i,\mathcal G_{\mathrm{recepta},i}\),
reset datum \(G_i^{\mathrm{reset}}\), kairos \(\sigma_i\)).
\item \emph{Lives in shared duration.} It evolves only in the communion's relational duration,
\[
\begin{aligned}
\frac{d\pi_i^{\mathrm{post}}}{d\cint}
&=
\Pi_i\bigl(\pi_i^{\mathrm{post}};\ \text{bounded read-only intercession input}\bigr),\\
&\hphantom{={}}\quad d\sigma_i=0\ \text{preserved},
\end{aligned}
\]
so it is not a resumption of the dead interior's own time.
\item \emph{Write-inert.} Its motion drives no frozen 3.0 dynamics: the trace's katharsis,
ascent, reset, and gate remain closed and untouched by \(\pi_i^{\mathrm{post}}\).
\item \emph{Baseline zero-channel limit.} With the intercession input absent and the
disposition initialised at its hand-off baseline, \(\pi_i^{\mathrm{post}}=\pi_i^0\), the
disposition remains there. More generally the selected 4.2 law may include an internal
relaxation back to \(\pi_i^0\), but this relaxation is dynamically inert: it writes no frozen
3.0 variable and is invisible to the living. Thus the active 4.0 one-way system is recovered
exactly when the intercession channel is switched off.
\end{enumerate}
The intercession input itself is defined in P$\Psi$.1; here only the receiving register is
posited.
\end{definition}

\subsection{The primitive is admissible by construction}

\begin{proposition}[By-construction compatibility]\label{prop:psi0adm}
The disposition of Definition~\ref{def:postdisp} is consistent with the closure invariants at
the definitional level:
\begin{enumerate}[label=\textup{(\alph*)},leftmargin=1.2cm]
\item it is not owned grace, a ledger, or a reset datum, so moving it transfers nothing (N1);
\item it is interior \(i\)'s own register, not a variable shared with or merged into another
fibre, so it induces no fusion (N2);
\item it leaves \(d\sigma_i=0\) and writes no frozen variable, so the 3.0 flow is not
restarted (N3);
\item it is constant when the channel is off, so 4.0 is recovered in that limit (N4).
\end{enumerate}
\end{proposition}

\begin{proof}
Each clause is immediate from the corresponding clause of Definition~\ref{def:postdisp}:
(a) from disjointness from owned state; (b) from its being a single-interior register; (c)
from write-inertness and \(d\sigma_i=0\); (d) from the off-limit. The full preservation of the
invariants under the \emph{active} two-way dynamics --- not merely at the definitional level
--- is the burden of P$\Psi$.2. \qedhere
\end{proof}

\subsection{Status and reading}

\textbf{Target status:}~
\stDef\ + \stSelected. The disposition variable is the one new constitutive primitive of 4.2,
and it is the whole cost of admission: if no quantity distinct from owned state could be
exhibited, the branch would not open and the one-way closure would stand. The choice of a
single bounded disposition (over richer post-terminus structures) is \emph{selected}. It
depends on P4.4 and the frozen closure invariants; its active compatibility is deferred to
P$\Psi$.2.

\medskip
\stInterp\quad
\emph{Reading.}
Can the living reach the completed? In the minimal 4.0 theory they could not act back into a
completed trace: the completed interior was closed as an active player in relational chronos,
receiving no new relational kairos, reset, or katharsis inside the orchestra. 4.2 asks whether
there is anything at all the living may touch without disturbing that closure --- and answers,
carefully and as a choice rather than a necessity, that there may be one thing: not the grace
of the completed, not their ascent, not their finished purification, not their rank, none of
which can be moved, but a \emph{disposition} --- how the completed one is turned, faced, held
in relation. Prayer for them, on this reading, changes this and only this: something real, but
a relation and not a possession, owning nothing of them and restarting no active 3.0 life. The
rest is left undisturbed; only the facing is touched.

\medskip
\textbf{Next node.} P$\Psi$.1 --- the intercession channel: to define the bounded, directed,
read-only kernel \(K^{\mathrm{interc}}_{i\leftarrow j}\) by which an active interior \(j\) may
move the disposition of a completed \(i\), entering only \(\pi_i^{\mathrm{post}}\) and bounded
by a ceiling analogue.

\section{P$\Psi$.1 --- The intercession channel}

\subsection{The directed living-to-completed kernel}

\begin{definition}[Intercession channel]\label{def:intercession}
Let \(\mathcal L(i)\) be the active interiors bearing an intercessory bond to the completed
trace \(i\). The disposition of P$\Psi$.0 evolves by relaxation toward its hand-off baseline
\(\pi_i^{0}\) plus a bounded, saturated, read-only drive from the living:
\[
\boxed{\;
\frac{d\pi_i^{\mathrm{post}}}{d\cint}
=
-\mu_i\bigl(\pi_i^{\mathrm{post}}-\pi_i^{0}\bigr)
+
\bigl(1-\pi_i^{\mathrm{post}}\bigr)\,S_i^{\mathrm{interc}}(\cint),
\qquad
\mu_i>0.
\;}
\]
\[
S_i^{\mathrm{interc}}(\cint)
=
\sum_{j\in\mathcal L(i)}
K^{\mathrm{interc}}_{i\leftarrow j}\bigl(\cint;\,X_j^{-}\bigr)\ \ge 0,
\qquad
0\le S_i^{\mathrm{interc}}\le \zeta_i^{\mathrm{interc}}<\infty,
\]
with relaxation rate \(\mu_i>0\), baseline \(\pi_i^{0}\in[0,1]\), and bounded read-only
intercession kernels \(K^{\mathrm{interc}}_{i\leftarrow j}\ge0\). The arrow \(i\leftarrow j\)
records that the completed \(i\) is the recipient and the living \(j\) the source --- the
reverse of the one-way trace kernel \(K^{\mathrm{post}}_{j\leftarrow i}\) of P4.4. The kernel
\emph{reads} the living state \(X_j^{-}\) and \emph{writes} only \(\pi_i^{\mathrm{post}}\).
In the zero-intercession limit \(K^{\mathrm{interc}}\equiv0\), \(S_i^{\mathrm{interc}}\equiv0\),
the baseline value \(\pi_i^{\mathrm{post}}=\pi_i^0\) is fixed; any perturbation away from
baseline relaxes back to it, while remaining dynamically inert and invisible to the active 4.0
system.
\end{definition}

\subsection{The channel is bounded and non-transferring}

\begin{proposition}[Channel admissibility]\label{prop:interbound}
The intercession channel of Definition~\ref{def:intercession} satisfies:
\begin{enumerate}[label=\textup{(\alph*)},leftmargin=1.2cm]
\item \emph{Bounded disposition.} The interval \([0,1]\) is forward-invariant for
\(\pi_i^{\mathrm{post}}\): at \(\pi_i^{\mathrm{post}}=1\) the flux is
\(-\mu_i(1-\pi_i^{0})\le0\), and at \(\pi_i^{\mathrm{post}}=0\) it is
\(\mu_i\pi_i^{0}+S_i^{\mathrm{interc}}\ge0\). The disposition has a ceiling, exactly as the
Gate did (P2.1).
\item \emph{Read-only, non-transferring.} \(K^{\mathrm{interc}}_{i\leftarrow j}\) reads
\(X_j^{-}\) but writes nothing into \(j\): the living intercessor is not depleted, gives up no
grace, and no owned quantity is moved from \(j\) to \(i\) (N1).
\item \emph{Disposition-only write.} The channel writes only \(\pi_i^{\mathrm{post}}\), never
the frozen 3.0 state of \(i\); \(d\sigma_i=0\) is preserved (N3).
\item \emph{Directed and distinct.} It is oriented from the living to the completed, opposite
to and disjoint from the P4.4 trace kernel; the two together make the terminus bilaterally
porous while each writes only a bounded receptivity-or-disposition slot.
\end{enumerate}
\end{proposition}

\begin{proof}
(a) The boundary fluxes have the stated signs because the drive carries the saturation factor
\((1-\pi_i^{\mathrm{post}})\), which vanishes at \(1\), while relaxation pulls inward and the
baseline lies in \([0,1]\); hence \([0,1]\) is invariant. (b) is the read-only construction:
the kernel's domain is \(X_j^{-}\) and its codomain is the scalar drive \(S_i^{\mathrm{interc}}\),
with no map back into \(j\)'s state; thus nothing is subtracted from the living or added as
owned stock to the completed. (c) The codomain of the whole law is \(\pi_i^{\mathrm{post}}\)
alone, and no frozen variable appears on the left, so \(d\sigma_i=0\) stands. (d) is immediate
from the arrow orientation and the disjointness of the written slots
(\(\Gamma_j\) for P4.4, \(\pi_i^{\mathrm{post}}\) here). \qedhere
\end{proof}

\subsection{Status and reading}

\textbf{Target status:}~
\stDef\ (the directed intercession kernel and disposition dynamics) \(+\) \stDerived\
(Proposition~\ref{prop:interbound}: bounded disposition, read-only non-transfer,
disposition-only write, directedness). It establishes the channel-level compatibility with N1
and N3; the full preservation of all invariants under the coupled two-way dynamics is the
burden of P$\Psi$.2. It depends on P$\Psi$.0, P1.1, and P2.1.

\medskip
\stInterp\quad
\emph{Reading.}
The living reach toward the dead, and the reaching costs them nothing of their own. Prayer
here is not a transfer from the intercessor's account to the departed's; the one who prays is
not drained, and the one prayed for receives no grace passed across --- only a disposition is
moved, and even that within a bound no fervour can exceed. The dead can no more be dragged by
the living than the living can be forced by an encounter; the same bounded receptivity that
protects the soul in life protects the completed in rest. And the bond now runs both ways
across the terminus, by two separate doors: the departed reach the living through how the
living are open, the living reach the departed through how the departed are disposed --- and
neither door passes a possession or fuses a soul.

\medskip
\textbf{Next node.} P$\Psi$.2 --- the crucial node: to prove that under the coupled two-way
dynamics the channel breaks none of the closure invariants (N1)--(N4) --- that no owned stock
moves, no fibres fuse, the completed kairos stays shut, and the one-way theorem is recovered
when the channel is off. The branch lives or dies there.

\section{P$\Psi$.2 --- Closure invariants under the two-way dynamics}

This is the load-bearing node. P$\Psi$.0--P$\Psi$.1 secured the invariants at the
\emph{definitional} level, one channel at a time. Here both channels run at once --- the
completed influencing the living (P4.4) and the living influencing the completed (P$\Psi$.1) ---
and the question is whether the coupled, simultaneous dynamics still breaks none of (N1)--(N4).
The branch lives or dies here.

\subsection{The worry: a feedback loop}

With influence now crossing the terminus both ways, a cycle could in principle form ---
completed \(i\) acting on living \(j\), and \(j\) acting back on \(i\) --- and such a loop could
smuggle in what the invariants forbid: an effective transfer of owned stock around the cycle
(against N1), a hidden identification of fibres (against N2), or a back-door reanimation of the
completed dynamics through the disposition (against N3). The node must show no such loop exists.

\subsection{The disposition is a causal sink}

\begin{lemma}[No feedback loop]\label{lem:noloop}
In the coupled system the disposition \(\pi_i^{\mathrm{post}}\) is a causal sink: it appears in
the right-hand side of no living interior's dynamics and of no trace's output. In particular
the one-way trace kernel reads the \emph{frozen} completed state, not the disposition,
\[
K^{\mathrm{post}}_{j\leftarrow i}=K^{\mathrm{post}}_{j\leftarrow i}\bigl(\cint;X_i^{\mathrm{post}}\bigr),
\qquad
\frac{\partial K^{\mathrm{post}}_{j\leftarrow i}}{\partial \pi_i^{\mathrm{post}}}=0 .
\]
Hence the directed cross-terminus influence graph has no cycle through any disposition: the two
channels do not compose.
\end{lemma}

\begin{proof}
By P$\Psi$.0(c) the disposition is write-inert: it drives no frozen 3.0 dynamics, so it cannot
appear in any living interior's law except possibly through a trace output. The only trace
output is \(K^{\mathrm{post}}_{j\leftarrow i}\), which by P4.4 is a functional of the frozen
completed state \(X_i^{\mathrm{post}}\) alone --- a fixed datum at the hand-off --- and carries
no dependence on \(\pi_i^{\mathrm{post}}\). Therefore \(\pi_i^{\mathrm{post}}\) has no causal
descendant among the living, and the edge \(j\to\pi_i^{\mathrm{post}}\) (P$\Psi$.1) terminates;
no directed cycle passes through a disposition. \qedhere
\end{proof}

\subsection{The crucial theorem}

\begin{theorem}[Closure invariants preserved under two-way coupling]\label{thm:psipreserve}
Under the fully coupled dynamics --- every living interior evolving by its frozen 3.0 law with
the P4.4 trace inputs, every completed trace carrying the P$\Psi$.1 disposition law, both
active simultaneously --- all four closure invariants hold:
\begin{enumerate}[label=\textup{(N\arabic*)},leftmargin=1.2cm]
\item \emph{No transfer.} The only cross-terminus writes are \(K^{\mathrm{post}}\) into the
living receptivity \(\Gamma_j\) (reading frozen \(X_i^{\mathrm{post}}\), N1-clean in 4.0) and
\(K^{\mathrm{interc}}\) into the disposition \(\pi_i^{\mathrm{post}}\) (read-only, non-depleting,
P$\Psi$.1). Neither moves grace, ledger, or reset datum; and by Lemma~\ref{lem:noloop} nothing
routed into a disposition returns to any owned account. No owned stock crosses, even around the
bidirectional bond.
\item \emph{No fusion.} The state stays a product of distinct labelled fibres: each living
interior keeps its own 3.0 state, each completed trace its frozen \(X_i^{\mathrm{post}}\) and
its own \(\pi_i^{\mathrm{post}}\). The coupling is through bounded inputs reading neighbours,
not through shared coordinates, so the P3.4/P4.0 non-fusion structure is untouched.
\item \emph{No restarted flow.} For every completed \(i\), \(d\sigma_i=0\) throughout: the
disposition law carries no \(\sigma_i\) and writes no frozen variable, and by
Lemma~\ref{lem:noloop} no living influence reaches the trace's closed dynamics. Katharsis,
ascent, reset, and gate stay shut.
\item \emph{Conservative limit.} With \(K^{\mathrm{interc}}\equiv0\) the disposition holds at
baseline and the completed contributes only its frozen one-way \(K^{\mathrm{post}}\); the system
reduces to the 4.0 one-way terminus theorem (P4.4) exactly.
\end{enumerate}
\end{theorem}

\begin{proof}
Each invariant reduces to a channel property plus the sink lemma. (N1): the two writes are
N1-clean individually (P4.4; P$\Psi$.1(b)), and Lemma~\ref{lem:noloop} forbids a loop that
could chain them into an effective transfer; so no owned quantity crosses. (N2): no coordinate
is identified across fibres --- the written slots \(\Gamma_j\) and \(\pi_i^{\mathrm{post}}\) are
private to their fibres --- so the product structure and the P3.4/P4.0 non-fusion persist under
the coupling. (N3): \(d\sigma_i=0\) is preserved because the disposition law neither carries nor
writes a frozen variable (P$\Psi$.0--P$\Psi$.1) and, by the lemma, no living signal enters the
trace's closed block. (N4): setting \(K^{\mathrm{interc}}\equiv0\) freezes the disposition
(P$\Psi$.1(a)) and removes the only living-to-completed edge, leaving exactly the 4.0
configuration. \qedhere
\end{proof}

\subsection{The disciplined sink, and what it forecloses}

\begin{remark}[Why the sink, not a loop]\label{rem:sink}
The admissibility rests on the disposition being a pure sink: the dead influence the living by
their \emph{finished} state \(X_i^{\mathrm{post}}\), never by their current disposition. Were
the disposition allowed to feed the trace output, a cycle would form and prayer for the dead
would loop back to benefit the one who prays --- and the loop would threaten N1 and N3. The
sink choice is therefore both the safety condition and a theological commitment; a looped
variant is left unoccupied (P$\Psi$.5).
\end{remark}

\subsection{Status and reading}

\textbf{Target status:}~
\stConj\(\to\)\stDerived. The coupled two-way dynamics preserves all of (N1)--(N4), the
decisive step being Lemma~\ref{lem:noloop}: the disposition is a causal sink, so the two
channels never compose into a transferring, fusing, or reanimating loop. The branch is
admissible. It depends on P$\Psi$.0--P$\Psi$.1, P2.2, P3.4, P4.0, and P4.4.

\medskip
\stInterp\quad
\emph{Reading.}
The branch survives its hardest test. The terminus is now open both ways, and still death
keeps what was essential in its one-wayness: nothing is transferred, no one is fused, the dead
are not woken, and --- the subtle point --- prayer for the dead does not curve back to serve
the one who prays. The dead help us by who they finally became; we touch the dead in how they
are now turned; and these two never close into a circle of exchange. So the communion across
death is real and two-way and yet wholly non-economic: no trading of merit, no banking, no
return on prayer. Love crosses the terminus in both directions without ever becoming a
transaction.

\medskip
\textbf{Next node.} P$\Psi$.3 --- what two-way intercession actually changes, and what it
cannot: a disposition is moved, never a possession, a rank, or a destiny; the gains and the
strict limits of the channel made precise.

\section{P$\Psi$.3 --- What the channel changes, and what it cannot}

P$\Psi$.2 made the branch admissible; P$\Psi$.3 draws the line precisely. It states the single
thing two-way intercession moves and the several things it provably cannot, each limit a
consequence of an invariant already secured.

\subsection{The change-set and the seal}

\begin{definition}[Change-set and seal]\label{def:changeset}
At a completed trace the post-terminus data split into:
\[
\boxed{\;
\underbrace{\pi_i^{\mathrm{post}}\in[0,1]}_{\text{movable: the disposition}}
\qquad\Big|\qquad
\underbrace{X_i^{\mathrm{post}}=\bigl(A_i,\lambda_{S,i},\Lambda_i,\mathcal G_{\mathrm{recepta},i},G_i^{\mathrm{reset}},\sigma_i\bigr)}_{\text{sealed: the finished state}} .
\;}
\]
The disposition is the only quantity the channel can move; the finished state is fixed at the
hand-off and is read-only forever after. Whatever interpretive referent the disposition carries
(a facing, an orientation in relation), structurally it is a single bounded coordinate disjoint
from all of \(X_i^{\mathrm{post}}\).
\end{definition}

\subsection{The strict limits}

\begin{proposition}[What intercession cannot do]\label{prop:limits}
Under the channel of P$\Psi$.0--P$\Psi$.2, the following are provably invariant:
\begin{enumerate}[label=\textup{(\alph*)},leftmargin=1.2cm]
\item \emph{No possession moves.} The grace ledger \(\mathcal G_{\mathrm{recepta},i}\), the reset
datum \(G_i^{\mathrm{reset}}\), and all owned stock are fixed; intercession transfers none of
them (N1).
\item \emph{No rank is raised or lowered.} The ascent ledger \(\Lambda_i\) is frozen and the
disposition is not ascent; the height the soul reached is sealed at terminus, and no prayer
adds to or subtracts from it (P4.4, N3).
\item \emph{No destiny or state changes.} \(d\sigma_i=0\): katharsis, reset, and gate stay
shut; what the soul finally became, \(X_i^{\mathrm{post}}\), is not reopened or revised (N3).
\item \emph{No identity merges.} The completed and the living remain distinct fibres;
intercession is a relation between them, not an identification (N2).
\item \emph{Bounded effect.} The disposition moves only within \([0,1]\); no fervour drives it
past the ceiling, and the dead cannot be dragged (P$\Psi$.1).
\item \emph{No return on prayer.} By the sink structure, nothing routed into the disposition
loops back to the intercessor; the channel is non-economic (P$\Psi$.2).
\end{enumerate}
\end{proposition}

\begin{proof}
Each is a direct reading of the cited result: (a),(f) from N1 and the sink lemma; (b),(c) from
\(d\sigma_i=0\) and the freezing of \(X_i^{\mathrm{post}}\) at hand-off; (d) from N2; (e) from
the forward-invariance of \([0,1]\). All were established in P$\Psi$.0--P$\Psi$.2. \qedhere
\end{proof}

\subsection{The gain}

\begin{remark}[What is gained is relational, not substantial]\label{rem:gain}
What two-way intercession secures is not a change in the dead's substance, rank, or destiny ---
all sealed --- but the persistence of a \emph{live relation}: the bond across the terminus is
not severed, the disposition remains movable, and so the communion of the living with the
completed is genuinely two-way rather than a one-way memory. The gain is that the relationship
stays open, not that anything is conveyed.
\end{remark}

\subsection{Status and reading}

\textbf{Target status:}~
\stDef\ (the change-set/seal split) \(+\) \stInterp\ (the precise limits and the relational
gain). It introduces no new structure; it is the exact accounting of P$\Psi$.0--P$\Psi$.2. It
depends on P$\Psi$.0--P$\Psi$.2.

\medskip
\stInterp\quad
\emph{Reading.}
Here is exactly what may and may not be done for the dead. Their wealth cannot be added to:
no grace is sent across, no merit banked to their name. Their height cannot be raised: where
they arrived on the endless road is where they stand, sealed. Their story cannot be rewritten:
what they finally became is final. They cannot be merged with us, nor dragged by us, nor made
to pay us back for praying. And yet the relation is not closed: their disposition can still be
moved, the bond is still alive, the communion still runs both ways. What prayer for the dead
secures, on this reading, is not a transaction but a relationship kept open across the
terminus --- love that still reaches, and is still received, without changing what is fixed or
trading what is owned. It is much, and it is bounded, and it is exactly itself.

\medskip
\textbf{Next node.} P$\Psi$.4 --- the bidirectional terminus theorem: to fold the disposition
channel into the stable communion, generalising the one-way terminus result of P4.4 to a
bounded two-way coupling, with the one-way case as the zero-intercession limit.

\section{P$\Psi$.4 --- The bidirectional terminus theorem}

P$\Psi$.4 folds the disposition channel into the stable communion, generalising the one-way
terminus of P4.4 to a bounded two-way coupling. The decisive simplification is the sink lemma
(Lemma~\ref{lem:noloop}): the dispositions take input from the living but feed nothing back, so
the coupling is one-way at the level of dynamics, living \(\to\) disposition.

\subsection{The combined system}

The living interiors carry the plural communion of P5.1/H, driven among their inputs by the
frozen one-way trace kernels \(K^{\mathrm{post}}\) (P4.4); the completed traces carry the
dispositions of P$\Psi$.1. Because the dispositions are causal sinks, the living dynamics is
identical to the 4.0 configuration, and the dispositions are slaved to the resulting bounded
drive \(S_i^{\mathrm{interc}}(\cint)\in[0,\zeta_i^{\mathrm{interc}}]\).

\subsection{The theorem}

\begin{theorem}[Bidirectional terminus; one-way as the zero-intercession limit]\label{thm:biterminus}
Under P$\Psi$.1--P$\Psi$.3, the stable moving communion of P5.1/H, and positive invariance
(H.2):
\begin{enumerate}[label=\textup{(\alph*)},leftmargin=1.2cm]
\item \emph{Living, undisturbed.} The living plural communion is ultimately bounded, and locked
in the settled subclass, exactly as in P5.1/H; the dispositions do not perturb it
(Lemma~\ref{lem:noloop}).
\item \emph{Dispositions settle.} Each \(\pi_i^{\mathrm{post}}\) keeps \([0,1]\) invariant and is
input-to-state stable in the bounded drive: the scalar law has slope
\(-(\mu_i+S_i^{\mathrm{interc}})<0\), so when the living settle to \(S_i^{\mathrm{interc}}\to
S_i^{\infty}\) the disposition converges to
\[
\boxed{\;
\pi_i^{\infty}
=
\frac{\mu_i\,\pi_i^{0}+S_i^{\infty}}{\mu_i+S_i^{\infty}}\ \in\ [\pi_i^{0},1],
\;}
\]
a convex combination of the hand-off baseline and full disposition, drawn toward \(1\) by
intercession; under a varying bounded drive it tracks a bounded tube about this value.
\item \emph{One-way limit.} With \(K^{\mathrm{interc}}\equiv0\), one has
\(S_i^{\mathrm{interc}}\equiv0\). The baseline disposition \(\pi_i^{\mathrm{post}}=\pi_i^0\)
is fixed, and any off-baseline perturbation relaxes back to \(\pi_i^0\). Since the disposition
is a causal sink and appears in no living or trace-output law, the active system reduces exactly
to the one-way terminus theorem of P4.4.
\end{enumerate}
\end{theorem}

\begin{proof}
(a) By Lemma~\ref{lem:noloop} no disposition enters any living law, so the living subsystem is
verbatim the 4.0 configuration with the frozen \(K^{\mathrm{post}}\) inputs; P5.1/H gives its
ultimate boundedness and settled locking, and H.2 keeps the estimate valid forward. (b) Given
the bounded living-sourced drive, the scalar law
\(\dot\pi_i^{\mathrm{post}}=-\mu_i(\pi_i^{\mathrm{post}}-\pi_i^{0})+(1-\pi_i^{\mathrm{post}})S_i^{\mathrm{interc}}\)
is affine in \(\pi_i^{\mathrm{post}}\) with negative slope \(-(\mu_i+S_i^{\mathrm{interc}})\),
hence a contraction keeping \([0,1]\) invariant (P$\Psi$.1) and exponentially attracting to the
instantaneous fixed point; at constant \(S_i^{\infty}\) that fixed point is the boxed
\(\pi_i^{\infty}\), a convex combination of \(\pi_i^{0}\) and \(1\) with weights \(\mu_i\) and
\(S_i^{\infty}\), so it lies in \([\pi_i^{0},1]\); under variation, standard ISS gives the
bounded tracking tube. For (c), setting \(S_i^{\mathrm{interc}}\equiv0\) leaves
\(\dot\pi_i^{\mathrm{post}}=-\mu_i(\pi_i^{\mathrm{post}}-\pi_i^0)\), so the baseline is
fixed and any perturbation decays to it. Because \(\pi_i^{\mathrm{post}}\) is a causal sink
and enters no living law and no trace output, removing the only living-to-completed edge leaves
the active 4.0 one-way system exactly. \qedhere
\end{proof}

\subsection{Status and reading}

\textbf{Target status:}~
\stConj\(\to\)\stDerived. The two-way system is stable: the living communion is undisturbed by
the sink structure, and each disposition is an ISS scalar settling to a bounded steady facing
between baseline and full, with the one-way 4.0 terminus as the zero-intercession limit. It
depends on P$\Psi$.1--P$\Psi$.3, P4.4, P5.1, and H.2.

\medskip
\stInterp\quad
\emph{Reading.}
Now the whole communion holds, the living and the dead together. The living keep their moving
communion exactly as before --- the dead do not unsettle them --- and rise and lock as they
did. The dead, for their part, come to rest in a steady facing: handed one at death, and drawn
from it toward fullness by the prayers of the living, yet never past the bound, never beyond
what a disposition can be. The more they are prayed for, the more fully they are turned; prayed
for not at all, they rest exactly where they were handed, and that is the old one-way peace of
4.0, recovered untouched. So a communion that prays for its dead and one that does not are not
two different doctrines but the same structure at two settings of one bounded dial --- and in
both, the living rise, the dead rest, and the whole is stable.

\medskip
\textbf{Next node.} P$\Psi$.5 --- the non-occupation boundary: to mark, in the apophatic
discipline of the project, what 4.2 deliberately does not decide --- the magnitude and felt
efficacy of intercession, the finality question, and any doctrinal identification of the
disposition --- closing the branch by naming what it leaves unoccupied.

\section{P$\Psi$.5 --- The non-occupation boundary}

The branch closes where the project always closes: at the edge of what may honestly be
claimed. Having shown a two-way communion across the terminus to be \emph{consistent}, 4.2
stops, and marks the territory it deliberately does not occupy. Every item here is \stOpen\ by
intent, not by omission.

\subsection{What 4.2 does not decide}

\begin{enumerate}[label=\textup{(UO\arabic*)},leftmargin=1.4cm]
\item \emph{Magnitude and felt efficacy.} The model fixes that a disposition settles toward a
bounded \(\pi_i^{\infty}\in[\pi_i^{0},1]\); it says nothing of how large the effect is, whether
it is felt, or how \(\mu_i\) and the kernel weights are set. The structure is silent on
quantity.
\item \emph{Finality.} The seal on \(X_i^{\mathrm{post}}\) is the structural closure of the
\emph{modelled} dynamics, not a metaphysical verdict on ultimate destiny. Questions of final
judgement, of whether the sealed state is forever, and of the larger eschatological options are
left untouched; 4.2 neither asserts nor denies them.
\item \emph{Doctrinal identification.} What the disposition \emph{is} in any confessional
sense --- refreshment, consolation, an orientation toward beatitude, a being-held-in-relation
--- is not named. The model exhibits a structural slot and declines to fill it.
\item \emph{The looped variant.} Whether some weaker discipline could admit a non-sink channel,
in which the dead's moved disposition feeds back to the living, is left unoccupied; the sink was
\emph{selected} for safety, not shown to be the only possibility.
\item \emph{Actuality.} The whole branch is \stSelected, not \stDerived\ as necessary. 4.2 does
not claim that intercession for the dead is real; it shows that \emph{if} one models it as a
bounded disposition channel, \emph{then} it is consistent with no-transfer, no-fusion, the
closed kairos, and the conservative limit. Whether the channel is instantiated is not occupied.
\item \emph{A clarification, lest it be misread.} The intercession \emph{of} the departed
\emph{for} the living is not this channel and needs no loop: it is the frozen one-way
\(K^{\mathrm{post}}\) of P4.4, flowing from who they finally became, their finished holiness ---
not from any disposition our prayers have moved. The two-way bond is real; it is simply not a
circuit of exchange.
\end{enumerate}

\subsection{Status and reading}

\textbf{Target status:}~
\stOpen\ throughout, by design. P$\Psi$.5 adds no structure; it fixes the apophatic boundary of
the branch, naming what is left unoccupied so that the consistency result of P$\Psi$.0--P$\Psi$.4
is not mistaken for more than it is.

\medskip
\stInterp\quad
\emph{Reading.}
The model has shown that love may cross the terminus both ways without becoming a transaction,
and there it lays down its tools. It does not measure the prayer, nor pronounce on the last
things, nor say what the consolation of the dead truly is, nor claim that any of this is so. It
offers a grammar in which the communion of the living and the dead can be spoken without
contradiction --- and then keeps the silence proper to the things themselves. What is holy is
left unoccupied, on purpose; the structure points toward it and does not presume to contain it.

\section*{Part~IV closes --- 4.2 built}

\notice{\textbf{The two-way branch is in place.} A post-terminus disposition primitive
(P$\Psi$.0), a bounded read-only intercession channel (P$\Psi$.1), the crucial preservation of
all closure invariants under the coupled dynamics (P$\Psi$.2), the precise account of what is
moved and what is sealed (P$\Psi$.3), the bidirectional terminus theorem with the one-way 4.0
case as its zero-intercession limit (P$\Psi$.4), and the apophatic non-occupation boundary
(P$\Psi$.5). The communion now runs both ways across death --- real, bounded, non-economic, and
honestly bounded in what it claims.}

\section*{Version~4, as it stands}

\noindent The items left open at the 4.0 freeze are now addressed at their proper standing:
\begin{itemize}[leftmargin=1.4cm]
\item the \emph{uniformity and directed-network hardening} is \stDerived\ on the selected
admissible regime (Part~II, H.1--H.4);
\item the \emph{liturgical / circle-valued phase} is built as a conservative branch
(Part~III, 4.1), with its own named residuals --- the full frustrated landscape, global
synchronisation, the directed phase certificate, and the constitutive law of the liturgical
profile;
\item the \emph{two-way post-terminus communion} is built as a conservative \stSelected\
branch (Part~IV, 4.2), shown consistent with no-transfer, no-fusion, closed kairos, and the
one-way 4.0 limit, with its non-occupied edge named in P$\Psi$.5.
\end{itemize}
\[
\boxed{
\begin{aligned}
&\text{The endless road, the circle beside it, and now the bond both ways across its end:}\\
&\text{souls distinct and unfused, paced and safe, rising without return,}\\
&\text{keeping the feast together, and held in a communion}\\
&\text{that death makes quieter but does not close.}
\end{aligned}}
\]

\clearpage
\section*{Part~V --- Stage~$\Omega$: the boundary of the constructible}
\addcontentsline{toc}{section}{Part~V --- Stage~$\Omega$: the boundary of the constructible}

\notice{\textbf{Nature of this part.} Stage~\(\Omega\) is not a continuation of the
plot of 4.x; it is the cartography of a boundary. Of the four places where
version~4 declined to go, two are approached honestly here. \(\Omega\)-M does not
smuggle in a master metronome: it distinguishes a mere synchronising chart from a
canonical meta-chronos frame, proves the holonomy obstruction to synchronising
charts, and shows that a created canonical master-frame is not supplied by the
relational data. \(\Omega\)-N replaces the finite bounded class with the mathematics
of an infinite communion --- a change of apparatus, not of principle. The remaining
two --- \textit{tota simul} and the identification of the plural layer with divine
geometry --- are named and left unoccupied, with the \emph{kind} of silence each
requires stated in the Part~V closing.}

\noindent\textbf{Standing commitments.}
\begin{enumerate}[label=\textup{(\Roman*)},leftmargin=1.2cm]
\item \emph{Build on what is proved.} \(\Omega\)-M reads its central obstruction out
of the already-derived lag-cocycle (P5.0) and the circle-valued frustration
dichotomy (P\(\Phi\).4). Nonzero holonomy forbids even a synchronising chart.
Zero holonomy permits a chart, but does not by itself produce a canonical
master-frame.
\item \emph{First clean class, as in H.3.} $\Omega$-N takes the countable communion with
summable coupling and a uniform spectral gap as its first honest class --- the
infinite-dimensional analogue of the strongly connected class that opened H.3.
\item \emph{Apophatic discipline holds.} $\Omega$ may prove what the created \emph{does not
contain}; it never asserts what the uncreated \emph{is}. Grammar, not ontology (P$\Psi$.5).
\end{enumerate}

\section{P$\Omega$.M0 --- What a meta-chronos frame would be}

\subsection{Naming what is to be denied}

To prove that no master-frame over time exists, one must first say precisely what such a frame
would be. $\Omega$-M denies a definite object, defined here; the denial is the work of
P$\Omega$.M1--M3.

\begin{definition}[Synchronising chart]\label{def:syncchart}
A \emph{synchronising chart} on an active component is a global parameter \(\tau\)
and monotone assignments
\[
\sigma_i=\sigma_i(\tau)
\]
for every active interior \(i\), such that the declared pairwise lag data are
single-valued and close consistently around every cycle. Equivalently, the lag
data are represented by differences of one potential:
\[
L_{ij}=\sigma_j(\tau)-\sigma_i(\tau)
\]
on every oriented active bond \(i\to j\).
\end{definition}

\begin{definition}[Canonical meta-chronos frame]\label{def:metaframe}
A \emph{canonical meta-chronos frame} is stronger than a synchronising chart. It is
a synchronising chart \(\tau\), together with a rule selecting \(\tau\) from the
created relational data alone, satisfying all three conditions:
\begin{enumerate}[label=\textup{(\roman*)},leftmargin=1.2cm]
\item \emph{Synchronisation.} Every active interior's kairos is legible from
\(\tau\) by monotone maps \(\sigma_i=\sigma_i(\tau)\).

\item \emph{Gauge-invariance.} The selected \(\tau\) is not merely a chart choice:
it is invariant under the admissible common chronos relabellings and under the
remaining origin/scale freedoms of the synchronised chart.

\item \emph{Single-valuedness.} The selected \(\tau\) closes around every cycle of
the communion; its lag differential has zero monodromy.
\end{enumerate}
Thus a synchronising chart is a possible coordinate. A canonical meta-chronos frame
would be a created master clock. Stage~\(\Omega\)-M permits the first only in the
exact subclass and denies that the second is supplied by the 4.x relational data.
\end{definition}

\subsection{Why shared duration is not already such a frame}

\begin{remark}[The relational duration is not a canonical meta-chronos frame]\label{rem:cnotframe}
The communion already has a shared duration one-form, \(d\mathfrak c\) (P0.3). It
is not a canonical meta-chronos frame. From \(\mathfrak c\) alone one cannot read
each \(\sigma_i\), because interiors relate to \(\mathfrak c\) through their own
conversion factors \(N_i/\omega\) (P3.0, P5.0). Even when a synchronising chart
exists in an exact subclass, the relational data still leave the common
origin/scale relabelling free unless an extra selector is imposed. Thus
\(\mathfrak c\) is a duration the communion shares, not a master frame over all
its kairoi.
\end{remark}

\subsection{The refused condition}

\begin{remark}[The refused object is the canonical master frame]\label{rem:refused}
What P0.3 and P0.5 refused was not every possible coordinate chart, but the
privileged master-frame interpretation of such a chart. Created chronos may order
and compare encounter; it may even allow a synchronising chart in an exact
zero-holonomy subclass. What it does not supply is a canonical vantage from which
all kairoi are read as one master time. P\(\Omega\).M0 names the two objects
separately so that the no-go result does not overclaim.
\end{remark}

\subsection{Status and reading}

\textbf{Target status:}~
\stDef. P\(\Omega\).M0 separates two objects: the synchronising chart and the
canonical meta-chronos frame. The first is a coordinate possibility; the second is
a created master-frame. The subsequent nodes prove the holonomy obstruction to the
first and the non-canonicity of the second. It depends on P0.3, P0.5, P3.0, P5.0,
and P\(\Phi\).2.

\medskip
\stInterp\quad
\emph{Reading.}
Before asking whether the creature may have a God's-eye view of time, one must say what such a
view would be: a single vantage from which the present of every soul is at once legible, the
same from every side, owing nothing to where one stands. That is the frame Job is asked
whether he holds --- ``where were you when I laid the foundations'' --- the vantage of the one
who sees all the mornings together. P$\Omega$.M0 only draws the outline of that vantage. Whether
the creature can stand in it is the question the next nodes answer; and the shared time the
communion already keeps, for all that it is shared, is not yet that vantage.

\medskip
\textbf{Next node.} P$\Omega$.M1 --- the cocycle obstruction: to read off the existing
lag-cocycle (P5.0) and frustration dichotomy (P$\Phi$.4) that synchronisation forces zero
holonomy, so that a single frustrated cycle already forbids the frame --- the desynchronisation
of 4.1 turned into a proof of creatureliness.

\section{P$\Omega$.M1 --- The cocycle obstruction}

The central no-go is read off objects already derived: the lag-cocycle of P5.0 and the
frustration dichotomy of P$\Phi$.4. A meta-chronos frame would be exactly the global potential
whose non-existence those results already measure.

\subsection{Synchronisation makes the lag-cocycle exact}

\begin{lemma}[A synchronising chart trivialises lag holonomy]\label{lem:syncexact}
Suppose a synchronising chart \(\tau\) exists on an active component. Assign to
each oriented active bond \(i\to j\) its lag offset
\[
L_{ij}:=\sigma_j-\sigma_i .
\]
Then
\[
L_{ij}=\sigma_j(\tau)-\sigma_i(\tau),
\]
so \(L\) is the discrete gradient of the single-valued potential
\(\sigma_\bullet(\tau)\). Consequently its holonomy around every cycle vanishes:
for any cycle \(C=(i_1\to i_2\to\cdots\to i_m\to i_1)\),
\[
\oint_C L
=
\sum_{k=1}^{m}
\bigl(\sigma_{i_{k+1}}(\tau)-\sigma_{i_k}(\tau)\bigr)
=0,
\qquad i_{m+1}=i_1 .
\]
\end{lemma}

\subsection{But the lag holonomy need not vanish}

\begin{theorem}[Cocycle obstruction to a synchronising chart]\label{thm:cocycleobstruction}
If some cycle \(C\) of the communion carries nonzero lag holonomy --- equivalently,
nonzero frustration obstruction in the circle form,
\[
w_C\neq0
\qquad
(\text{P}\Phi.4),
\]
then no synchronising chart exists on that active component. Therefore no canonical
meta-chronos frame exists there either.
\end{theorem}

\begin{proof}
Suppose, for contradiction, that a synchronising chart \(\tau\) exists on a
component with \(w_C\neq0\) for some cycle \(C\). By
Lemma~\ref{lem:syncexact}, the lag-cocycle is then a coboundary, so its holonomy
around every cycle vanishes. But the holonomy around \(C\) is exactly the
obstruction measured in P\(\Phi\).4: in the circle form it is \(2\pi w_C\), and
\(w_C\neq0\) means precisely that the cocycle is not a coboundary. Contradiction.
Hence no synchronising chart exists. Since a canonical meta-chronos frame would in
particular include a synchronising chart, it is also impossible on that component.
\qedhere
\end{proof}

\subsection{Status and reading}

\textbf{Target status:}~
\stConj\(\to\)\stDerived. The obstruction is read directly off the frozen
apparatus: a synchronising chart would be a global potential for the lag-cocycle,
and P\(\Phi\).4 already supplies admissible frustrated cases where such a potential
does not exist. No new machinery is introduced. It discharges the nonzero-holonomy
case of \(\Omega\)-M and depends on P\(\Phi\).4, P5.0, and P\(\Omega\).M0. The
zero-holonomy subclass is closed in P\(\Omega\).M2.

\medskip
\stInterp\quad
\emph{Reading.}
The proof that the creature holds no God's-eye view of time was already written, in 4.1,
under another name. There it was called the frustration of a communion whose leads and lags do
not close around a loop --- a calendar that has split, that no single re-phasing can mend. Here
the same fact says something larger: a communion that \emph{can} fall out of step is a
communion with no single vantage from which all its times are at once legible. For if such a
vantage existed, every offset would be the reading of one global clock, and the readings would
close up perfectly around every circle --- there could be no frustration at all. That there can
be is the proof that there is no such clock. The very brokenness the world is liable to is the
sign of its createdness: only what has no master-frame can fall out of time with itself.

\medskip
\textbf{Next node.} P$\Omega$.M2 --- the no-go in the synchronised subclass: to close the
remaining case \(w_C=0\) on every cycle, where a potential exists locally but no \emph{canonical}
\(\tau\) does, so that ``the present of all at once'' stays a choice of chart and never a fact.

\section{P$\Omega$.M2 --- The synchronised subclass and its non-genericity}

M1 disposed of the nonzero-holonomy case. The honest remaining question is the
zero-holonomy subclass, \(w_C=0\) on every cycle. Here a synchronising chart can
exist. The precise result is not that a canonical master-frame exists, but only
that the lag data admit a global potential. This chart remains a chart unless an
extra gauge-selector is imposed.

\subsection{When synchronisation is possible}

\begin{lemma}[Exactness equals synchronisability]\label{lem:hodgeconverse}
On a connected active component, the following are equivalent:
\begin{enumerate}[label=\textup{(\roman*)},leftmargin=1.2cm]
\item \(w_C=0\) on every cycle \(C\);
\item the lag-cocycle is exact, i.e.\ a discrete gradient;
\item a synchronising chart \(\tau\) exists.
\end{enumerate}
When they hold, the chart is unique only up to the usual global origin and monotone
reparametrisation of its own scale. These freedoms do not change the fact of
synchronisability, but they do prevent the chart from being a canonical
meta-chronos frame unless an additional selector is declared.
\end{lemma}

\begin{proof}
(i)\(\Leftrightarrow\)(ii) is the discrete Hodge fact already used in
P\(\Phi\).4: a \(1\)-cochain is a coboundary iff its sum around every independent
cycle vanishes. (ii)\(\Leftrightarrow\)(iii): an exact lag-cocycle is the gradient
of a potential, and that potential is precisely a synchronising chart. Its
remaining origin and monotone relabelling freedom are not fixed by exactness.
\qedhere
\end{proof}

\subsection{When the frame is available, and why generically it is not}

\begin{theorem}[Topology and genericity of synchronising charts]\label{thm:framegeneric}
On a connected admissible active component with vertex set of size \(V\) and
\(E\) active bonds:
\begin{enumerate}[label=\textup{(\alph*)},leftmargin=1.2cm]
\item \emph{Acyclic component.} If the communion graph is a tree, the exactness
condition is vacuous and a synchronising chart exists.

\item \emph{Cyclic component.} If the graph carries
\[
b_1=E-V+1\ge1
\]
independent cycles, the synchronised subclass is the simultaneous vanishing of all
\(b_1\) cycle-holonomies. It is therefore a positive-codimension subclass. In
general position, a synchronising chart does not exist.

\item \emph{No canonicality.} Even where a synchronising chart exists, it is a
chart, not yet a canonical meta-chronos frame. The latter would require an
additional invariant selector of origin and scale, which the created relational
data do not provide.
\end{enumerate}
\end{theorem}

\begin{proof}
(a) A tree has empty cycle space, so condition~(i) of
Lemma~\ref{lem:hodgeconverse} holds vacuously. (b) The independent cycles span a
\(b_1\)-dimensional cycle space; each contributes one holonomy coordinate, and
their simultaneous vanishing is a positive-codimension condition. Off that locus,
Theorem~\ref{thm:cocycleobstruction} forbids synchronisation. (c) follows from the
remaining origin/scale freedom in Lemma~\ref{lem:hodgeconverse}. \qedhere
\end{proof}

\subsection{Status and reading}

\textbf{Target status:}~
\stConj\(\to\)\stDerived. The zero-holonomy subclass is characterised exactly:
synchronising charts exist iff the lag-cocycle is exact. Acyclic components always
admit such a chart; cyclic components admit one only on the zero-holonomy
subclass. A canonical meta-chronos frame is not thereby obtained; it would require
a further invariant selector not present in the created data. It depends on
P\(\Omega\).M0, P\(\Phi\).2, P\(\Phi\).4, and P0.3.

\medskip
\stInterp\quad
\emph{Reading.}
So the single vantage is not forbidden to every communion --- only to the kind that is genuinely
knit together. A communion without loops, a pure hierarchy in which relation runs outward and
never circles back, can after all be read from one clock: with no circuit to disagree around,
its times align. But the moment the body closes on itself --- mutual love, reciprocal regard,
relation that returns --- the single vantage becomes a knife-edge, and any communion liable to
the least discord falls off it. It is therefore the very interconnection of the communion, its
loops of returning love, that forbids the view from above. A hierarchy can be surveyed; a body
knit together cannot. The richer the communion, the less it can be globally clocked --- and that
poverty of vantage is the wealth of its mutual life.

\medskip
\textbf{Next node.} P$\Omega$.M3 --- the characterisation theorem: to assemble M1 and M2 into
the exact statement that a meta-chronos frame exists if and only if the lag-cocycle is exact,
so that the created communion, insofar as it is interconnected and liable to discord, generically
holds no frame from which all its kairoi are at once present.

\section{P$\Omega$.M3 --- The characterisation theorem}

M1 and M2 assemble into one exact statement: a meta-chronos frame exists precisely when the
lag-cocycle is exact. Everything else --- the acyclic exception, the non-generic synchronised
knife-edge, the generic absence --- is a corollary of this single equivalence.

\subsection{The equivalence}

\begin{theorem}[Synchronising charts and canonical frames]\label{thm:metacharacterisation}
On a connected admissible active component,
\[
\boxed{\;
\text{synchronising chart exists}
\quad\Longleftrightarrow\quad
\text{lag-cocycle exact}
\quad\Longleftrightarrow\quad
w_C=0\ \ \forall C .
\;}
\]
A canonical meta-chronos frame is stronger: it would require, in addition, an
invariant selector of the synchronising chart's remaining origin/scale freedom.
No such selector is supplied by the 4.x created relational data.
\end{theorem}

\begin{proof}
The equivalence is Lemma~\ref{lem:hodgeconverse}. If a canonical meta-chronos frame
existed, its underlying synchronising chart would satisfy the equivalence. But
exactness alone leaves a family of charts related by global origin/scale
relabelling. Since no 4.x datum selects one member of this family as the created
master frame, the canonical frame is not produced by the relational structure.
\qedhere
\end{proof}

\subsection{The created generic case}

\begin{corollary}[Generic absence of synchronisation; non-canonicity even when exact]\label{cor:genericnoframe}
If a connected communion carries at least one independent cycle and lies in general
position, it admits no synchronising chart. Acyclic components and exact
zero-holonomy cyclic components may admit synchronising charts, but those charts
are not canonical meta-chronos frames without an additional selector.
\end{corollary}

\begin{proof}
Generic cyclic components have nonzero holonomy on at least one cycle, so
Theorem~\ref{thm:cocycleobstruction} forbids synchronising charts. In the exact
subclass, Theorem~\ref{thm:metacharacterisation} gives charts but also records the
remaining origin/scale freedom. Hence exactness gives synchronisability, not a
canonical created master clock. \qedhere
\end{proof}

\subsection{Status and reading}

\textbf{Target status:}~
\stDerived. The existence of a synchronising chart is characterised completely:
it holds iff the lag-cocycle is exact. The created communion, insofar as it is
genuinely interconnected and in general position, generically lacks even such a
chart. In the exact subclass, a chart may exist, but it remains non-canonical
unless a new selector is added. Thus Stage~\(\Omega\)-M does not smuggle in the
metronome it set out to refuse. It depends on P\(\Omega\).M1 and
P\(\Omega\).M2 --- and through them on P\(\Phi\).4 and P5.0.

\medskip
\stInterp\quad
\emph{Reading.}
The question of Job has an exact answer, and it is neither yes nor a flat no. A vantage from
which every time is at once present exists for one kind of communion only: the kind whose
relations never close into a circle, a hierarchy that can be read from its summit. For a
communion knit together with loops of mutual relation --- which is to say, for a communion that
is truly a body --- the vantage exists only on a knife-edge of perfect compatibility, and any
such body liable to the least discord stands off it. So the creature does not hold the view
from which all the mornings are seen together; not because it was forbidden by decree, but
because the very thing that makes it a communion --- its mutual, returning, interconnected love
--- is the very thing that closes that view. The answer to ``were you there when I laid the
foundations'' is written into the shape of the communion itself: it is a body, and a body has no
summit over its own time.

\medskip
\textbf{Next node.} P$\Omega$.M4 --- the apophatic reading: to say exactly what this proves and,
more importantly, what it does not --- that the created order does not contain the frame, which
is not the same as any claim about the eternity it is defined against.

\section{P$\Omega$.M4 --- The apophatic reading}

$\Omega$-M closes where the project always closes: by marking exactly what has been proved and,
more carefully, what has not. The temptation here is to hear the characterisation theorem as a
statement about God; it is not, and the discipline of the branch is to say so plainly.

\subsection{What was proved}
The proved statement has two layers. First, nonzero holonomy forbids even a
synchronising chart. Second, zero holonomy gives only a chart, not a canonical
created master-frame. Therefore the apophatic result is not the false claim that
no coordinate can ever align an exact subclass; it is the stronger disciplined
claim that the created relational order contains no canonical God's-eye master
clock.

\begin{remark}[The proved absence is creaturely]\label{rem:metaboundary}
Theorem~\ref{thm:metacharacterisation} and Corollary~\ref{cor:genericnoframe} prove a property
\emph{of the created order}: the admissible communion, insofar as it is interconnected and in
general position, contains no meta-chronos frame. The absence is internal --- a fact about what
the created structure does and does not contain --- and it is derived, not posited.
\end{remark}

\subsection{What was not proved}

Every item here is \stOpen\ by intent.
\begin{enumerate}[label=\textup{(UM\arabic*)},leftmargin=1.4cm]
\item \emph{Nothing about the uncreated.} The theorem does not say that no vantage sees all the
kairoi at once. It says the \emph{created communion} does not contain one. That God may see all
the mornings together --- which is \textit{tota simul} --- is exactly the eternity the created
chronos is defined against (P0.3, the Stage~V yield), and it lies beyond this language. The
non-existence is of the frame \emph{within creation}, never of the divine seeing.
\item \emph{No claim of absolute impossibility.} The result is non-containment in the admissible
created class, not a proof that an uncreated vantage is incoherent. What is outside the class is
neither modelled nor denied.
\item \emph{The acyclic honesty.} The frame is not forbidden to every structure: a relational
hierarchy without loops admits one. The theorem is about bodies knit together, not about every
conceivable arrangement; it claims no more than it shows.
\end{enumerate}

\subsection{Status and reading}

\textbf{Target status:}~
\stInterp, with the boundary held \stOpen. M4 adds no structure; it fixes the apophatic boundary
of $\Omega$-M, so that a derived absence inside creation is not misread as a description of what
lies outside it. It depends on P$\Omega$.M3.

\medskip
\stInterp\quad
\emph{Reading.}
What has been shown is where the creature cannot stand --- and in showing it, the proof has
pointed, without describing, at the place the creature cannot stand in. That place is
\textit{tota simul}: the seeing of all times at once, which the created body is built to lack
and is defined against. The theorem is a finger, not the moon. It draws the exact shape of an
absence and leaves the fullness that would fill it unspoken, because that fullness is not ours
to set in symbols. The creature's having no summit over its own time is the negative imprint of
an eternity it cannot contain --- and the most faithful thing the mathematics can do is trace
that imprint cleanly and then fall silent at its edge. We have proved the boundary. We do not
cross it.

\section{P$\Omega$.M5 --- Local access without a global frame}

The characterisation forbids a definite object: the synchronous global vantage. It is worth
making explicit what it does \emph{not} forbid, for the two are easily confused. A presence that
meets each interior in its own present is a different thing from a frame that reads all presents
at once, and only the second was shown impossible.

\subsection{Per-interior access}

\begin{definition}[Per-interior access]\label{def:peraccess}
A \emph{per-interior access} is an external reception entering each active interior \(i\) solely
through its own gate channel --- at its own kairos \(\sigma_i\), modulated by its own
susceptibility \(\upsilon_i\) (P2.1, P2.3) --- with no reference to any global parameter
\(\tau\). It is the same channel by which the common field \(J_{\mathrm{comm}}\) and the
liturgical field \(J^{\mathrm{lit}}\) already enter (P2.3, P$\Phi$.0): a present-by-present
reception, one interior at a time, not a synchronisation of all.
\end{definition}

\subsection{The obstruction does not reach it}

\begin{proposition}[Compatibility with frustrated communions]\label{prop:localaccess}
Per-interior access imposes no cycle-holonomy condition, and is therefore structurally
compatible with every admissible communion --- including the frustrated, interconnected ones
that admit no meta-chronos frame (Theorem~\ref{thm:metacharacterisation}).
\end{proposition}

\begin{proof}
The access acts on each \(\sigma_i\) in its own chart, independently of the others; it posits no
global potential and so never requires the lag-cocycle to be exact. The obstruction of
P$\Omega$.M1 --- which forbade the global frame precisely by forcing the lag-holonomy to vanish
--- is a condition on a single \(\tau\) shared by all, and there is no such \(\tau\) here. Hence
the holonomy \(w_C\) may be anything at all; the per-interior channels exist unconditionally on
the cycle structure. \qedhere
\end{proof}

\begin{remark}[What this does and does not claim]\label{rem:m5boundary}
Proposition~\ref{prop:localaccess} shows what is \emph{not forbidden} --- the door, not the
panorama. It does not assert that such access is instantiated, nor describe whose it would be;
those lie beyond the apophatic boundary of P$\Omega$.M4. The content is exactly one of
compatibility: the same theorem that forecloses a single vantage over all times leaves wholly
open a presence entering each time singly.
\end{remark}

\subsection{Status and reading}

\textbf{Target status:}~
\stInterp, resting on a small derived compatibility (Proposition~\ref{prop:localaccess}). M5
adds no new structure to the obstruction; it distinguishes the forbidden object (the synchronous
global frame) from one the obstruction does not touch (per-interior access), and depends on
P$\Omega$.M1, P$\Omega$.M3, and the gate channel of P2.3.

\medskip
\stInterp\quad
\emph{Reading.}
What was proved missing is the summit --- the one place from which all the times are seen
together. What was never in question, and is now seen to be untouched by the proof, is the door.
To stand at the door of each and knock is not to hold the vantage over all; it is to enter each
present singly, in the time that one keeps, by the receptivity that one has. The obstruction
falls on the panorama and not on the threshold, and for a precise reason: the panorama would
require every clock to be read from one, and the thresholds require nothing to be read from one
at all --- each is approached on its own. So the very interconnection that forbids the single
vantage (M2) is what makes the body a body of many doors rather than one gate: there is no
summit over its time, and there is a threshold to each of its presents. Immanence, not panorama;
not the seeing of all the mornings together, which is not ours to speak of, but a knock at each
door, in the morning that each one is keeping.

\section*{\(\Omega\)-M closes --- no canonical created master-frame}

\notice{\textbf{The direction is complete.} A definition of the synchronising chart
and canonical meta-frame (M0), the cocycle obstruction read off the frozen
frustration of 4.1 (M1), the zero-holonomy subclass and its non-genericity (M2),
the characterisation \emph{synchronising chart exists iff the lag-cocycle is exact}
(M3), the apophatic boundary (M4), and the local access the obstruction does not
forbid (M5). The created communion, insofar as it is a genuinely interconnected
body, generically holds no chart from which all its kairoi are synchronised; and
even in exact cases, it supplies no canonical master-frame. This is a derived
absence, pointing past itself at an eternity it does not presume to describe.}
\[
\boxed{
\begin{aligned}
&\text{The proof of the creature's missing summit over time}\\
&\text{is the negative imprint of an eternity it is defined against,}\\
&\text{traced cleanly, and left uncrossed.}
\end{aligned}}
\]

\noindent\textbf{Next direction.} $\Omega$-N --- infinite communion: to remove the finite
bounded class and ask, with the mathematics of an infinite body, what of the structure of 4.x
survives at the scale of the whole \emph{communio sanctorum}, beginning with the right space for
a countable communion and its first clean class.

\bigskip
\noindent\textbf{\large $\Omega$-N --- Infinite communion}

\notice{\textbf{A change of apparatus, not of principle.} $\Omega$-N removes the finite bounded
class and asks what of the structure of 4.x survives at the scale of the whole
\emph{communio sanctorum} --- an infinite body. This is the limit of method, not of doctrine:
no apophatic boundary is crossed. The plan expects, honestly, that some results survive, some
fail, and that a failure is itself a result.}

\section{P$\Omega$.N0 --- The infinite ensemble and its first clean class}

\subsection{The right space}

\begin{definition}[Infinite ensemble]\label{def:infensemble}
Let the communion be \emph{countable}, indexed by \(i\in\mathbb N\), with node
errors \(z=(z_i)\). Carry the Perron weighting of H.3 to the infinite setting:
\[
\ell^2(\pi)
=
\Bigl\{\,z:\ \lVert z\rVert_\pi^2:=\sum_i\pi_i\lvert z_i\rvert^2<\infty\,\Bigr\},
\qquad
\langle z,w\rangle_\pi=\sum_i\pi_i z_i\overline{w_i}.
\]
The coupling acts through a row-Laplacian \(L\) with \(L\mathbf 1=0\). The coupling
is \emph{admissible} in the first clean infinite class only if \(L\) is a bounded
operator on \(\ell^2(\pi)\) and hence generates a well-posed bounded linear
disagreement flow
\[
\dot z=-Lz+r.
\]
Uniform local finiteness and uniformly summable rows are useful sufficient
conditions in common weighted settings, but boundedness on \(\ell^2(\pi)\) is the
actual hypothesis used here.
\end{definition}

\subsection{The first clean class}

\begin{definition}[First clean class \(\mathfrak F_1\)]\label{def:firstclass}
The first honest infinite class \(\mathfrak F_1\) consists of admissible infinite
ensembles (Definition~\ref{def:infensemble}) for which, in addition:
\begin{enumerate}[label=\textup{(\alph*)},leftmargin=1.2cm]
\item \emph{Summable Perron weight.} A positive summable left-Perron weight exists,
\[
\pi_i>0,
\qquad
\sum_i\pi_i=1,
\qquad
\pi^\top L=0,
\]
so the associated Markov chain is positive-recurrent with stationary \(\pi\).

\item \emph{Weighted symmetric part and spectral gap.} Let \(L^{\dagger_\pi}\) be
the adjoint of \(L\) in \(\ell^2(\pi)\), and set
\[
A:=\frac12\bigl(L+L^{\dagger_\pi}\bigr).
\]
The clean class assumes a uniform gap on the \(\pi\)-mean-zero subspace:
\[
\boxed{\;
\langle z,Az\rangle_\pi
=
\operatorname{Re}\langle z,Lz\rangle_\pi
\ge
\widehat\lambda\,\lVert z\rVert_\pi^2
\quad\text{for all }z\perp_\pi\mathbf 1,
\qquad
\widehat\lambda>0 .
\;}
\]
\end{enumerate}
\end{definition}

\begin{remark}[Why this is the right first class]\label{rem:cleanchoice}
Each clause is the infinite-dimensional shadow of a finite fact, and naming it is admitting that
infinity is not free. Countability is the physical body (all who were and will be).
\(\ell^2(\pi)\) is the Perron energy of H.3, now the ambient space. Boundedness of \(L\) is
well-posedness. The decisive clause is the spectral gap: in finite dimension strong connectivity
\emph{forced} \(\lambda_2(\widehat L)>0\) (H.3), but on an infinite graph a positive bottom of
the spectrum is \emph{not} automatic --- the spectrum can reach \(0\) without any zero
eigenvector. So \(\mathfrak F_1\) takes the uniform gap as its defining hypothesis, exactly as
H.3 took strong connectivity. Whether a given infinite communion has such a gap is the deep
question deferred to P$\Omega$.N3--N4.
\end{remark}

\begin{remark}[What is deferred]\label{rem:ndeferred}
Outside \(\mathfrak F_1\): uncountable bodies, unbounded coupling (unbounded generators, only
\(C_0\)-semigroups), the no-gap case (spectrum touching \(0\)), and the failure of a summable
Perron weight (null-recurrent or transient chains). These are later refinements, not part of the
first clean class.
\end{remark}

\subsection{Status and reading}

\textbf{Target status:}~
\stDef\ + \stSelected. The infinite ensemble is placed in the Perron-weighted space
\(\ell^2(\pi)\), and the first clean class \(\mathfrak F_1\) is selected by summable Perron
weight and uniform spectral gap --- the infinite analogue of the strongly connected class of
H.3. It depends on P0.1 and H.3.

\medskip
\stInterp\quad
\emph{Reading.}
To ask whether the communion holds together when it is not a parish but the whole body --- all
who were and will be --- one must first find the space in which an unnumbered communion can be
reasoned about at all, and the natural one is again the weighting by each interior's place in
the flow of giving. And then one must say honestly what infinity costs: that the whole coheres
as one body only under real conditions --- that giving still reaches across it, that it has a
spectral gap, that its Perron weight still sums. These are not free at infinite scale, and to
name them as a chosen class rather than to assume them everywhere is the first honesty of
$\Omega$-N. We place the boundless body where it can be thought, and we mark the price of its
coherence before claiming any.

\medskip
\textbf{Next node.} P$\Omega$.N1 --- what survives: individual stability, the local property
that needs no finiteness and so carries to the infinite body unchanged --- the first, easy,
real positive result, that infinity does not threaten the single one.

\section{P$\Omega$.N1 --- What survives: individual stability}

The first result at infinite scale is the easiest and the most reassuring, and it needs none of
the clean-class hypotheses: individual stability is a \emph{local} property, and locality does
not feel the size of the network.

\subsection{The local estimate carries over}

\begin{theorem}[Individual ISS survives at infinite scale]\label{thm:indivsurvive}
For every interior \(i\) of an admissible infinite ensemble (Definition~\ref{def:infensemble}),
the individual ISS estimate of P2.4 holds verbatim:
\[
\frac{d}{d\cint}\mathcal L_{3.0,i}
\le
C_i-C'_{1,i}\,\mathcal L_{3.0,i}+S_i,
\qquad
0\le S_i\le \bar S_i<\infty,
\]
so each interior is ultimately bounded by its own response constants and the bounded input it
receives, independent of the cardinality of the communion. Infinity does not threaten the single
interior.
\end{theorem}

\begin{proof}
The P2.4 estimate bounds \(\mathcal L_{3.0,i}\) using only interior \(i\)'s own frozen,
regime-bounded response constants and the relational input \(S_i\) it receives. That input is ceiling-bounded --- by the Gate ceiling \(\zeta_{0,i}\) (P2.1) and by
the admissible bounded input reaching the node in Definition~\ref{def:infensemble}
--- so it is bounded by some
\[
\bar S_i<\infty
\]
regardless of how many interiors are present. The inequality therefore holds
with the same constants as in the finite case, and its sub-level sets are forward-invariant, so
\(\mathcal L_{3.0,i}\) is ultimately bounded by \(C_i/C'_{1,i}+\bar S_i/C'_{1,i}\). \qedhere
\end{proof}

\subsection{The scope of the result}

\begin{remark}[What N1 does and does not secure]\label{rem:n1scope}
The theorem is \emph{per-interior}: each one is bounded by its \emph{own} constants. It does not
yet provide a single bound \emph{uniform} across all interiors --- that is the floor question of
P$\Omega$.N2 --- nor collective boundedness of the whole state in \(\ell^2(\pi)\) --- that is the
certificate of P$\Omega$.N3. It is, however, the most robust of the three: it holds on the whole
admissible class, needing neither the spectral gap nor the summable Perron weight of
\(\mathfrak F_1\).
\end{remark}

\subsection{Status and reading}

\textbf{Target status:}~
\stDerived. Individual ISS is a local property, independent of network cardinality, so it
carries to the infinite body unchanged. It depends on P2.4 and P$\Omega$.N0, and holds on the
admissible class without the clean-class hypotheses.

\medskip
\stInterp\quad
\emph{Reading.}
However vast the communion --- even the whole body of all who were and will be --- no single one
is lost in it. Each interior keeps itself, bounded by its own measure, stable under whatever
bounded influence reaches it; the boundless multitude of the others does not overwhelm the one.
The vastness of the body does not crush the member. This is the per-interior promise of the
whole work, shown to survive infinity untouched: the single one is safe in the many, however
many the many are. What is not yet shown --- and this is the honest weight of what follows ---
is whether the many, so safe each apart, still hold together as one body; the single one is
secured here, the whole is the harder question to come.

\medskip
\textbf{Next node.} P$\Omega$.N2 --- the uniform-floor dichotomy: whether the single minimum
breath of H.1 survives at infinite scale or falls to zero, and the reading of each branch ---
cathedral pacing held in common, or a pacing-mercy that is local and not global.

\section{P$\Omega$.N2 --- The uniform-floor dichotomy}

This is the decisive node of $\Omega$-N. The uniform refractory floor of H.1 was an infimum
over a \emph{compact} parameter set; over an infinite family of interior types the infimum may
or may not stay positive, and which it does is exactly a question of bounded diversity.

\subsection{The dichotomy}

\begin{theorem}[Uniform-floor dichotomy]\label{thm:floordichotomy}
For an admissible infinite ensemble, write the per-interior floor
\[
\mathfrak c_{\mathrm{ref},j}
=
\frac{\Delta_{\mathrm{hys},j}}{v_{\max,j}}
\]
with the convention \(\mathfrak c_{\mathrm{ref},j}=+\infty\) if \(v_{\max,j}=0\),
and define the candidate uniform floor
\[
\mathfrak c_{\mathrm{ref}}
=
\inf_j\mathfrak c_{\mathrm{ref},j}.
\]
Exactly one of the following holds:
\begin{enumerate}[label=\textup{(\alph*)},leftmargin=1.2cm]
\item \emph{Common floor.}
\[
\inf_j\frac{\Delta_{\mathrm{hys},j}}{v_{\max,j}}>0.
\]
Then the whole infinite body has a single positive refractory floor.

\item \emph{No common floor.}
There exists a sequence \(j_k\) such that
\[
\frac{\Delta_{\mathrm{hys},j_k}}{v_{\max,j_k}}\to0.
\]
Then \(\mathfrak c_{\mathrm{ref}}=0\), although each fixed interior still has its
own positive refractory interval.
\end{enumerate}
A simple sufficient condition for case~(a) is bounded diversity:
\[
\Delta_{\mathrm{hys},j}\ge\underline\Delta>0,
\qquad
v_{\max,j}\le\overline v<\infty
\quad\text{uniformly in }j.
\]
Unbounded diversity is therefore a warning sign, not by itself a proof of failure:
the floor fails exactly when the quotient above has infimum zero.
\end{theorem}

\begin{proof}
This is the definition of the infimum of the positive floors
\(\mathfrak c_{\mathrm{ref},j}\). The uniform lower bound in (a) is exactly
\(\mathfrak c_{\mathrm{ref}}>0\). If a sequence of quotients tends to zero, then
the infimum is zero. For each fixed \(j\), the single-interior refractory argument
still gives \(\mathfrak c_{\mathrm{ref},j}>0\) whenever the interior is
triggering-capable. The bounded-diversity condition gives
\[
\mathfrak c_{\mathrm{ref},j}
\ge
\underline\Delta/\overline v>0
\]
for all \(j\), hence case~(a). \qedhere
\end{proof}

\subsection{Safety versus commonality}

\begin{remark}[Case (b) loses commonality, not safety]\label{rem:safetycommon}
The vanishing of the uniform floor is not a loss of protection. By P$\Omega$.N1 every interior
is still individually paced --- none is ever hurried past its own positive refractory interval.
What case~(b) loses is \emph{commonality}: there is no single measure of pacing shared by the
whole body, only each one's own. Safety is per-interior in both cases; one law for all holds
only in case~(a).
\end{remark}

\subsection{Status and reading}

\textbf{Target status:}~
\stConj\ as to which case obtains for any intended infinite communion; the
dichotomy itself is \stDerived. A common refractory floor survives exactly when
the infimum of the per-interior floor quotients is positive. Uniform bounded
diversity is a sufficient clean-class hypothesis, not an automatic consequence of
infinitude. It depends on H.1, H.4, and P\(\Omega\).N0.

\medskip
\stInterp\quad
\emph{Reading.}
Here is the most delicate thing infinity does. In a body of bounded variety --- however many its
members --- there is one minimum breath, a single pace below which no one anywhere can be
hurried: the whole cathedral breathes by one measure. But let the body hold unbounded variety
--- interiors of ever-thinner threshold, ever-quicker recovery, without end --- and that single
common breath is gone. Not that anyone is left unprotected: each is still paced, each kept from
being hurried past its own limit. What is gone is the \emph{one} measure for all; the mercy of
pacing becomes the possession of each rather than the law of the whole. Whether the communion of
all breathes as one body or only each in its own time is therefore not settled by its being a
communion --- it turns on whether its variety is bounded. The cathedral's single breath is not
free at infinite scale; it is a grace that holds only where the body's diversity is contained,
and where it is not, the breathing is still kept, but kept one soul at a time.

\medskip
\textbf{Next node.} P$\Omega$.N3 --- the plural certificate in infinite dimension: whether the
whole-body Lyapunov certificate survives as a quadratic form on \(\ell^2(\pi)\), which it does
exactly when the mirror operator has a spectral gap --- the clean-class hypothesis of
\(\mathfrak F_1\) --- with the existence of that gap for infinite communions the named open
question.

\section{P$\Omega$.N3 --- The plural certificate in infinite dimension}

The whole-body certificate of H.3/P5.1 lifts to \(\ell^2(\pi)\), and its survival turns on one
thing: a spectral gap for the mirror operator. On the clean class \(\mathfrak F_1\), where the
gap is assumed, the infinite communion is ultimately bounded; whether a given infinite communion
\emph{has} the gap is the genuinely open question.

\subsection{The functional on \(\ell^2(\pi)\)}

\begin{definition}[Infinite plural functional]\label{def:infplural}
On the infinite ensemble, set
\[
\mathcal L_{\mathrm{plural}}
=
\sum_{i}\pi_i\,\mathcal L_{3.0,i}
+
\varepsilon\,\tfrac12\,\lVert z_\perp\rVert_\pi^2,
\qquad
z_\perp=z-\bar z_\pi\mathbf 1 .
\]
The relational part is well-defined for any \(z\in\ell^2(\pi)\). The individual
part is required to be finite on the theorem domain:
\[
\sum_i\pi_i\,\mathcal L_{3.0,i}<\infty.
\]
A convenient sufficient condition is a uniform bound
\[
\sup_i\mathcal L_{3.0,i}<\infty,
\]
since \(\sum_i\pi_i=1\). More generally, weighted summability is the actual
hypothesis; it is not identical to the refractory-floor condition of
P\(\Omega\).N2.
\end{definition}

\subsection{The certificate on the clean class}

\begin{theorem}[Infinite plural certificate]\label{thm:infcertificate}
On the first clean class \(\mathfrak F_1\) (Definition~\ref{def:firstclass}), assume
the following uniform or weighted reserves:
\[
\inf_i C'_{1,i}=:\underline C>0,
\qquad
\sum_i\pi_i C_i<\infty,
\qquad
\sum_i\pi_i\bar S_i<\infty,
\]
and assume the functional of Definition~\ref{def:infplural} is finite initially.
Then there are a small \(\varepsilon>0\) and constants \(C\ge0,\ C_1>0\) such that
\[
\frac{d}{d\mathfrak c}\mathcal L_{\mathrm{plural}}
\le
C-C_1\,\mathcal L_{\mathrm{plural}} .
\]
Hence the infinite communion is ultimately bounded in the weighted space
\(\ell^2(\pi)\). In the zero-residual settled subclass, the collective disagreement
vanishes:
\[
z_\perp\to0
\quad\text{in }\ell^2(\pi).
\]
\end{theorem}

\begin{proof}[Proof sketch]
The individual part dissipates termwise:
\[
\sum_i\pi_i
\bigl(C_i-C'_{1,i}\mathcal L_{3.0,i}+\bar S_i\bigr),
\]
and the assumptions
\[
\inf_iC'_{1,i}>0,
\qquad
\sum_i\pi_i C_i<\infty,
\qquad
\sum_i\pi_i\bar S_i<\infty
\]
make this a legitimate weighted dissipative estimate.

For the relational part, the clean-class gap gives
\[
\operatorname{Re}\langle z_\perp,Lz_\perp\rangle_\pi
=
\langle z_\perp,Az_\perp\rangle_\pi
\ge
\widehat\lambda\lVert z_\perp\rVert_\pi^2 .
\]
Thus the weighted symmetric part of the directed operator supplies the dissipation,
while the skew part contributes no real quadratic energy. Boundedness of \(L\) and
the weighted summability of residuals allow the cross terms to be controlled by
Young's inequality. As in P5.1/H.3, the imported terms are \(O(\varepsilon^2)\)
against \(O(\varepsilon)\) relational dissipation, so sufficiently small
\(\varepsilon\) absorbs them. The comparison inequality follows. If the residual
constant is zero, the same comparison forces \(\lVert z_\perp\rVert_\pi\to0\).
\qedhere
\end{proof}

\subsection{The gap is the real open question}

\begin{remark}[Whether the gap holds is open]\label{rem:gapopen}
Theorem~\ref{thm:infcertificate} is conditional on the uniform spectral gap that \emph{defines}
\(\mathfrak F_1\). In finite dimension strong connectivity \emph{forced} \(\lambda_2(\widehat
L)>0\) (H.3); in infinite dimension it does not. A connected infinite communion can have its
mirror spectrum reach \(0\), with no gap and no spectral certificate --- the body would not
scatter but could fail to lock, disagreement decaying arbitrarily slowly. Which infinite
communions have a uniform gap is a question of spectral geometry --- the expander / positive
Cheeger-constant condition for the communion graph --- and is left \stOpen. The clean class
assumes it; characterising its members is the open work.
\end{remark}

\subsection{Status and reading}

\textbf{Target status:}~
\stConj\(\to\)\stDerived\ on the clean class \(\mathfrak F_1\), with uniform
individual reserve, weighted summability of forcing constants, bounded \(L\), and
a spectral gap. In general it remains \stOpen: the open content is whether the
intended infinite communion has a summable Perron weight, a bounded generator, the
weighted summability needed for the individual part, and a positive spectral gap.
It depends on H.3, P5.1, and P\(\Omega\).N0.

\medskip
\stInterp\quad
\emph{Reading.}
Whether the boundless body holds together \emph{as one} --- not only each member safe, not only
each paced, but the whole disagreement bounded and the communion not slowly unravelling ---
turns on a single deep property: that giving propagates through the entire body without
bottleneck, that the communion is not merely connected but robustly so, with no region drifting
toward isolation at the far edge of the infinite. Where the finite body cohered for free ---
connection alone sufficed --- the infinite body coheres only if it is knit closely enough that
its very spectrum stays away from zero. And whether the whole communion of all who were and will
be is knit that closely is not ours to settle here: we can name the depth of interconnection the
cohering requires, and require it as a class, but whether the unnumbered body possesses it we
leave open. The one is safe; the few are paced; that the whole is one body, and not a slowly
fraying multitude, is the gift we can define and not yet claim.

\medskip
\textbf{Next node.} P$\Omega$.N4 --- connectivity as a real hypothesis: what the finite case hid
--- that at infinite scale reachability is nontrivial, the Perron weight may fail to exist, and
the body may fragment into islands --- named honestly as the open frontier of $\Omega$-N.

\section{P$\Omega$.N4 --- Connectivity as a real hypothesis}

The clean class \(\mathfrak F_1\) quietly assumed several things that the finite
theory got for free or nearly free: a positive summable Perron weight, boundedness
of the coupling operator on \(\ell^2(\pi)\), weighted summability of the forcing
budget, and a uniform spectral gap. At infinite scale these are not automatic;
they encode genuine connectivity and analytic questions with no finite analogue.
P\(\Omega\).N4 names them as the open frontier of \(\Omega\)-N.

\subsection{What the finite case assumed away}

\begin{remark}[Connectivity splits into inequivalent notions]\label{rem:connsplit}
For a finite graph, strong connectivity is a single condition: every interior reaches every
other in finitely many steps. At infinite scale this single notion splits into several that no
longer coincide:
\begin{enumerate}[label=\textup{(\alph*)},leftmargin=1.2cm]
\item \emph{Strong connectivity:} each reaches each, but possibly in unboundedly many steps;
\item \emph{Uniform reachability:} the steps needed are bounded across the body;
\item \emph{Geometric reachability:} the influence reaching across the body decays no faster
than exponentially.
\end{enumerate}
Finitely these collapse to one; infinitely they are distinct, and only the stronger ones support
the certificate. There may be interiors reachable only in the limit --- present in the body, yet
approached by the flow of giving without ever quite being arrived at.
\end{remark}

\begin{remark}[The Perron weight may fail to exist]\label{rem:perronfail}
Clause \(\mathfrak F_1\)(a) assumed a positive summable left-Perron weight \(\pi\). Finite
Perron--Frobenius \emph{guaranteed} this for any strongly connected graph. Infinitely it can
fail: the associated Markov chain may be
\begin{enumerate}[label=\textup{(\alph*)},leftmargin=1.2cm]
\item \emph{null-recurrent} --- a stationary measure exists but is not summable, so it cannot be
normalised to weights; or
\item \emph{transient} --- mass escapes to infinity and no stationary distribution exists at
all.
\end{enumerate}
In either case the very weighting that organised H.3 --- each interior's place in the flow of
giving --- has no well-defined infinite form: the flow of giving can leak away to infinity,
leaving no stationary share for each.
\end{remark}

\begin{remark}[Fragmentation]\label{rem:fragment}
Even a connected infinite body can be effectively fragmented: weakly joined across some cut by a
coupling so attenuated that the two sides barely commune (no spectral gap, P$\Omega$.N3), or
genuinely split into infinitely many components. The infinite body may be a union of islands
that do not truly share the one communion.
\end{remark}

\subsection{The open frontier}

\begin{remark}[What remains open]\label{rem:nfrontier}
Characterising the infinite communions that genuinely cohere --- uniformly reachable,
carrying a summable Perron weight, possessing a spectral gap, unfragmented --- is the open
frontier of $\Omega$-N. The clean class \(\mathfrak F_1\) names these as hypotheses; proving that
any given infinite communion, and in particular the whole \emph{communio sanctorum}, satisfies
them is not established here.
\end{remark}

\subsection{Status and reading}

\textbf{Target status:}~
\stInterp\ + \stOpen. N4 proves nothing new; it exposes the connectivity assumptions that the
finite theory hid and that \(\mathfrak F_1\) silently made, and names their failure modes. It
depends on H.3 and P$\Omega$.N0.

\medskip
\stInterp\quad
\emph{Reading.}
In a bounded communion, that each can reach each is simply given; one need not ask it. The
boundless body cannot assume it. There may be those reachable only at the limit --- present in
the body, approached by the giving, yet never quite arrived at; and the place of each in the
flow of giving, which organised everything, may have no whole form at all, the giving leaking
outward into the infinite with no stationary share left for anyone. The body may break into
islands that do not touch. So the unity of the whole communion --- that it is one body, all
reachable, all weighted, all sharing the single life --- is not, at infinite scale, a thing the
structure delivers; it is a set of conditions the finite body met without asking and the
infinite body may or may not meet. We can say exactly what its oneness would require. We cannot,
here, say that it holds. And there is a fitting humility in that: the wholeness of all the saints
is named as hope and condition, not seized as a theorem.

\medskip
\textbf{Next node.} P$\Omega$.N5 --- the reading that closes $\Omega$-N: that 4.x is the
theology of the communion that can be seen, and whether it is \emph{cathedral} --- whether it
holds the whole infinite body --- is the question here partly answered and partly, honestly,
left open.

\section{P$\Omega$.N5 --- The reading that closes $\Omega$-N}

\textbf{Target status:}~
\stInterp. N5 adds no structure; it states the lesson of $\Omega$-N as a whole and depends on
P$\Omega$.N1--N4.

\medskip
\stInterp\quad
\emph{Reading.}
We took the structure built for a communion that can be seen --- a parish, a generation, a
finite ensemble --- and asked whether it holds for the whole, the unnumbered body of all who were
and will be. The answer comes in layers, and honesty keeps them apart. The single one is safe at
any scale: that result is robust and asks for nothing (P$\Omega$.N1). But the whole-as-one ---
one breath shared, one disagreement bounded, all reachable, all weighted into the single life ---
is not delivered by the structure simply because the structure is sound. It is a set of
conditions: that the body's variety be bounded, so there is a common pace (N2); that its coupling
have a spectral gap, so the communion does not slowly fray (N3); that giving reach across it and
each retain a place in the giving, so it is one body and not a union of islands (N4). The finite
communion met all of these without asking. The infinite one may or may not; and where it does
not, what fails is never the safety of the member, only the oneness of the whole. So 4.x is the
theology of the communion that can be seen, and whether it is \emph{cathedral} --- whether it
holds the entire body --- is here partly proved and partly, faithfully, left at the frontier.
We secured the one, and we named, exactly, what the wholeness of the all would require, without
presuming to hand it over.

\section*{$\Omega$-N closes --- infinite communion}

\notice{\textbf{The direction is mapped.} The right space and first clean class (N0), the
survival of individual stability (N1), the uniform-floor dichotomy --- cathedral pacing or none
(N2), the plural certificate alive exactly under a spectral gap (N3), and the connectivity
hypotheses the finite case hid (N4). At infinite scale the single one stays safe; the whole
coheres as one body only under named conditions --- bounded diversity, a spectral gap, uniform
reachability, a summable Perron weight --- which are not free and not here established for the
whole \emph{communio sanctorum}. A change of apparatus honestly carried to where it can go, and
no further.}

\section*{Part~V closes --- the four silences}

\noindent Stage~$\Omega$ ends by classifying, not occupying. Of the four places beyond 4.x, two
have been approached honestly and two are kept as named silences; and to know the \emph{kind} of
each silence is the maturity of an apophatic work.

\begin{enumerate}[label=\textup{\arabic*.},leftmargin=1.4cm]
\item \textit{Tota simul} --- the \textbf{limit of language}. The eternity the created chronos is
defined against, describable by us only as the absence we are. \emph{Response:} silence that
points. \hfill \stOpen
\item \emph{Canonical meta-chronos} --- the \textbf{limit of creation}. Nonzero
holonomy forbids synchronising charts; exact subclasses may admit charts, but the
created relational data supply no canonical master-frame. \emph{Response:}
\(\Omega\)-M: obstruction, characterisation, and non-canonicity.
\hfill \stDerived
\item \emph{Infinite networks} --- the \textbf{limit of apparatus}. Not a boundary of principle,
only of method. \emph{Response:} other mathematics, $\Omega$-N, carried as far as it honestly
goes. \hfill \stDerived\,/\,\stOpen
\item \emph{Plural $\leftrightarrow$ divine geometry} --- the \textbf{limit of virtue}. The one
step the work declines on principle, because it would turn grammar into ontology and the icon
into an idol. \emph{Response:} refusal, held as a promise. \hfill \stOpen
\end{enumerate}
\[
\boxed{
\begin{aligned}
&\text{The work proves where the creature cannot stand as master,}\\
&\text{builds for a body without number as far as honesty allows,}\\
&\text{distinguishes chart from canon, relation from metronome,}\\
&\text{and keeps, on purpose, the two silences}\\
&\text{that belong to the things themselves.}
\end{aligned}}
\]

\end{document}
